\chapter*{计数原理专题练习}
\begin{enumerate}
    \item 分类加法计数原理的应用
    \begin{enumerate}
        \item 书架的第 1 层放有 3 本不同的计算机书，第 2 层放有 3 本不同的文艺书，第 3 层放有 2 本不同的体育书.从书架上任取 1 本书，不同的取法种数为\blkda{8}

        \begin{jieda}
            书架的第 1 层放有 3 本不同的计算机书, 第 2 层放有 3 本不同的文艺书,
            
            第 3 层放有 2 本不同的体育书.从书架上任取 1 本书，不同的取法种数为 $3 + 3 + 2 = 8$ .
        \end{jieda}
        
        \item 书架上放有 2 本不同的科学类图书, 3 本不同的文学类图书和 5 本不同的历史类图书，小李从中任选 1 本阅读，不同的选法共有\blkda{10}种.
        
        \item 现有 3 幅不同的油画, 4 幅不同的国画, 5 幅不同的水彩画, 从这些画中选一幅布置房间，则不同的选法共有\blkda{12}种.
        \begin{jieda}
            依题意,不同的选法共有 $3 + 4 + 5 = {12}$ 种.
        \end{jieda}

        \item 在所有的两位数中，个位数字小于十位数字且为偶数，那么这样的两位数有多少个?
        \begin{jieda}
            当个位数字是 8 时, 十位数字取 9 , 只有 1 个;
            
            当个位数字是 6 时, 十位数字可取 7, 8, 9, 共 3 个;
            
            当个位数字是 4 时, 十位数字可取 5, 6, 7, 8, 9, 共 5 个;
            
            同理可知, 当个位数字是 2 时, 共 7 个,
            
            当个位数字是 0 时，共 9 个.
            
            由分类加法计数原理知,符合条件的两位数共有 $1 + 3 + 5 + 7 + 9 = {25}$ (个).
        \end{jieda}

        \item 定义“各位数字之和为 8 的三位数叫幸运数”，比如 116, 431，则所有幸运数的个数为\blkda{36}
        \begin{jieda}
            按照百位数字进行分类讨论:
            
            当百位数是 1 , 后两位相加为 7 , 有 8 种; 当百位数是 2 , 后两位相加为 6 , 有 7 种;
            
            当百位数是 3 , 后两位相加为 5 , 有 6 种; 当百位数是 4 , 后两位相加为 4 , 有 5 种;
            
            当百位数是 5 , 后两位相加为 3 , 有 4 种; 当百位数是 6 , 后两位相加为 2 , 有 3 种;
            
            当百位数是 7 , 后两位相加为 1 , 有 2 种; 当百位数是 8 , 后两位相加为 0 , 有 1 种;
            
            总共有 $8 + 7 + 6 + 5 + 4 + 3 + 2 + 1 = {36}$ 种.
        \end{jieda}
        
        \item 数学中“凸数”是一个位数不低于 3 的奇位数，是最中间的数位上的数字比两边的数字都大的数，则没有重复数字且大于 564 的三位数中“凸数”的个数为\blkda{65}
        \begin{jieda}
            最高位是 5 的“凸数”, 中间数分别为 7, 8, 9, 分别有 6, 7, 8 个, 共有 21 个;
            
            最高位是 6 的“凸数”，中间数分别为 7, 8, 9, 分别有 6, 7, 8 个, 共有 21 个;
            
            最高位是 7 的“凸数”, 中间数分别为 8, 9 , 分别有 7, 8 个，共有 15 个;
            
            最高位是 8 的“凸数”，中间数为 9 ，有 8 个，
            
            所以没有重复数字且大于 564 的三位数中“凸数”的个数为 ${21} + {21} + {15} + 8 = {65}$ .
        \end{jieda}
        
    \end{enumerate}
    
    \item 分步乘法计数原理的应用
    \begin{enumerate}
        \item 回答下列问题:
        \begin{enumerate}
            \item 5 封不同的信投入 3 个不同的邮筒的投法有多少种？
            \item 5 个同学争夺 3 个比赛的冠军，每个比赛冠军只有 1 人，冠军获得情况共有多少种？
        \end{enumerate}
        \begin{jieda}
            \begin{enumerate}
                \item 5 封不同的信投入 3 个不同的邮筒的投法有 $3 \times  3 \times  3 \times  3 \times  3 = {3}^{5} = {243}$ 种
                
                \item 5 个同学争夺 3 个比赛的冠军，冠军获得情况共有 $5 \times  5 \times  5 = {5}^{3} = {125}$ 种.
            \end{enumerate}
        \end{jieda}

        \item 有0,1,2,3,4五个数字，问:
        \begin{enumerate}
            \item 可以组成多少个无重复数字的四位密码?
            \item 可以组成多少个无重复数字的四位数?
        \end{enumerate}
        \begin{jieda}
            \begin{enumerate}
                \item 完成“组成无重复数字的四位密码”这件事，可以分四个步骤:
                
                第 1 步，选取左边第一个位置上的数字，有 5 种选取方法；
                
                第 2 步，选取左边第二个位置上的数字，有 4 种选取方法；
                
                第 3 步，选取左边第三个位置上的数字，有 3 种选取方法；
                
                第 4 步，选取左边第四个位置上的数字，有 2 种选取方法.
                
                由分步乘法计数原理知,可组成不同的四位密码共有 $N = 5 \times  4 \times  3 \times  2 = {120}$ 个.

                \item 完成“组成无重复数字的四位数”这件事，可以分四个步骤:

                第 1 步，从 1, 2, 3, 4 中选取一个数字做千位数字，有 4 种不同的选取方法；
                
                第 2 步, 从 1, 2, 3, 4 中剩余的三个数字和 0 共四个数字中选取一个数字做百位数字, 有 4 种不同的选取方法;
                
                第 3 步，从剩余的三个数字中选取一个数字做十位数字，有 3 种不同的选取方法；
                
                第 4 步, 从剩余的两个数字中选取一个数字做个位数字, 有 2 种不同的选取方法.
                
                由分步乘法计数原理知，可组成不同的四位数共有 $N = 4 \times  4 \times  3 \times  2 = {96}$ (个).
            \end{enumerate}
        \end{jieda}
        
        \item 编号为 1, 2, 3, 4 的四位同学参观某博物馆, 该博物馆共有编号为 1,2,3,4 的四个门,若规定编号为 1,2,3,4 的四位同学进入博物馆不能走与自己编号相同的门, 则四位同学用不同的方式进入博物馆的方法种数为\blkda{81}
        \begin{jieda}
            因不能走与自己编号相同的门, 安排编号为 1 的同学进入博物馆有 3 种选法;
            
            同理编号为 2,3,4 的同学进入博物馆各有 3 种方法,
            
            由分步乘法计数原理,共有 $3 \times  3 \times  3 \times  3 = {81}$ 种方法. 
        \end{jieda}
        
        \item 以下运算结果为 ${3}^{4}$ 的是\dotfill{(\kh{D})}
        \xx{3 封不同的信投入 4 个不同的邮筒的投法}{4 个运动员争夺 3 个项目的冠军的冠军情况数 (每个项目只有一个冠军)}{3 块地种植 4 种不同的蔬菜的种法}{4 个同学购买 3 种不同的书籍，每人购买 1 本的种数}
        \begin{jieda}
            对于 ${A},3$ 封不同的信投入 4 个不同的邮筒的投法有 $4 \times  4 \times  4 = {4}^{3} = {64}$ 种,故 ${A}$ 错误; 
            
            对于 B，4 个运动员争夺 3 个项目的冠军(每个项目只有一个冠军)，共有 $4 \times  4 \times  4 = {4}^{3} = {64}$ 种可能，故 B 错误;
            
            对于 ${C},3$ 块地种植 4 种不同的蔬菜的种法有 $4 \times  4 \times  4 = {4}^{3} = {64}$ 种,故 ${C}$ 错误;
            
            对于 D,4 个同学购买 3 种不同的书籍,每人购买 1 本的种数有 $3 \times  3 \times  3 \times  3 = {3}^{4} = {81}$ 种,故 D 正确. 故选: D.
        \end{jieda}
        
        \item 高二某班级 4 名同学要参加足球、篮球、乒乓球比赛，每人限报一项， 其中甲同学不能报名足球，乙、丙、丁三位同学所报项目都不相同，则不同的报名种数有\blkda{12}种.
        \begin{jieda}
            由甲同学不能报名足球, 可得甲有 2 种报名方式
            
            乙、丙、丁三位同学所报项目都不相同，可得乙有 3 种报名方式，丙有 2 种报名方式，丁只有 1 种报名方式
            
            共分步计数原理可得共有 $2 \times  3 \times  2 \times  1 = {12}$ 种.
        \end{jieda}
        
        \item 如图，无人机光影秀中，有 8 架无人机排列成如图所示，每架无人机均可以发出 4 种不同颜色的光, 1至 5 号的无人机颜色必须相同，6、7 号无人机颜色必须相同，8 号无人机与其他无人机颜色均不相同，则这 8 架无人机同时发光时，一共可以有\blkda{36}种灯光组合.
        
        \begin{center}
        \includegraphics[width=0.3\textwidth]{bodyxb/src/bo_d67fm877aajc739ac4t0_2_143_406_369_326_0.jpg}
        \end{center}
        \begin{jieda}
            8号4种选择，当8号颜色确定后，6,7号有3种选择，1至5号有3中选择，因此一共$4\times3\times3=36$种选择.
        \end{jieda}
        
    \end{enumerate}
    \item 步中有类和类中有步
    \begin{enumerate}
        \item 为了纪念我国成功举办北京冬奥会，中国邮政发行《北京举办 2022 年冬奥会成功纪念》邮票，图案分别为冬奥会会徽“冬梦”、冬残奥会会徽“飞跃”、冬奥会吉祥物“冰墩墩”、冬残奥会吉祥物“雪容融”及“志愿者标志”，现将一套 5 枚邮票任取 3 枚，要求取出的邮票既含会徽邮票又含吉祥物邮票，则不同的取法种数为\blkda{8}
        \begin{jieda}
            由题意可知: 会徽邮票有 2 枚, 吉祥物邮票有 2 枚, “志愿者标志”有 1 枚，若任取 3 枚，取出的邮票既含会徽邮票又含吉祥物邮票，则有:
            \begin{dingautolist}{172}
                \item 若会徽邮票有 2 枚，吉祥物邮票有 1 枚，共有 $1 \times  2 = 2$ 种；
                
                \item 若会徽邮票有 1 枚，吉祥物邮票有 2 枚，共有 $2 \times  1 = 2$ 种；
                
                \item 若会徽邮票有 1 枚，吉祥物邮票有 1 枚， “志愿者标志”有 1 枚，共有 $2 \times  2 \times  1 = 4$ 种；
            \end{dingautolist}
            故共有 $2 + 2 + 4 = 8$ 种.
        \end{jieda}
        \item 现有印有数字 0, 1, 2, 6, 12, 20, 22, 26 的卡片，每种卡片均相同且有若干张. 若从中任选几张卡片并摆成一排, 则数字 20220126 的摆放方式共有\blkda{18}种.
        \begin{jieda}
            依题意,
            
            摆放 20 的方式有: 2，0 或 20 两种方式；
            
            摆放 220 的方式有: 2，2，0 或 22，0 或 2，20 三种方式；
            
            摆放 126 的方式有:1，2，6 或 12，6 或 1，26 三种方式；
            
            由分步计数原理知,数字 20220126 的摆放方式共有: $2 \times  3 \times  3 = {18}$ 种方式.
        \end{jieda}
        
        \item 中国有十二生肖，又叫十二属相，每一个人的出生年份对应了十二种动物(鼠、牛、虎、兔、龙、蛇、马、羊、猴、鸡、狗、猪)中的一种. 现有十二生肖的吉祥物各一个，已知甲同学喜欢牛、马，乙同学喜欢牛、狗和羊，丙同学所有的吉祥物都喜欢，让甲乙丙三位同学依次从中选一个作为礼物珍藏，若各人所选取的礼物都是自己喜欢的，则不同的选法有\blkda{50}种.
        \begin{jieda}
            根据题意, 分 2 种情况讨论:
            
             \ding{172} 若甲选择牛，此时乙的选择有 2 种，丙的选择有 10 种，此时有 $2 \times  {10} = {20}$ 种不同的选法
            
             \ding{173} 若甲选择马，此时乙的选择有 3 种，丙的选择有 10 种，此时有 $3 \times  {10} = {30}$ 种不同的选法
            
            则共有 ${20} + {30} = {50}$ 种选法.
        \end{jieda}
        
        \item 在如图所示的九宫格中填入数字和字母,已知三个字母: $a,b,c$ 都填到九宫格中且不能在同一行同一列, 其他每格只能从数字 1,2,3 中选择一个填入, 有公共边的两个格数字不相同，则不同的填法种数为\blkda{5184}.
        
        \begin{center}
        \includegraphics[width=0.2\textwidth]{bodyxb/src/bo_d67fm877aajc739ac4t0_2_123_1536_122_284_0.jpg}
        \end{center}

        \begin{jieda}
            先填 $a,b,c$ ,再填1,2,3,由分步乘法计数原理求解即可.

            先填 $a,b,c$ ,对于 $a$ ,有 9 个位置可以选择,
            
            除去 $a$ 所在的行和列, $b$ 有 4 个位置可以选择,
            
            再除去 $b$ 所在的行和列, $c$ 只剩 1 个位置,
            
            故填 $a,b,c$ 共有 $9 \times  4 \times  1 = {36}$ 种; 再填 1,2,3,如下图,
            
            \begin{center}
            \begin{tabular}{|c|c|c|}
            \hline
            $a$ &  \ding{175}  &  \ding{176}  \\
            \cline{1-3}
             \ding{172}  & $b$ &  \ding{177}  \\
            \cline{1-3}
             \ding{173}  &  \ding{174}  & $c$ \\
            \cline{1-3}
            \hline
            \end{tabular}
            \end{center}
            
            由于公共边的两个格数字不相同从 \ding{172} 号位置开始:
            
             \ding{172} 有 3 种选择， \ding{173} 有 2 种选择， \ding{174} 有 2 种选择， \ding{175} 有 3 种选择， \ding{176} 有 2 种选择， \ding{177}  有 2 种选择，故填 1,2,3 有 $3 \times  2 \times  2 \times  3 \times  2 \times  2 = {144}$ 种

            采用相同的方法可以判断，对于$a, b, c$其他的排列方式，填$1,2,3$也是$144$种
            
            所以共有 ${36} \times  {144} = {5184}$ 种.
        \end{jieda}
    \end{enumerate}
    
    \item 代数中的计数问题
    \begin{enumerate}
        \item 2025 有\blkda{15}个不同的正因数
        \begin{jieda}
            分解质因数可得 ${2025} = {5}^{2} \times  {3}^{4}$ ,所以 2025 的不同的正因数有 $3 \times  5 = {15}$ 个.
        \end{jieda}
        
        \item 世界数学三大猜想:“费马猜想”、“四色猜想”、“哥德巴赫猜想”，其中“四色猜想”和“费马猜想”已经分别在 1976 年和 1994 年荣升为“四色定理”和“费马大定理”. 如今， 哥德巴赫猜想仍未解决. 哥德巴赫猜想描述为: 任何不小于 4 的偶数, 都可以写成两个质数之和. (质数是指在大于 1 的自然数中, 除了 1 和它本身以外不再有其他因数的自然数). 在不超过 17 的质数中, 随机选取两个不同的数，其和为奇数取法有\blkda{6}种.
        \begin{jieda}
            不超过 17 的质数有:2、3、5、7、11、13、17，共 7 个，
            
            在这 7 个数中随机选取两个不同的数，其和为奇数，则 2 必取,
            
            然后在剩余 6 个奇数中任选一个即可,
            
            所以，不同的取法种数为 6 种.
        \end{jieda}
        
        \item “回文联”是对联中的一种，既可顺读，也可倒读. 比如，一副描绘厦门鼓浪屿景色的回文联:雾锁山头山锁雾，天连水尾水连天，由此定义“回文数”， $n$ 为自然数，且 $n$ 的各位数字反向排列所得自然数 $n$ ’ 与 $n$ 相等,这样的 $n$ 称为“回文数”,如: 1221,2413142. 则所有 6 位数中是“回文数”且各位数字不全相同的共有\blkda{891}个.
        \begin{jieda}
            6 位 “回文数”中个位与十万位数字相同且不为 0，十位与万位数字相同, 百位与千位数字相同
            
            第一步，确定个位与十万位数字，有 9 种可能，
            
            第二步，确定十位与万位数字，有 10 种可能，
            
            第三步，确定百位与千位数字，有 10 种可能，
            
            则 6 位 “回文数”共有 $9 \times  {10} \times  {10} = {900}$ (个),
            
            又 6 位 “回文数”中各位数字全相同的共有 9 个,
            
            则所有 6 位数中是“回文数”且各位数字不全相同的共有 900-9 = 891 (个).
        \end{jieda}
        
        \item 集合 $A$ 是\{ 1,2,3,4,5,6,7,8,9,10\}的子集，且 $A$ 中的元素有完全平方数， 则满足条件的集合 $A$ 共有\blkda{896}个.
        \begin{jieda}
            集合 $\{ 1,2,3,4,5,6,7,8,9,{10}\}$ 中的完全平方数有1,4,9,
            
            令 $B = \{ 1,4,9\} ,C = \{ 2,3,5,6,7,8,{10}\}$ ,
            
            则集合 $B = \{ 1,4,9\}$ 的非空子集有 ${2}^{3} - 1 = 7$ 个,
            
            集合 $C = \{ 2,3,5,6,7,8,{10}\}$ 的子集有 ${2}^{7} = {128}$ 个,
            
            则满足条件的集合 $A$ 为集合 $B$ 的非空子集与集合 $C$ 的子集的并集.
            
            故一共有 $7 \times  {128} = {896}$ 个.
        \end{jieda}
    \end{enumerate}
    
    \item 几何计数问题
    \begin{enumerate}
        \item 在直角坐标系 ${xOy}$ 中,已知 $\bigtriangleup {AOB}$ 三边所在直线的方程分别为 $x =$ 0， $y = {0,{2x}} + {3y} = {30}$ ，则 $\bigtriangleup  {AOB}$ 内部和边上整点(即横、纵坐标均为整数的点)的总数是\blkda{91}个.
        \begin{jieda}
            由题设,直线 ${2x} + {3y} = {30}$ 分别交 $x\text{ 、 }y$ 轴于 $\left( {{15},0}\right) \text{ 、 }\left( {{10},0}\right)$ ,
            
            \begin{center}
            \includegraphics[width=0.3\textwidth]{bodyxb/src/bo_d67fm877aajc739ac4t0_7_141_331_362_292_0.jpg}
            \end{center}
            
            以高为 10,宽为 15 的矩形内 (含边) 整数点有 176 个，其中直线 ${2x} + {3y} = {30}$ 上的整数点有 $\left( {{15},0}\right) \text{ 、 }\left( {{12},2}\right)$ . $\left( {9,4}\right) \text{ 、 }\left( {6,6}\right) \text{ 、 }\left( {3,8}\right) \text{ 、 }\left( {0,{10}}\right)$ ,共 6 个,
            
            所以，矩形对角线 ${AB}$ 两侧的三角形中整点的个数为 $\dfrac{{176} - 6}{2} = {85}$ 个，
            
            综上， $\bigtriangleup {AOB}$ 中整点的个数为 ${85} + 6 = {91}$ 个.
        \end{jieda}
        
        \item 从正十五边形的顶点中选出 3 个构成钝角三角形，则不同的选法有\blkda{315}.
        
        \begin{jieda}
            如图所示,以 $A$ 为钝角顶点,在直径 $A{A}^{\prime }$ 的左边取点 ${B}_{1}$ ,右边依次取 ${C}_{1},{C}_{2},\cdots ,{C}_{6}$ ,得到 6 个钝角三角形，当取 ${C}_{7}$ 时， $\bigtriangleup  {B}_{7}A{C}_{7}$ 为锐角三角形；
            
            同理,直径 $A{A}^{\prime }$ 的左边取点 ${B}_{2}$ ,右边依次取 ${C}_{1},{C}_{2},\cdots ,{C}_{5}$ ,得到 5 个钝角三角形,当取 ${C}_{6},{C}_{7}$ 时, $\bigtriangleup {B}_{2}A{C}_{6}$ 、 $\bigtriangleup {B}_{2}A{C}_{7}$ 为锐角三角形;
            
            在直径 $A{A}^{\prime }$ 的左边取点 ${B}_{6}$ 时，得到一个钝角 $\bigtriangleup {B}_{6}A{C}_{l}$ ,
            
            在直径 $A{A}^{\prime }$ 的左边取点 ${B}_{7}$ 时，没有钝角三角形.
            
            故以 $A$ 为钝角顶点的三角形共有 $6 + 5 + 4 + 3 + 2 + 1 = {21}$ (个) .
            
            以其余 14 个点为钝角顶点的三角形也各有 21 个,
            
            所以总共有 ${15} \times  {21} = {315}$ (个) 钝角三角形.
            
            \begin{center}
            \includegraphics[width=0.3\textwidth]{bodyxb/src/bo_d67fm877aajc739ac4t0_9_160_601_461_463_0.jpg}
            \end{center}
        \end{jieda}
        
        \item 过三棱柱中任意两个顶点连线作直线，在所有这些直线连线中构成异面直线的对数为\blkda{36}
        \begin{jieda}
            如图, \begin{center}
            \includegraphics[width=0.3\textwidth]{bodyxb/src/bo_d67fm877aajc739ac4t0_8_150_213_364_360_0.jpg}
            \end{center}
            
            分以下几类:
            
            棱柱侧棱与底面边之间所构成的异面直线有: $3 \times  2 = 6$ 对;
            
            棱柱侧棱与侧面对角线之间所构成的异面直线有: $3 \times  2 = 6$ 对;
            
            底面边与侧面对角线之间所构成的异面直线有: $6 \times  2 = {12}$ 对;
            
            底面边与底面边之间所构成的异面直线有: $3 \times  2 = 6$ 对;
            
            侧面对角线与侧面对角线之间所构成的异面直线有: $\dfrac{6 \times  2}{2} = 6$ 对;
            
            所以共有 $6 + 6 + {12} + 6 + 6 = {36}$ 对.
        \end{jieda}
        
        \item 如果一条直线与一个平面垂直，那么称此直线与平面构成一个 “正交线面对”.在一个正方体中，由两个顶点确定的直线与含有四个顶点的平面构成的“正交线面对”的个数是\blkda{36}个.
        \begin{jieda}
            正方体的两个顶点确定的直线有棱、面对角线、体对角线
            
            对于每一条棱，都与两个侧面构成“正交线面对”，这样的“正交线面对”有 $2 \times  {12} = {24}$个
            
            对于每一条面对角线，都与一个对角面构成“正交线面对”，这样的“正交线面对”有 12 个
            
            不存在四个顶点确定的平面与体对角线垂直，
            
            所以正方体中“正交线面对”共有 ${24} + {12} = {36}$ 个.
        \end{jieda}
        
        
    \end{enumerate}
    
    \item 数字排列问题

    \begin{enumerate}
        \item 用数字 0，1，2，3，4，5 组成的有重复数字的三位数且是偶数的个数为\blkda{38}
        \begin{jieda}
            由题意可知, 这三位数是偶数, 则说明其个位数为偶数, 即 0, 2, 4, 有 3 种选择, 而由于这是一个三位数，所以百位数不能是 0，有 5 种选择，因为存在重复数字，由此分类讨论:  
            \begin{dingautolist}{172}
                \item 当个位数为 0 时，则百位数有 5 种选择
                
                十位数有两种情况：与百位数一样, 只有一种选择；与个位数一样, 也只有一种选择

                \item 当个位数为 2 时，
                
                如果百位数为 2 , 则十位数有 6 种选择,
                
                如果百位数不为 2 , 则百位数有 4 种选择, 此时十位数可以与百位数或个位数相同, 有 2 种选择
                
                \item 当个位数为 4 时,
                
                如果百位数为 4 ，则十位数有 6 种选择，
                
                如果百位数不为 4 ，则百位数有 4 种选择，十位数可以与百位数或个位数相同，有 2 种选择
            \end{dingautolist}
            综上所述， $5 \times  1 + 5 \times  1 + 1 \times  6 + 4 \times  2 + 1 \times  6 + 4 \times  2 = {38}$ .
        \end{jieda}
        
        \item 由 0, 1, 2, 3, 5 组成的无重复数字的五位偶数共有\blkda{42}个.
        \begin{jieda}
            若五位数的个位数是0,则有 ${n}_{1} = 4 \times  3 \times  2 \times  1 = {24}$ 种情形;
            
            若五位数的个位数是 2 , 由于 0 不排首位,
            
            因此只有 1,3,5 有 3 种情形,中间的三个位置有 $3 \times  2 \times  1 = 6$ 种情形,
            
            依据分步计数原理,可得 ${n}_{2} = 3 \times  6 = {18}$ 种情形.
            
            由分类计数原理可得所有无重复五位偶数的个数为 $n = {n}_{1} + {n}_{2} = {24} + {18} = {42}$ .
        \end{jieda}
        
        \item 用 1, 2, 3, 4 四个数字组成无重复数字的四位数，其中比 2000 大的偶数共有\blkda{8}个
        \begin{jieda}
            比 2000 大, 故千位为 2,3,4,
            
            若千位为 2,则个位为 4,有 $2 \times  1 = 2$ (个) 符合题意的四位数;
            
            若千位为 3,则个位为 2 或 4,有 $2 \times  2 \times  1 = 4$ (个) 符合题意的四位数;
            
            若千位为 4,则个位为 2,有 $2 \times  1 = 2$ (个) 符合题意的四位数.
            
            根据分类加法计数原理得,一共有 $2 + 4 + 2 = 8$ (个) 符合题意的四位数.
        \end{jieda}
        
        \item 用 0、1、2、3、4 五个数字，可以组成没有重复数字的三位奇数的个数为\blkda{18}
        \begin{jieda}
            由题意可知, 末位数字为 1 或 3 ，首位数字有 3 种选择，则中间的数位有 3 种选择, 由分步乘法计数原理可知,可以组成没有重复数字的三位奇数的个数为 $2 \times  {3}^{2} = {18}$ 
        \end{jieda}

    \end{enumerate}
    \item 涂色问题
    \begin{enumerate}
        \item 已知五个区域 $A,B,C,D,E$ 依次相邻,如图所示,现在给这 5 个区域涂色，要求相邻的两个区域不能涂相同的颜色，现有 5 种不同的颜色可供选择，则不同的涂色方法有\blkda{1280}
        
        \begin{center}
        \includegraphics[width=0.2\textwidth]{bodyxb/src/bo_d67fm877aajc739ac4t0_11_143_669_341_151_0.jpg}
        \end{center}
        \begin{jieda}
            依题意,分 5 步依次对 $A,B,C,D,E$ 涂色，所以不同的涂色方法有 $5 \times  4 \times  4 \times  4 \times  4 = {1280}$ (种).
        \end{jieda}
        
        \item 如图，用 4 种不同的颜色，对四边形中的四个区域进行着色，要求有公共边的两个区域不能用同一种颜色，则不同的着色方法有\blkda{48}
        
        \begin{center}
        \includegraphics[width=0.2\textwidth]{bodyxb/src/bo_d67fm877aajc739ac4t0_12_146_1534_244_237_0.jpg}
        \end{center}
        \begin{jieda}
            将四个区域标记为 $A,B,C,D$ ,如下图所示:
            
            \begin{center}
            \begin{tabular}{|c|c|}
            \hline
            \multicolumn{2}{|c|}{$A$} \\
            \cline{1-2}
            $B$ & $C$ \\
            \cline{1-2}
            \multicolumn{2}{|c|}{$D$} \\
            \cline{1-2}
            \hline
            \end{tabular}
            \end{center}
            
            第一步涂A:4种涂法，第二步涂B:3种涂法，第三步涂C:2种涂法，第四步涂D:2种涂法，
            
            根据分步乘法计数原理可知,一共有 $4 \times  3 \times  2 \times  2 = {48}$ 种着色方法.
        \end{jieda}
        
        \item 用红、黄、蓝三种不同颜色给如图所示的 4 块区域 $A$ 、 $B$ 、 $C$ 、 $D$ 涂色，要求同一区域用同一种颜色，有共公边的区域使用不同颜色，则共有涂色方法\blkda{18}种.
        
        \begin{center}
        \includegraphics[width=0.2\textwidth]{bodyxb/src/bo_d67fm877aajc739ac4t0_12_144_675_310_233_0.jpg}
        \end{center}
        \begin{jieda}
            先涂 $A$ ,有 3 种涂法,再涂 $B$ 有 2 种涂法
            
            涂 $C$ 时,与 $A$ 同色,有 1 种涂法,此时 $D$ 有 2 种涂法;当 $C$ 与 $A$ 异色时有 1 种涂法,这是 $D$ 有 1 种涂法
            
            所以共有 $3 \times  2 \times  \left( {1 \times  2 + 1 \times  1}\right)  = {18}$ 种.
        \end{jieda}
        
        \item 用红、黄、蓝、白、黑五种颜色在“田”字型的 4 个小方格内涂色，每格涂一种颜色，相邻两格涂不同的颜色，则不同的涂色方法共有\blkda{260}种.
        \begin{jieda}
            将“田”字型的 4 个小方格分别编号为 A, B, C, D，如下图所示:
            
            \begin{center}
            \begin{tabular}{|c|c|}
            \hline
            $A$ & $B$ \\
            \cline{1-2}
            $C$ & $D$ \\
            \cline{1-2}
            \hline
            \end{tabular}
            \end{center}
            
            根据题意, 将问题分成 4 步进行,
            
            第一步:涂方格 A，共有 5 种颜色选择，
            
            第二步:涂方格 B，需与 A 不同，共有 4 种颜色选择，
            
            第三步:涂方格C，若方格C与方格B颜色相同，只有 1 种选择；若方格C与方格B颜色不同，则有 3 种选择;
            
            第四步:涂方格 D，当方格 C 与方格 B 颜色相同，方格 D 有 4 种颜色选择；当方格 C 与方格 B 颜色不同， 方格 D 有 3 种颜色选择,
            
            因此可得不同的涂色方法共有 $5 \times  4 \times  \left( {1 \times  4 + 3 \times  3}\right)  = {260}$ 种.
        \end{jieda}
        
    \end{enumerate}
    
    \item 两个计数原理的综合应用
    \begin{enumerate}
        \item 某农学院计划从 10 种不同的水稻品种和 7 种不同的小麦品种中, 选 5 种品种种植在如图所示五块实验田中，要求仅选两种小麦品种且需种植在相邻两块实验田中，其他三块实验田选种水稻品种，则不同种法有\blkda{120960}种.
        
        \begin{center}
        \begin{tabular}{|c|c|c|c|c|c|}
        \hline
        1 & 2 & 3 & 4 & 5 & \phantom{X} \\
        \cline{1-6}
        \hline
        \end{tabular}
        \end{center}
        \begin{jieda}
            由题得相邻两块实验田分成 1 和 2; 2 和 3; 3 和 4; 4 和 5 四类;
            
            第一类在 1 和 2 上种植小麦, “ 1 ”有 7 种选择, “ 2 ”有 6 种选择, 剩下 3 块实验田种植水稻,
            
            分别有 10,9,8 种选择,所以共计 $7 \times  6 \times  {10} \times  9 \times  8 = {30240}$ 种;
            
            第二、三、四类和第一类种数相同. 综上总计有 ${30240} \times  4 = {120960}$ 种方法.
        \end{jieda}
        
        \item 某单位有 5 位同事各有一辆私家车，车牌尾数分别是 0, 1, 2, 3, 5，为遵守所在城市元月 15 日至 18 日 4 天的限行规定(奇数日车牌尾数为奇数的车通行，偶数日车牌尾数为偶数的车通行)，五人商议拼车出行，每天任选一辆符合规定的车，但甲的车(车牌尾数为 2)最多只能用一天，则不同的用车方案种数是 \blkda{27}
        \begin{jieda}
            15 日至 18 日, 有 2 天奇数日和 2 天偶数日, 车牌尾数中有 3 个奇数和 2 个偶数. 通过按日期分步, 按第一个偶数日是否开甲的车分 2 类,
            
            第一类: $3 \times  1 \times  3 \times  1 = 9$ ,第二类: $3 \times  1 \times  3 \times  2 = {18}$ ,共27种.
        \end{jieda}
        \item 从0,1,2,3,4五个数字中选出 3 个数字组成一个三位数.
        \begin{enumerate}
            \item 可以组成多少个三位数？
            \item 可以组成多少个无重复数字的三位数？
            \item 可以组成多少个无重复数字的三位偶数?
        \end{enumerate}
        \begin{jieda}
            \begin{enumerate}
                \item 三位数的首位不能为 0 ，但可以有重复数字，首先考虑首位的排法，除 0 外共有 4 种方法,第二、三位可以排 0,因此,根据分步乘法计数原理共有 $4 \times  5 \times  5 = {100}$ (个).
                
                \item 三位数的首位不能为 0，首先考虑首位的排法，除 0 外共有 4 种方法，第二位可以排 0，除首位排的数字共有 4 种方法, 第三位除前两位排的数字共有 3 种方法, 因此, 根据分步乘法计数原理共有 $4 \times  4 \times  3 = {48}$ (个).
                \item 偶数末位数字可取 0,2,4 , 因此,可以分两类:
                
                一类是末位数字是 0，则有 $4 \times  3 = {12}$ (种) 排法;
                
                一类是末位数字不是 0，则末位有 2 种排法，即 2 或 4，再排首位，因 0 不能在首位，
                
                所以有 3 种排法,十位有 3 种排法,因此有 $2 \times  3 \times  3 = {18}$ (种) 排法.
                
                因此有 ${12} + {18} = {30}$ (种) 排法. 即可以排成 30 个无重复数字的三位偶数.
            \end{enumerate}
        \end{jieda}
        
        \item 高二(1)班、(48)班、(62)班分别有 7, 5, 9 人参加创新技能大赛笔试.
        \begin{enumerate}
            \item 如果选一人当组长，那么有多少种不同的选法？
            \item 如果老师任组长，每班选一名副组长，那么有多少种不同的选法？
            \item 如果推选两名学生参赛，要求这两人来自不同的班级，那么有多少种不同的选法？
        \end{enumerate}
        \begin{jieda}
            \begin{enumerate}
                \item 事件选一人当组长可分三类方案完成，
                
                第一类，组长从 (1) 班选出，有 7 种选法、
                
                第二类，组长从 (48) 班选出，有 5 种选法、
                
                第三类，组长从 (62) 班选出，有 9 种选法，
                
                根据分类加法计数原理,选一人当组长有 $7 \neq  5 + 9 = {21}$ 种选法
                \item 如果老师任组长，每班选一名副组长，则需要分三步，
                
                第一步, 从 (1) 班选一名同学担任副组长, 有 7 种选法,
                
                第二步, 从 (48) 班选一名同学担任副组长, 有 5 种选法,
                
                第三步, 从 (62) 班选一名同学担任副组长, 有 9 种选法,
                
                根据分步乘法计数原理,每班选一名副组长共有 $7 \times  5 \times  9 = {315}$ 种选法
                \item 事件推选两名学生参赛，要求这两人来自不同的班级，可分为三类方案，
                
                第一类,若两人来自(1)班和(48)班，有 $7 \times  5 = {35}$ 种选法，
                
                第二列，若两人来自(1)班和(62)班，有 $7 \times  9 = {63}$ 种选法，
                
                第三类，若两人来(48)班和(62)班，有 $5 \times  9 = {45}$ 种选法，
                
                综上可知,这两人来自不同的班级的不同的选法有 ${35} + {63} + {45} = {143}$ 种.
            \end{enumerate}
        \end{jieda}
    \end{enumerate}
    \item 元素 (位置) 有限制的排列问题
    \begin{enumerate}
        \item 从 6 人(包含甲)中选派出 3 人参加A，B，C 这三项不同的活动， 且每项活动有且仅有 1 人参加，若甲不参加 $A$ 和 $B$ 活动，则不同的选派方案有\blkda{80}种.
        \begin{jieda}
            当甲被选中时,不同的选派方案有 $P_{5}^{2} = {20}$ 种;
            
            甲没被选中时，不同的选派方案有 $P_{5}^{3} = {60}$ 种.
            
            故满足条件的不同的选派方案有 ${20} + {60} = {80}$ 种.
        \end{jieda}
        \item 某工厂生产一种产品需经过一，二，三，四共 4 道工序，现要从 $A,B,C,D,E,F$ 这 6 名员工中选出 4 人，安排在 4 道工序上工作(每道工序安排一人)，如果员工A不能安排在第四道工序，则不同的安排方法共有\blkda{300}种.
        \begin{jieda}
            从 6 名员工中任选 4 人，安排在 4 道工序上工作的安排方法数为 $P_{6}^{4}$ 种,
            
            其中员工 ${A}$ 在第四道工序工作的安排方法数为 $P_{5}^{3}$ 种,
            
            所以不同的安排方法共有 $P_{6}^{4} - P_{5}^{3} = {300}$ (种).
        \end{jieda}
        
        \item 某大学 2023 年继续开展基础学科招生改革试点(以下简称强基计划), 以“为国选才育才”为宗旨，探索多维度考核评价模式，选拔一批有志向、有兴趣、有天赋的青年学生进行专门培养，为国家重大战略领域输送后备人才. 某市通过初审考核，甲、乙、丙、丁、戊五名同学成功入围该大学强基计划复试，参加学科基础素质测试，决出第一到第五名的名次(无并列名次). 甲和乙去询问成绩，回答者对甲说:“很遗憾，你和乙都没有得到冠军”，对乙说:“你当然不会是最差的”从这两个回答分析，5 人的名次排列可能有多少种不同情况有\blkda{54}种
        \begin{jieda}
            依题意甲、乙都没有排在第一名, 且乙没有排在第五名,
            
             \ding{172} 甲排在第五名，则有 $P_{3}^{1}P_{3}^{3} = {18}$ 种排法;
            
             \ding{173} 甲没有排在第五名，则甲、乙有 $P_{3}^{2}$ 种排法，其余人全排列，故有 $P_{3}^{2}P_{3}^{3} = {36}$ 种排法；
            
            综上可得一共有 ${18} + {36} = {54}$ 种不同的排法.
        \end{jieda}

        
        \item 由0,1,2,3,4这五个数字.
        \begin{enumerate}
            \item 能组成多少个无重复数字的五位数？
            
            \item 能组成多少个无重复数字的五位偶数？
            
            \item 组成无重复数字的五位数中比 21034 大的数有多少个？
        \end{enumerate}
        \begin{jieda}
            \begin{enumerate}
                \item 先排数字 0,0 只能占除最高位外的其余四个数位，有 $P_{4}^{1}$ 种排法，
                
                再排四个非 0 数字有 $P_{4}^{4}$ 种,由分步乘法计数原理得 $P_{4}^{1}P_{4}^{4} = 4 \times  {24} = {96}$ ,
                
                所以能组成 96 个无重复数字的五位数
                \item 当个位数字为 0 时，则可以组成 $P_{4}^{4} \equiv  {24}$ 个无重复数字的五位偶数，
                
                当个位数字为 2 或 4 时，则可以组成 ${{C}}_{2}^{1}{{C}}_{3}^{1}P_{3}^{3} = {36}$ 个无重复数字的五位偶数，
                
                即可以组成 ${24} + {36} = {60}$ 个无重复数字的五位偶数

                \item 计算比 21034 大的五位数的个数分两类:
                
                万位比 2 大的五位数个数是 $P_{2}^{1}P_{4}^{4}$ ,
                
                万位是 2 的五位数中,千位比 1 大的有 $P_{2}^{2}P_{3}^{3}$ 个,千位是 1,百位比 0 大的有 $P_{2}^{2}P_{2}^{2}$ 个,千位是 1,百位是 0, 十位比 3 大的有 1 个,
                
                由分类加法计数原理得 $P_{2}^{1}P_{4}^{4} + P_{2}^{2}P_{3}^{3} + P_{2}^{2}P_{2}^{2} + 1 = {65}$ ,
                
                所以组成无重复数字的五位数中比 21034 大的数有 65 个.
            \end{enumerate}
        \end{jieda}
        
        \item 用 0, 1, 2, 3, 4, 5 这 6 个数字可以组成多少个符合下列条件的无重复的数字? (列式并计算)
        \begin{enumerate}
            \item 六位数
            \item 六位奇数
            \item 能被 5 整除的六位数
            \item 组成的六位数按从小到大顺序排列，第 265 个数是多少？
        \end{enumerate}
        \begin{jieda}
            \begin{enumerate}
                \item 先排首数，有 $P_{5}^{1} = 5$ 种，最后排其它有 $P_{5}^{5} = {120}$ 种，根据分步计数原理得，六位数有 $5 \times  {120} = {600}$ 种.
                
                \item 先排个位数，有 $P_{3}^{1} = 3$ 种，由 0 不能在首位，则排首位有 4 种，最后排其它有 $P_{4}^{4} = {24}$ 种，根据分步计数原理得,六位奇数有 $3 \times  4 \times  {24} = {288}$ 个.

                \item 能被 5 整除的六位数，则个位数是 0 或 5
                
                个位数是 0,则有 $P_{5}^{5}$ 种;

                个位数是 5 , 先排首位, 0 不作为首位,则有 $P_{4}^{1}$ 种排法,其余位置有 $P_{4}^{4}$ 种排法,

                所以共有 $P_{5}^{5} + P_{4}^{1}P_{4}^{4} = {216}$ 个.

                \item 首位数字不能为 0，首位数字为 1 有 $P_{5}^{5} = {120}$ 种
                
                首位数字为 2,有 $P_{5}^{5} = {120}$ 种
                
                首位数字为 3,万位数字上为 0,有 $P_{4}^{4} = {24}$ 种,此时所有 6 位数有 ${120} + {120} + {24} = {264}$ 个
                
                所以第 264 个数是 305421，第 265 个数是 310245.
            \end{enumerate}
        \end{jieda}
    \end{enumerate}

    \item 相邻问题的排列问题
    \begin{enumerate}
        \item 有 5 本不同的书，其中语文书 2 本，数学书 2 本，物理书 1 本. 若将其随机摆放到书架的同一层上，则相同科目的书相邻的排法有\blkda{24}种
        \begin{jieda}
            将 2 本语文书捆绑、 2 本数学书捆绑，则相同科目的书相邻的排法种数为 $P_{2}^{2}P_{2}^{2}P_{3}^{3} = 2 \times  2 \times  6 = {24}$ 种.
        \end{jieda}

        \item 有5本不同的书和2本相同的书，若将所有的书随机摆放到书架的同一层上，则相同的书相邻的排法有\blkda{720}种
        \begin{jieda}
            $P_{6}^{6} = {720}$ 种.
        \end{jieda}

        \item 春节是团圆的日子，为了烘托这一喜庆的气氛，某村组织了“村晚”. 通过海选，现有 6 个自编节目需要安排演出，为了更好地突出演出效果，对这 6 个节目的演出顺序有如下要求:“杂技节目”排在后三位，“相声”与“小品”必须相继演出，则不同的演出方案有\blkda{120}种
        \begin{jieda}
            先将“相声”与“小品”排在一起，有 $P_{2}^{2}$ 种排法，再与其它 4 个节目排序，有 $P_{5}^{5}$ 种排法，最后考虑杂技节目在前三位或在后三位情况一样,所以有 $\dfrac{P_{2}^{2}P_{5}^{5}}{2} = {120}$ 种.
        \end{jieda}
        \item 北京时间 2023 年 10 月 26 日 19 时 34 分，神舟十六号航天员乘组(景海鹏，杜海潮，朱杨柱 3 人)顺利打开“家门”，欢迎远道而来的神舟十七号航天员乘组(汤洪波，唐胜杰，江新林 3 人) 人驻“天宫”. 随后，两个航天员乘组拍下“全家福”，共同向全国人民报平安. 若这 6 名航天员站成一排合影留念, 唐胜杰与江新林相邻, 景海鹏不站最左边, 汤洪波不站最右边, 则不同的排法有\blkda{156}种.
        \begin{jieda}
            第一步，唐胜杰、江新林 2 人相邻，有 $P_{2}^{2} = 2$ 种排法；
            
            第二步，分景海鹏站最右边与景海鹏不站最左边与最右边两种情况讨论
            
            第一种情况:景海鹏站最右边，共有 $P_{4}^{4} = {24}$ 种排法；
            
            第二种情况:景海鹏不站最左边与最右边，则共有 $P_{3}^{1}P_{3}^{1}P_{3}^{3} = {54}$ 种排法，
            
            故总共有 $2 \times  \left( {{24} + {54}}\right)  = {156}$ 种排法.
        \end{jieda}
        
        \item 为积极落实“双减”政策，丰富学生的课外活动，某校开设了陶艺、 剪纸、插花等 5 门课程.分别安排在周一到周五，每天一节，其中陶艺课不排在周一，剪纸和插花课相邻的课程的安排方案种数为 \blkda{36}
        \begin{jieda}
            剪纸和插花课相邻的安排方法有 $P_{4}^{4}P_{2}^{2} = {48}$ 种,
            
            剪纸和插花课相邻且陶艺课排在周一的安排方法有 $P_{3}^{3}P_{2}^{2} = {12}$ ,
            
            故陶艺课不排在周一,剪纸和插花课相邻的课程安排方法一共有 ${48} - {12} = {36}$
        \end{jieda}
    \end{enumerate}
    \item 不相邻排列问题
    \begin{enumerate}
        \item “四书五经”是我国 9 部经典名著《大学》《论语》《中庸》《孟子》《周易》《尚书》《诗经》《礼记》《春秋》的合称. 为弘扬中国传统文化，某校计划在读书节活动期间举办“四书五经”知识讲座，每部名著安排 1 次讲座，若要求《大学》《论语》《周易》均不相邻，则排法种数为\blkda{$P_{6}^{6}P_{7}^{3}$}
        \begin{jieda}
            先排《中庸》《孟子》《尚书》《诗经》《礼记》《春秋》这 6 次经典名著的讲座，共有 $P_{6}^{6}$ 种排法;
            
            再从 7 个空位中选 3 个，排《大学》《论语》《周易》这 3 次讲座，有 $P_{7}^{3}$ 种排法，
            
            故总共有 $P_{6}^{6}P_{7}^{3}$ 种排法
        \end{jieda}

        
        \item 有5本不同的书和2本相同的书，若将所有的书随机摆放到书架的同一层上，则相同的书不相邻的排法有\blkda{1800}种
        \begin{jieda}
            $P_{5}^{5}C_6^2 = 1800$ 种.
        \end{jieda}
        
        \item 有5本相同的书和2本不同的书，若将所有的书随机摆放到书架的同一层上，则不同的书不相邻的排法有\blkda{15}种
        \begin{jieda}
            $C_6^2 = 15$ 种.
        \end{jieda}

        \item 由0,1,2,3,4,5组成没有重复数字的六位数,其中任意两个偶数都不相邻，则满足条件的六位数的个数为\blkda{108}
        \begin{jieda}
            先排 3 个奇数,有 $P_{3}^{3} = 6$ 种排法,
            
            排完奇数后形成 4 个空,插入余下 3 个偶数,有 $P_{4}^{3} = {24}$ 种排法,
            
            但此时 0 放在首位的情况有 $P_{3}^{2} = 6$ 种,故满足条件的排法有 $6 \times  \left( {{24} - 6}\right)  = {108}$ .
        \end{jieda}
        
        \item 一位语文老师在网上购买了四书五经各一套，四书指《大学》《中庸》《论语》《孟子》，五经指《诗经》《尚书》《礼记》《周易》《春秋》，他将 9 本书整齐地放在同一层书架上，若四书，五经必须分别排在一起，且《大学》和《春秋》不能相邻，则不同方式的排列种数为\blkda{5472}
        \begin{jieda}
            四书、五经必须分别排在一起,共有 $P_{5}^{5}P_{4}^{4}P_{2}^{2} = {5760}$ 种,
            
            若《大学》和《春秋》相邻，则不符合条件，共有 $P_{3}^{3}P_{4}^{4}P_{2}^{2} = {288}$ 种，
            
            则共有 5760-288 = 5472 种.
        \end{jieda}
        
        \item 中国古代儒家提出的“六艺”指:礼、乐、射、御、书、数. 某校国学社团预在周六开展“六艺”课程讲座活动，周六这天准备排课六节，每艺一节，排课有如下要求:“乐” 排在“书”与“数”的前面，“礼”和“射”不相邻且不排在最后面，则针对“六艺”课程讲座活动的不同排课顺序共有\blkda{96}种.
        \begin{jieda}
            若“乐”排在第一节，则从御、书、数种选一节排最后一节，要满足“礼”和“射”不相邻，则有 $3 \times  {A}_{2}^{2}{A}_{3}^{2} = {36}$ 种方法,
            
            若“乐”排在第 2 节，则从书、数种选一节排最后一节或者“御”安排最后一节，要满足“礼”和“射”不相邻，则有 ${{C}}_{2}^{1}{{C}}_{2}^{1}P_{3}^{3} + {{C}}_{2}^{1}P_{3}^{3} + {{C}}_{2}^{1}P_{2}^{2} = {40}$ 种方法,
            
            若“乐”排在第 3 节，则从书、数种选一节排最后一节，要满足“礼”和“射”不相邻，则有 ${{C}}_{2}^{1}{{C}}_{2}^{1}{{C}}_{2}^{1}{{C}}_{2}^{1} = {16}$ 种方法,
            
            若“乐”排在第 4 节，则“书”与“数”排最后两节，要满足“礼”和“射”不相邻，则有 $P_{2}^{2}P_{2}^{2} = 4$ 种方法，
            
            故总的方法一共有 ${36} + {40} + {16} + 4 = {96}$
        \end{jieda}
    \end{enumerate}
    \item 相邻、不相邻排列问题（综合）
    \begin{enumerate}
        \item 某种产品的加工需要经过 6 道工序.
        \begin{enumerate}
            \item 若其中某 2 道工序不能放在最前面也不能放在最后面，问有多少种加工顺序？
            
            \item 若其中某 3 道工序必须相邻. 问有多少种加工顺序？
            
            \item 若其中某 3 道工序两两不能相邻，问有多少种加工顺序？
        \end{enumerate}
        \begin{jieda}
            \begin{enumerate}
                \item 先从另外 4 道工序中任选 2 道工序放在最前面与最后面,有 $P_{4}^{2} = {12}$ 种不同的排法，再将其余的 4 道工序全排列,有 $P_{4}^{4} = {24}$ 种不同的排法，由分步乘法计数原理可得,共有 ${12} \times  {24} = {288}$ 种加工顺序.

                \item 先排这 3 道工序，有 $P_{3}^{3} = 6$ 种不同的排法，再将它们看作一个整体，与其余的 3 道工序全排列，有 $P_{4}^{4} = {24}$ 种不同的排法，由分步乘法计数原理可得,共有 $6 \times  {24} = {144}$ 种加工顺序.
                
                \item 先排其余的 3 道工序，有 $P_{3}^{3} = 6$ 种不同的排法，有 4 个空档，再将这 3 道工序插入空档，有 $P_{4}^{3} = {24}$ 种不同的排法，由分步乘法计数原理可得，共有 $6 \times  {24} = {144}$ 种加工顺序.
            \end{enumerate}
        \end{jieda}
        \item 某次介绍会需要安排 6 个产品的介绍顺序，其中 3 个产品来自 $A$ 公司,2 个产品来自 $B$ 公司,1 个产品来自 $C$ 公司.
        \begin{enumerate}
            \item 求 $B$ 公司的 2 个产品的介绍顺序相邻的方案数
            \item 求同一个公司产品的介绍顺序不相邻， $C$ 公司的产品既不是第一个介绍，也不是最后一个介绍的方案数.
        \end{enumerate}
        \begin{jieda}
            \begin{enumerate}
                \item 将 $B$ 公司的 2 个产品的介绍顺序捆绑在一起，与其他 4 个产品进行全排列,共有 $P_{2}^{2}P_{5}^{5} = 2 \times  5 \times  4 \times  3 \times  2 \times  1 = {240}$ 种，故 $B$ 公司的 2 个产品的介绍顺序相邻的方案数有 240 种
                \item 先排 $A$ 公司的 3 个产品有 $P_{3}^{3} = 6$ 种排法
                
                由于同一个公司产品的介绍顺序不相邻, 故分两类情况:
                
                一是 $B$ 公司的 2 个产品和 $C$ 公司的 1 个产品都在 $A$ 公司的 3 个产品之间,
                
                即 $B$ 公司的 2 个产品中的 1 个和 $C$ 公司的 1 个产品相邻,
                
                共有 ${{C}}_{2}^{1}P_{2}^{2}P_{2}^{2} = 8$ 种排法,
                
                二是 $B$ 公司的 2 个产品中的 1 个和 $C$ 公司的 1 个产品在 $A$ 公司的 3 个产品之间,
                
                另一个在第一个或最后一个,共有 ${{C}}_{2}^1{{C}}_{2}^1P_{2}^{2} = 8$ ,
                
                所以共有 $6 \times  \left( {8 + 8}\right)  = {96}$ 种方案.
            \end{enumerate}
        \end{jieda}
        \item 现有 8 名师生站成一排照相，其中老师 2 人，男学生 4 人，女学生 2 人，在下列情况下，各有多少种不同的站法？
        \begin{enumerate}
            \item 老师站在最中间，2名女学生分别在老师的两边且相邻，4 名男学生两边各 2 人
            \item 4 名男学生互不相邻，男学生甲不能在两端
            \item 2 名老师之间必要有男女学生各 1 人.
        \end{enumerate}
        \begin{jieda}
            \begin{enumerate}
                \item 由题意可得共 $P_{2}^{2}P_{2}^{2}P_{4}^{4} = 2 \times  2 \times  {24} = {96}$ 种不同的站法.
                
                \item 先排老师和女学生共有 $P_{4}^{4}$ 种站法，再排男学生甲有 ${{C}}_{3}^{1}$ 种站法，最后排剩余的 3 名男学生有 $P_{4}^{3}$ 种站法，所以共有 $P_{4}^{4}{{C}}_{3}^{1}P_{4}^{3} = {24} \times  3 \times  {24} = {1728}$ 种不同的站法.
                \item 先任选一男学生一女学生站两位老师中间，有 ${{C}}_{2}^{1}{{C}}_{4}^{1}P_{2}^{2}$ 种站法，两老师的站法有 $P_{2}^{2}$ 种，再将一男学生一女学生两位老师进行捆绑与剩余的 4 个人进行全排列有 $P_{5}^{5}$ 种
                
                所以共有 ${{C}}_{2}^{1}{{C}}_{4}^{1}P_{2}^{2}P_{2}^{2}P_{5}^{5} = 2 \times  4 \times  2 \times  2 \times  {120} = {3840}$ 种不同的站法.
            \end{enumerate}
        \end{jieda}
        \item 有 3 名男生和 4 名女生相约一起去观看电影, 他们的座位在同一排且连在一起. (列出算式, 并计算出结果)
        \begin{enumerate}
            \item 女生必须坐在一起的排法有多少种？
            \item 女生互不相邻的排法有多少种？
            \item 甲、乙两位同学相邻且都不与丙同学相邻的排法有多少种？
        \end{enumerate}
        \begin{jieda}
            \begin{enumerate}
                \item 先将 4 名女生排在一起，有 $P_{4}^{4}$ 种排法，将排好的女生视为一个整体，再与 3 名男生进行排列,共有 $P_{4}^{4}$ 种排法
                
                由分步乘法计数原理,共有 $P_{4}^{4} \times  P_{4}^{4} = {24} \times  {24} = {576}$ (种) 排法.
                
                \item 先将 3 名男生排好，共有 $P_{3}^{3}$ 种排法，在这 3 名男生中间以及两边共 4 个空位中插入 4 名女生，共有 $P_{4}^{4}$ 种排法，再由分步乘法计数原理,可得共有 $P_{3}^{3} \times  P_{4}^{4} = 6 \times  {24} = {144}$ (种) 排法.

                \item 先将甲、乙、丙以外的其余 4 人排好，共有 $P_{4}^{4}$ 种排法，由于甲、乙相邻，则有 $P_{2}^{2}$ 种排法，最后将排好的甲、乙这个整体与丙分别插入原先排好的 4 人产生:的 5 个空隙中，共有 $P_{5}^{2}$ 种排法
                
                由分步乘法计数原理,可得共有 $P_{4}^{4} \times  P_{2}^{2} \times  P_{5}^{2} = {24} \times  2 \times  {20} = {960}$ (种) 排法.
            \end{enumerate}
        \end{jieda}
        \item 2024 龙年春节档新片《热辣滚烫》是一部充满正能量，讲述感人故事的电影，影片通过主人公杜乐莹的成长历程，让我们感受到了奋斗和坚持的力量，激励着每个人在面对困难时勇敢向前. 现有 4 名男生和 2 名女生相约一起去观看该影片，他们的座位在同一排且连在一起. (列出算式, 并计算出结果)
        \begin{enumerate}
            \item 女生互不相邻的坐法有多少种？
            \item 若甲不坐最左端，乙不坐最右端，则不同排列方式共有多少种？
            \item 若甲不坐在两端，乙和丙相邻，则不同排列方式共有多少种？
        \end{enumerate}
        \begin{jieda}
            \begin{enumerate}
                \item 不相邻问题插空法,先排 4 个男生共有 $P_{4}^{4}$ 种方法,把 2 个女生插空有 $P_{5}^{2}$ 种方法,所以不同排列方式共有 $P_{4}^{4}P_{5}^{2} = {24} \times  {20} = {480}$ 种:
                
                \item 方法一:“间接法”，不同排列方式共有 $P_{6}^{6} - 2P_{5}^{5} + P_{4}^{4} = {720} - {240} + {24} = {504}$ 种
                
                方法二:“直接法”，一类甲坐最右端，有 $P_{5}^{5} = {120}$ 种坐法:另一类甲坐中间四个位置中的一个，有 $P_{4}^{1}P_{4}^{1}P_{4}^{4} =$ 384 种坐法. 故有 $P_{5}^{5} + P_{4}^{7}P_{4}^{1}P_{4}^{4} = {120} + {384} = {504}$ 种不同坐法.
                
                \item 方法一:共有 6 个位置，因为甲不坐在两端，所以甲有 4 种坐法，
                
                当甲确定时，要求乙和丙相邻，共有 3 种可能，
                
                所以不同排列方式共有 $4P_{2}^{2} \times  3 \times  P_{3}^{3} = 4 \times  2 \times  3 \times  6 = {144}$ 种.
                
                方法二: 第一步乙、丙相邻共有 $P_{2}^{2}$ 种方法,第二步乙、丙与余下的三人全排列共有 $P_{4}^{4}$ 种方法,第三步把甲插入到中间的 3 个空挡，有 $P_{3}^{1}$ 种方法，故共有 $P_{2}^{2}P_{4}^{4}P_{3}^{1} = {144}$ 种不同的坐法.
            \end{enumerate}
        \end{jieda}

    \end{enumerate}
    
    \item 定序问题
    \begin{enumerate}
        \item 身高各不相同的六位同学 $A$ 、 $B$ 、 $C$ 、 $D$ 、 $E$ 、 $F$ 站成一排照相
        \begin{enumerate}
            \item $A$ 与 $C$ 同学不相邻，共有多少种站法？(结果用数字作答)
            \item A 、 C 、 D 三位同学从左到右按照由高到矮的顺序站，共有多少种站法？(结果用数字作答)
        \end{enumerate}
        \begin{jieda}
            \begin{enumerate}
                \item 先排列除 $A$ 与 $C$ 外的 4 个人，有 $P_{4}^{4}$ 种方法，4 个人排列共有 5 个空，利用插空法将 $A$ 和 $C$ 插入 5 个空，有 $P_{5}^{2}$ 种方法，所以共有 $P_{4}^{4}{A}_{5}^{2} = {480}$ 种方法.
                
                \item 对6个人全排列有 $P_{6}^{6}$ 种方法， $A$ 、 $C$ 、 $D$ 全排列有 $P_{3}^{3}$ 种方法，则 $A\text{ 、 }C\text{ 、 }D$ 从左到右按高到矮的排列有 $\dfrac{P_{6}^{6}}{P_{3}^{3}} = {120}$ 种方法.
            \end{enumerate}
        \end{jieda}
        
        \item 2024 年 3 月 12 日是我国第 46 个植树节，为建设美丽新重庆，重庆市礼嘉中学高二年级 7 名志愿者参加了植树节活动, 3 名男生和 4 名女生站成一排. (最后答案用数字作答) 
        \begin{enumerate}
            \item 甲不在中间也不在两端的站法有多少种？
            
            \item 男、女相间的站法有多少种？
            
            \item 甲、乙、丙三人从左到右顺序一定的站法有多少种？
        \end{enumerate}
        \begin{jieda}
            \begin{enumerate}
                \item 甲不在中间也不在两端，故甲可选 4 个位置，其余六人可排除 $P_{6}^{6}$ 种，

                故共有 $4P_{6}^{6} = {2880}$ 种
                
                \item 先排男生，共有 $P_{3}^{3}$ 种，则女生可在男生排完后的四个空中选择四个，即有 $P_{4}^{4}$ 种，故共有 $P_{3}^{3}P_{4}^{4} = {144}$ 种

                \item 全部排好共有$A_7^7$种，由甲、乙、丙三人顺序一定，共有故 $\dfrac{P_{7}^{7}}{P_{3}^{3}} = {840}$ 种.
            \end{enumerate}
        \end{jieda}
        \item 在 2024 年宾县一中纪念“五四”活动中，获得一等奖的某节目参演人员合影留念. 3 名男生和 4 名女生站成一排. (最后答案用数字作答)
        \begin{enumerate}
            \item 甲不在中间也不在两端的站法有多少种？
            
            \item 男、女相间的站法有多少种?
            
            \item 甲、乙、丙三人从左到右顺序一定的站法有多少种？
        \end{enumerate}
        \begin{jieda}
            \begin{enumerate}
                \item 甲不在中间也不在两端，故甲可选 4 个位置，其余六人可排除 $P_{6}^{6}$ 种，故共有 $4P_{6}^{6} = {2880}$ 种
                
                \item 先排男生，共有 $P_{3}^{3}$ 种，则女生可在男生排完后的四个空中选择四个，即有 $P_{4}^{4}$ 种，故共有 $P_{3}^{3}P_{4}^{4} = {144}$ 种
                
                \item 全部排好共有 $P_{7}^{7}$ 种，由甲、乙、丙三人顺序一定，共有故 $\frac{P_{7}^{7}}{P_{3}^{3}} = {840}$ 种.
            \end{enumerate}
        \end{jieda}
        
        \item 有 3 名男生与 4 名女生, 在下列不同条件下, 分别求排法种数.
        \begin{enumerate}
            \item 全体排成一排，女生必须站在一起
            
            \item 全体排成一排，男生互不相邻
            
            \item 全体排成一行，其中甲，乙，丙三人从左至右的顺序不变
        \end{enumerate}
        \begin{jieda}
            \begin{enumerate}
                \item 将女生看成一个整体，与3名男生在一起进行全排列，有 $P_{4}^{4}$ 种方法，再将 4 名女生进行全排列,也有 $P_{4}^{4}$ 种方法，故共有 $P_{4}^{4} \times  P_{4}^{4} = {576}$ 种排法.
                \item 男生不相邻，而女生不作要求，所以应先排女生，有 $P_{4}^{4}$ 种方法，再在女生之间及首尾空出的 5 个空位中任选 3 个空位排男生，有 $P_{5}^{3}$ 种方法，故共有 $P_{5}^{3} \times  P_{4}^{4} = {1440}$ 种排法.
                \item 从7个位置中选四个安排除甲，乙，丙以外的4个人，有 $P_{7}^{4}$ 种方法，剩下的三个位置从左至右依次安排甲，乙，丙，仅有一种安排，故共有 $P_{7}^{4} = {840}$ 种排法
            \end{enumerate}
        \end{jieda}
        \item 有 3 名男生和 4 名女生，根据下列不同的要求，求不同的排列方法种数.
        \begin{enumerate}
            \item 全体排成一行，其中 3 名男生必须排在一起
            \item 全体排成一行，3 名男生互不相邻
            \item 全体排成一行，其中甲、乙、丙三人从左至右的顺序不变
            \item 全体排成一行，其中甲不在最左边，乙不在最右边.
        \end{enumerate}
        \begin{jieda}
            \begin{enumerate}
                \item 捆绑法. 将男生看成一个整体，进行全排列，再与其他元素进行全排列，共有 $P_{3}^{3}P_{5}^{5} = {720}$ (种) 排法
                
                \item 插空法. 先排女生，然后在空位中插入男生，共有 $P_{4}^{4}P_{5}^{3} = {1440}$ (种) 排法
                
                \item 定序排列. 7 名学生排成一行，分两步:
                
                第一步，设固定甲、乙、丙从左至右顺序的排列总数为 $N$ ；
                
                第二步，对甲、乙、丙进行全排列. 由乘法原理得 $P_{7}^{7} = {N} \times  P_{3}^{3}$ ,
                
                所以 $N = \dfrac{P_{7}^{2}}{P_{3}^{3}} = {840}$ (种)
                
                \item 位置分析法. 先排最左边，除去甲外有 $P_{6}^{1}$ 种排法，余下的 6 个位置全排有 $P_{6}^{6}$ 种排法， 但应剔除乙在最右边的排法 $P_{5}^{1}P_{5}^{5}$ 种，则符合条件的排法共有 $P_{6}^{1}P_{6}^{6} - P_{5}^{1}P_{5}^{5} = {3720}$ (种).
            \end{enumerate}
        \end{jieda}
    \end{enumerate}

    
    \item 有关排列数的计算与证明
    \begin{enumerate}
        \item 求 $P_{3}^{2} + P_{4}^{2}$ 的值为\blkda{18}
        \begin{jieda}
            $P_{3}^{2} + P_{4}^{2} = 3 \times  2 + 4 \times  3 = {18}$ .
        \end{jieda}
        \item $P_{6}^{2} + P_{5}^{2} =$\blkda{50}
        \item $n \in  {{N}}^{ * },n < {20}$ ,则 $\left( {{21} - n}\right) \cdots \left( {{100} - n}\right)$ 等于\dotfill{(\kh{A})}
        \xx{$P_{{100} - n}^{80}$}{$P_{{100} - n}^{{20} - n}$}{$P_{{100} - n}^{87}$}{$P_{{20} - n}^{80}$}
        \begin{jieda}
            因 $n \in  {{N}}^{ * }$ 且 $n < {20},\left( {{21} - n}\right) \left( {{22} - n}\right) \cdots \left( {{100} - n}\right)$ 表示 80 个连续正整数的乘积, 其中最大因数为 ${100} - n$ ,最小因数为 ${21} - n$ ,由排列数公式的意义得结果为 $P_{{100} - n}^{80}$ ,
            
            所以 $\left( {{21} - n}\right) \left( {{22} - n}\right) \cdots \left( {{100} - n}\right)  = P_{{100} - n}^{80}$ .
        \end{jieda}
        \item 设 $n \in  {{N}}_{ + }$ ,且 $n < {19}$ ，则 $\left( {{19} - n}\right)  \cdot  \left( {{20} - n}\right) \cdots \cdots \left( {{2024} - n}\right)$ 等于\dotfill{(\kh{D})}
        \xx{$P_{{2024} - n}^{2005}$}{$P_{{2024} - n}^{{2006} - n}$}{$P_{{2024} - n}^{19}$}{$P_{{2024} - n}^{2006}$}
        \begin{jieda}
            先确定最大数,即 ${2024} - n$ ,再确定因数的个数,即 $\left( {{2024} - n}\right)  - \left( {{19} - n}\right)  + 1 = {2006}$ , 所以原式 $= P_{{2024} - n}^{2006}$ .
        \end{jieda}
        \item 计算:
        \begin{dingautolist}{172}
            \item $P_{4}^{1} + P_{4}^{2} + P_{4}^{3} + P_{4}^{4}$
            \item $4P_{4}^{2} + 5P_{5}^{3}$ 
            \item $P_{n}^{2} = 7P_{n - 4}^{2}$ ，求 $n$
        \end{dingautolist}
        \begin{jieda}
            \begin{dingautolist}{172}
                \item $P_{4}^{1} + P_{4}^{2} + P_{4}^{3} + P_{4}^{4} = 4 + 4 \times  3 + 4 \times  3 \times  2 + 4 \times  3 \times  2 \times  1 = {64}$
                \item $4{A}_{4}^{2} + 5{A}_{5}^{3} = 4 \times  4 \times  3 + 5 \times  5 \times  4 \times  3 = {348}$
                \item 由 $P_{n}^{2} = 7P_{n - 4}^{2}$ ,得 $n - 4 \geq  2,n \in  {{N}}^{ * }$ ,即 $n \geq  6,n \in  {{N}}^{ * }$ ,则 $n\left( {n - 1}\right)  = 7\left( {n - 4}\right) \left( {n - 5}\right)$ , 整理得 $\left( {{3n} - {10}}\right) \left( {n - 7}\right)  = 0$ ,所以 $n = 7$ .
            \end{dingautolist}
        \end{jieda}

        \item 求证:
        \begin{dingautolist}{172}
            \item $P_{7}^{4} + 4P_{7}^{3} = P_{8}^{4}$
            \item $P_{n}^{m} + mP_{n}^{m - 1} = P_{n + 1}^{m}$
        \end{dingautolist}
        \begin{jieda}
            \begin{dingautolist}{172}
                \item 证明: $P_{7}^{4} + 4P_{7}^{3} = \dfrac{7!}{3!} + \dfrac{4 \times  7!}{4!} = \dfrac{4 \times  7! \times  2}{4!} = \dfrac{8!}{4!} = P_{8}^{4}$ .
                \item 证明: 
                
            \end{dingautolist}
        \end{jieda}
    \end{enumerate}
    
    \item 排列数方程和不等式
    \begin{enumerate}
        \item 解方程：$P_{2n}^{3} = 2P_{n + 1}^{4}$
        \begin{jieda}
            因为 $P_{2n}^{3} = 2P_{n + 1}^{4}$ ,

            由 $\left\{  \begin{array}{l} {2n} \geq  3 \\  n + 1 \geq  4 \\  n \in  {{N}}^{ * } \end{array}\right.$ ，解得 $n \geq  3$ ，
            
            由原式可得 ${2n}\left( {{2n} - 1}\right) \left( {{2n} - 2}\right)  = 2\left( {n + 1}\right) n\left( {n - 1}\right) \left( {n - 2}\right)$ ，解得 $n = 5$ 或 $n = 0$ 或 $n = 1$.
            
            又因为 $n \geq  3$ ，所以 $n = 5$ .
        \end{jieda}
        \item 已知 $3P_{9}^{x} = 4P_{9}^{x - 1}$ ，则 $x$ 等于\blkda{6}
        \begin{jieda}
            由题意得 $3 \times  \dfrac{8!}{\left( {8 - x}\right) !} = 4 \times  \dfrac{9!}{\left( {{10} - x}\right) !}$ ,
            
            化简可得 $3 = 4 \times  \dfrac{9}{\left( {{10} - x}\right) \left( {9 - x}\right) }$ ,解得 $x = {13}$ 或 6 ,
            
            因为 $\left\{  \begin{matrix} x \leq  8 \\  x - 1 \leq  9 \end{matrix}\right.$ ,所以 $x \leq  8$ 且 $x \in  {N}^{ * }$ ,故 $x = 6$ 
        \end{jieda} 
        \item 解方程：$P_{2n}^{3} = 2P_{n + 1}^{4}$
        \begin{jieda}
            因为 $P_{2n}^{3} = 2P_{n + 1}^{4}$，所以$\left\{ \begin{array}{l} {2n} \geq  3 \\    n + 1 \geq  4, \\  n \in  {{N}}^{ * } \end{array}\right.$
            
            由原式可得 ${2n}\left( {{2n} - 1}\right) \left( {{2n} - 2}\right)  = 2\left( {n + 1}\right) n\left( {n - 1}\right) \left( {n - 2}\right)$ ,解得 $n = 5$ 或 $n = 0$ 或 $n = 1$ .
            
            又因为 $n \geq  3$ ，所以 $n = 5$ .
        \end{jieda}
        \item 不等式 $3P_{x}^{3} \leq  2P_{x + 1}^{2} + 6P_{x}^{2}$ 的解集为\blkda{$\{ 3,4,5\}$}
        \begin{jieda}
            易知 $x \geq  3,x \in  N$ .
            
            因为 $P_{x}^{3} = x\left( {x - 1}\right) \left( {x - 2}\right) ,P_{x + 1}^{2} = \left( {x + 1}\right) x,P_{x}^{2} = x\left( {x - 1}\right)$ ,
            
            所以原不等式可化为 ${3x}\left( {x - 1}\right) \left( {x - 2}\right)  \leq  {2x}\left( {x + 1}\right)  + {6x}\left( {x - 1}\right)$ ,
            
            所以 $3 \leq  x \leq  5$ ,
            
            所以原不等式的解集为 $\{ 3,4,5\}$ .
        \end{jieda}

        \item 不等式 $xP_{x}^{3} > 3P_{x}^{2}$ 的解集是\blkda[4]{$\left\{  {x \mid  x > 3,x \in  {{N}}_{ + }}\right\}$}
        \begin{jieda}
            由 $xP_{x}^{3} > 3P_{x}^{2}$ ,得 $\left\{  \begin{matrix} x\left\lbrack  {x\left( {x - 1}\right) \left( {x - 2}\right) }\right\rbrack   > 3\left\lbrack  {x\left( {x - 1}\right) }\right\rbrack  \\  x \geq  3 \\  x \in  {{N}}_{ + } \end{matrix}\right.$ ，解得 $x > 3,x \in  {{N}}_{ + }$
        \end{jieda}
        
        \item 不等式 $P_{8}^{x} < 6P_8^{x - 2}$ 的解集为\blkda{$\{ 8\}$}
        \begin{jieda}
            由 $P_{8}^{x} < P_{8}^{x - 2}$ ,可得 $\dfrac{8!}{\left( {8 - x}\right) !} < 6 \times  \dfrac{8!}{\left( {{10} - x}\right) !}$ 整理得 ${x}^{2} - {19x} + {84} < 0$ ,解得 $7 < x < {12}$ , 又因为 $\left\{  \begin{matrix} x \leq  8 \\  x - 2 \geq  0 \end{matrix}\right.$ ,解得 $2 \leq  x \leq  8$ ,

            综上可得 $7 < x \leq  8$ ，又由 $x \in  {{N}}^{ * }$ 所以 $x = 8$ .
        \end{jieda}

        \item 解不等式: $3P_{x}^{3} \leq  2P_{x + 1}^{2} + 6P_{x}^{2}$
        \begin{jieda}
            因为 $3P_{x}^{3} \leq  2P_{x + 1}^{2} + 6P_{x}^{2}$ ,
            
            所以 $\left\{  \begin{matrix} 3 \times  \dfrac{x!}{\left( {x - 3}\right) !} \leq  2 \times  \dfrac{\left( {x + 1}\right) !}{\left( {x - 1}\right) !} + 6 \times  \dfrac{x!}{\left( {x - 2}\right) !}\\  x \geq  3,x \in  {{N}}^{ * }\end{matrix}\right.$ 即 $\left\{  {\begin{matrix} 3 \leq  2 \times  \dfrac{x + 1}{\left( {x - 1}\right) \left( {x - 2}\right) } + \dfrac{6}{\left( x - 2\right) } \\  x \geq  3,x \in  {{N}}^{ * } \end{matrix}\text{ , }}\right.$
            
            整理得 $\left\{  \begin{matrix} \left( {{3x} - 2}\right) \left( {x - 5}\right)  \leq  0 \\  x \geq  3,x \in  {{N}}^{ * } \end{matrix}\right.$ ，解得 $\left\{  \begin{matrix} \dfrac{2}{3} \leq  x \leq  5 \\  x \geq  3,x \in  {{N}}^{ * } \end{matrix}\right.$ ，故 $x \in  \{ 3,4,5\}$ ，
            
            所以不等式解集为 $\{ 3,4,5\}$ .
        \end{jieda} 
    \end{enumerate}
    \item 组合计数问题
    \begin{enumerate}
        \item 某校计划在五四青年节期间举行歌唱比赛，高二年级某班从本班 5 名男生 4 名女生中选 4 人，代表本班参赛，按照学校要求女生至少参加 1 人至多参加 2 人，则选派方式共有 \blkda{100}种
        \begin{jieda}
            若恰有 1 名女生参加,则有 ${{C}}_{5}^{3}{{C}}_{4}^{1} = {10} \times  4 = {40}$ 种,
            
            若恰有 2 名女生参加，则有 ${{C}}_{5}^{2}{{C}}_{4}^{2} = {10} \times  6 = {60}$ 种，
            
            所以共有 ${40} + {60} = {100}$ 种不同的选派方式.
        \end{jieda}
        
        \item 学校夏季运动会需要从 4 名男生和 3 名女生中选取 4 名志愿者，则选出的志愿者中至少有 2 名女生的不同选法种数为\blkda{22}
        \begin{jieda}
            选出的志愿者中,有 2 个女生 2 个男生时,选法种数为 ${{C}}_{3}^{2}{{C}}_{4}^{2} = {18}$ 种,
            
            有 3 个女生 1 个男生时,选法种数为 ${{C}}_{3}^{3}{{C}}_{4}^{1} = 4$ 种,
            
            所以不同选法有 ${18} + 4 = {22}$ 种.
        \end{jieda}
        \item 小明准备从苹果、香橙、水蜜桃和圣女果等六种水果中买三种.
        \begin{enumerate}
            \item 若不买苹果，共有多少种买法？
            \item 若香橙和水蜜桃中至多买一种，共有多少种买法？
            \item 若香橙和圣女果中至少买一种，且香橙和苹果不同时买，共有多少种买法？
        \end{enumerate}
        \begin{jieda}
            \begin{enumerate}
                \item 若不买苹果，共有 ${{C}}_{5}^{3} = {10}$ 种买法.
                
                \item 若香橙和水蜜桃中至多买一种，共有 ${{C}}_{4}^{3} + {{C}}_{2}^{1}{{C}}_{4}^{2} = {16}$ 种买法.
                
                \item 当香橙和圣女果中只买香橙时，有 ${{C}}_{3}^{2}$ 种买法；

                当香橙和圣女果中只买圣女果时，有 ${{C}}_{4}^{2}$ 种买法；
                
                当香橙和圣女果都买时，有 ${{C}}_{3}^{1}$ 种买法.
                
                故买法总数为 ${{C}}_{3}^{2} + {{C}}_{4}^{2} + {{C}}_{3}^{1} = {12}$ 种.
            \end{enumerate}
        \end{jieda}
        
        \item 有男运动员 6 名，女运动员 4 名，其中男，女队长各 1 名.先选派 5 人外出参加比赛，在下列情形中各有多少种选派方法?
        \begin{enumerate}
            \item 任意选派 5 人
            \item 男运动员 3 名，女运动员 2 名
            \item 两个队长必须参加
            \item 至少有三名女运动员.
        \end{enumerate}
        \begin{jieda}
            \begin{enumerate}
                \item 若任意选派 5 人，则有 ${{C}}_{10}^{5} = {252}$ 种选派方法.
                \item 若男运动员3名，女运动员2名，则有 ${{C}}_{6}^{3}{{C}}_{4}^{2} = {120}$ 种选派方法.
                \item 若两个队长必须参加，则有 ${{C}}_{8}^{3} = {56}$ 种选派方法.
                \item 若只有 3 名女运动员参加，则有 ${{C}}_{4}^{3}{{C}}_{6}^{2} = {60}$ 种选派方法；
                
                若 4 名女运动员全部参加，则有 ${{C}}_{4}^{4}{{C}}_{6}^{1} = 6$ 种选派方法;
                
                可知至少有三名女运动员,则有 ${60} + 6 = {66}$ 种选派方法.
            \end{enumerate}
        \end{jieda}
        
        \item 若一个四位数的各位数字之和为 4，则称该四位数为“F 数”， 这样的“F 数”有\blkda{20}个
        \begin{jieda}
            由题意,可得 $4 = 4 + 0 + 0 + 0 = 3 + 1 + 0 + 0 = 2 + 2 + 0 + 0 = 2 + 1 + 1 + 0 = 1 + 1 +  1 + 1$
            
            当四位数为由4,0,0,0构成时,共有 1 种情况;
            
            当四位数为由3,1,0,0构成时,共有 ${{C}}_{2}^{1}{{C}}_{3}^{1} = 6$ 种情况;
            
            当四位数为由2,2,0,0构成时,共有 ${{C}}_{3}^{1} = 3$ 种情况;
            
            当四位数为由2,1,1,0构成时,共有 ${{C}}_{3}^{1}{{C}}_{3}^{1} = 9 $ 种情况;
            
            当四位数为由1,1,1,1构成时,共有 1 种情况,
            
            由分类计数原理,可得共有 $1 + 6 + 3 + 9 + 1 = {20}$ 种不同的“ $F$ 数”.
        \end{jieda}
        \item 现有如下定义:除最高数位上的数字外，其余每一个数字均比其左边的数字大的正整数叫“幸福数”(如 3467 和 1579 都是四位“幸福数”).
        \begin{enumerate}
            \item 求四位“幸福数”的个数；
            \item 如果把所有的四位“幸福数”按照从小到大的顺序排列，求第 125 个四位“幸福数”
        \end{enumerate}
        \begin{jieda}
            \begin{enumerate}
                \item 根据题意，四位“幸福数”中不能有 0 ，故只需在数字 $1, 2, 3, \cdots , 9$ 中任取 4 个，将其从小到大排列，即可得到一个四位“幸福数”，每种取法对应 1 个“幸福数”，则四位“幸福数”共有 ${{C}}_{9}^{4} = {126}$ 个.
                
                \item 对于所有的四位“幸福数”，1 在最高数位上的有 ${{C}}_{8}^{3} = {56}$ 个，2 在最高数位上的有 ${{C}}_{7}^{3} = {35}$ 个，

                3 在最高数位上的有 ${{C}}_{6}^{3} = {20}$ 个， 4 在最高数位上的有 ${{C}}_{5}^{3} = {10}$ 个， 5 在最高数位上的有 ${{C}}_{4}^{3} = 4$ 个.
                
                因为 ${56} + {35} + {20} + {10} + 4 = {125}$，所以第 125 个四位“幸福数”是最高数位为 5 的最大的四位“幸福数”, 为 5789.
            \end{enumerate}
        \end{jieda}
        \item 在平面直角坐标系中，确定若干个点，点的横、纵坐标均取自集合 $A =  \{  - 2,0,1,2\}$ ,这样的点共有 $n$ 个.
        \begin{enumerate}
            \item 求以这 $n$ 个点中的 2 个点为端点的线段的条数
            \item 求这 $n$ 个点能确定的直线的条数
            \item 若从这 $n$ 个点中选出 3 个点分别为三角形的 3 个顶点，求这样的三角形的个数
        \end{enumerate}
        \begin{jieda}
            \begin{enumerate}
                \item 点的横、纵坐标均有 4 种可能，则 $n = 4 \times  4 = {16}$ .所以所求线段的条数为 ${{C}}_{16}^{2} = {120}$
                \item 如图，在这 $n$ 个点中，仅有 4 点共线的直线有 9 条，仅有 3 点共线的直线有 6 条
                
                \begin{center}
                \includegraphics[width=0.3\textwidth]{bodyxb/src/bo_d67fn3v7aajc739ac5e0_10_142_1194_369_383_0.jpg}
                \end{center}
                
                所以这 $n$ 个点能确定的直线的条数为 ${{C}}_{16}^{2} - 9{{C}}_{4}^{2} - 6{{C}}_{3}^{2} + 9 + 6 = {63}$
                \item 从这 $n$ 个点中选出 3 个点，共有 ${{C}}_{16}^{3} = {560}$ 种选法.

                在同一条直线上的 3 个点不能构成三角形，所以所求的三角形的个数为 ${{C}}_{16}^{3} - 9{{C}}_{4}^{3} - 6{{C}}_{3}^{3} = {518}$ .
            \end{enumerate}
        \end{jieda}
        \item 四面体的顶点和各棱的中点共 10 个点. 在这 10 点中取 4 个不共面的点, 则不同的取法种数为 \blkda{141}
        \begin{jieda}
            从 10 个点中任取 4 个点有 ${{C}}_{10}^{4}$ 种取法,其中 4 点共面的情况有三类.
            
            第一类,取出的 4 个点位于四面体的同一个面上,有 $4{{C}}_{6}^{4}$ 种;
            
            第二类, 取任一条棱上的 3 个点及该棱所对棱的中点, 这 4 点共面, 有 6 种;
            
            第三类，由中位线构成的平行四边形(其两组对边分别平行于四面体相对的两条棱)，
            
            它的 4 顶点共面, 有 3 种.
            
            以上三类情况不合要求应减掉， $\therefore$ 不同的取法共有 ${{C}}_{10}^{4} - 4{{C}}_{6}^{4} - 6 - 3 = {141}$ 种.
        \end{jieda}
        \end{enumerate}
    \item 分组问题
    \begin{enumerate}
        \item 4个人分成$2+2$两组的分数方案数为\blkda{3}
        \item 4个人分成$3+1$两组的分数方案数为\blkda{4}
        \item 5个人分成$3+2$两组的分数方案数为\blkda{10}
        \item 5个人分成$2+2+1$两组的分数方案数为\blkda{15}
        \begin{jieda}
            $\dfrac{{{C}}_{5}^{2}{{C}}_{3}^{2}{{C}}_{3}^{1}}{P_{2}^{2}} = {15}$
        \end{jieda}
        \item 6个人分成$2+2+2$三组的分数方案数为\blkda{15}
        \begin{jieda}
            $\dfrac{{{C}}_{6}^{2}{{C}}_{4}^{2}{{C}}_{2}^{2}}{P_{3}^{3}} = {15}$
        \end{jieda}
    \end{enumerate}
    \item 分组分配问题
    \begin{enumerate}
        \item 某高校派出 4 名学生去三家公司实习，每位同学只能前往一家公司实习，并且每个公司至少有一名同学前来实习，则不同的安排方案有 \blkda{36}种.
        \begin{jieda}
        先分组，分成$2,1,1$三组，由$C_4^2 = 6$种分组方式

            再分配，有$A_3^3 = 6$种分配方式
            
            因此不同的方案有$36$种.
        \end{jieda}
        
        \item 将 5 名大学生分配到 3 个乡镇当村官.每个乡镇至少一名，则不同分配方案有\blkda{150}种
        \begin{jieda}
            依题意, 要使每个乡镇至少一名, 可以有 “2:2:1 "或 "3:1:1 "两种分配方案.

            按照“ 2:2:1 ”分配时，有 $\dfrac{{{C}}_{5}^{2}{{C}}_{3}^{2}{{C}}_{1}^{1}}{P_{2}^{2}} \cdot  P_{3}^{3} = {90}$ 种方法；
            
            按照“3:1:1 分配时,有 $\dfrac{{{C}}_{5 }^{3}{{C}}_{2}^{1}{{C}}_{1}^{1}}{P_{2}^{2}} \cdot  P_{3}^{3} = {60}$ 种方法.
            
            由分类加法计数原理, 可得不同分配方案有 90 + 60 = 150 种.
        \end{jieda}

        
        \item 甲、乙等 5 人去 $A,B,C$ 三个不同的景区游览，每个人去一个景区，每个景区都有人游览，若甲、乙两人不去同一景区游览，则不同的游览方法的种数为\blkda{114}
        \begin{jieda}
            去 $A,B,C$ 三个不同的景区游览,每个人去一个景区,每个景区都有人去游览,因此先分组再分配，5 个人可以分为 3 组,分别是 $\left( {1,1,3}\right) \text{ 、 }\left( {2,2,1}\right)$ 
            
            当为 $\left( {1,1,3}\right)$ 时，有 $\dfrac{{{C}}_{5}^{1}{{C}}_{4}^{1}{{C}}_{3}^{3}}{P_{2}^{2}} = {10}$ 种组合，
            
            当为 $\left( {2,2,1}\right)$ 时,有 $\dfrac{{{C}}_{5}^{2}{{C}}_{3}^{2}{{C}}_{3}^{1}}{P_{2}^{2}} = {15}$ 种组合,
            
            再分配到三个不同的景区,有 $\left( {{10} + {15}}\right)  \times  P_{3}^{3} = {150}$ 种;
            
            以上情况包含甲乙去同一景区，需要再减去此种情况，
            
            将甲乙捆绑起来作为一个元素，此时有四个元素去三个不同的景区，此时只有 $\left( {1,1,2}\right)$ 这种组合，因此有 $\dfrac{{{C}}_{4}^{1}{{C}}_{3}^{2}{{C}}_{2}^{2}}{P_{2}^{2}} = 6$ 种组合,再分配给三个不同的景区,有 $6 \times  P_{3}^{3} = {36}$ 种;
            
            因此满足题意的有: ${150} - {36} = {114}$ 种.
        \end{jieda}
        
        
        \item 甲、乙等 5 人计划去上海、苏州及青岛三个城市调查农民工薪资情况. 每个人只能去一个城市，并且每个城市都要有人去，则不同的分配方案共有种数为\blkda{150}
        \begin{jieda}
            把 5 人分组有两类情况:1:1:3 和 2:2:1.
            
            先把 5 人按 1:1:3 分组，有 ${{C}}_{5}^{3}$ 种分组方法，
            
            按 2:2:1 分组,有 $\dfrac{{{C}}_{5}^{2}{{C}}_{3}^{2}}{P_{2}^{2}}$ 种分组方法,
            
            因此不同分组方法数为 ${{C}}_{5}^{3} + \dfrac{{{C}}_{5}^{2}{{C}}_{3}^{2}}{P_{2}^{2}}$ ,
            
            再把三组人安排到三个城市，有 $P_{3}^{3}$ 种方法，
            
            所以不同分配方法种数是 $\left( {{{C}}_{5}^{3} + \dfrac{{{C}}_{5}^{2}{{C}}_{3}^{2}}{P_{2}^{2}}}\right) P_{3}^{3} = \left( {{10} + {15}}\right)  \times  6 = {150}$ .
        \end{jieda}
        

        \item 大连市普通高中创新实践学校始建于 2010 年 1 月，以丰富多彩的活动广受学生们的喜爱.现有 A, B, C, D, E 五名同学参加现代农业技术模块, 影视艺术创作模块和生物创新实验模块三个模块，每个人只能参加一个模块，每个模块至少有一个人参加，其中 A 不参加现代农业技术模块，生物创新实验模块因实验材料条件限制只能有最多两个人参加，则不同的分配方式共有\blkda{84}种.
        \begin{jieda}
            因为生物创新实验模块因实验材料条件限制只能有最多两个人参加，所以参加生物创新实验模块的为 1 人和 2 人两种情况，
            
            (1)当参加生物创新实验模块的为 1 人时，若这个人为 ${A}$ ，则一共有 $\dfrac{{{C}}_{4}^{2}}{P_{2}^{2}} \cdot  P_{2}^{2} + {{C}}_{4}^{3} \cdot  P_{2}^{2} = {14}$ 种不同的分配方式； 若这个人不是 $A$ ，则 $A$ 只能参加现代农业技术模块，一共有 ${{C}}_{4}^{1}\left( {1 + {{C}}_{3}^{2} \cdot  P_{2}^{2}}\right)  = {28}$ 种不同的分配方式；
            
            (2)参加生物创新实验模块的为2从时，若这两人中有 $A$ ，则一共有 ${{C}}_{4}^{1} \cdot  {{C}}_{3}^{2} \cdot  P_{2}^{2} = {24}$ ，
            
            若这两人中没有 $A$ ，则 $A$ 只能参加现代农业技术模块，一共有 ${{C}}_{4}^{2}\left( {1 + {{C}}_{2}^{1}}\right)  = {18}$ 种不同的分配方式；
            
            综上, 一共由 ${14} + {28} + {24} + {18} = {84}$ 种不同的分配方式
        \end{jieda}
        \item 某市教育局决定派出 8 名心理咨询专家(5 男 3 女)到甲、乙学校进行心理问题调研.
        \begin{enumerate}
            \item 每所学校均有 4 名专家参加调研，有多少种的安排方法？
            \item 每所学校至少有 3 人且必须有女专家参加调研，有多少种的安排方法？
        \end{enumerate}
        \begin{jieda}
            \begin{enumerate}
                \item 由题知,每所学校均有 4 名专家参加调研的安排方法有 ${{C}}_{8}^{4}{{C}}_{4}^{4} = {70}$ 种.
                
                \item 分成$(3,5)$两组共有$C_8^3 = 56$种方案，减去其中全男和全女的$C_5^3+C_3^3 = 11$，共$45$种方案

                分成$(4,4)$两组共有$\dfrac{C_8^4C_4^4}{A_2^2} = 35$种方案，减去其中女性都在同一组的$C_3^3C_5^1= = 5$种，共$30$种

                再分配，故一共$(45+30)\times A_2^2 = 150$种方案.
            \end{enumerate}
        \end{jieda}
        \item 6 本不同的书, 按下列要求各有多少种不同的分法?(要求:4 题最终答案均要用数字作答)
        \begin{enumerate}
            \item 分给甲、乙、丙三人，每人两本
            \item 分为三份，每份两本
            \item 分为三份，一份一本，一份两本，一份三本
            \item 分给甲、乙、丙三人，每人至少一本.
        \end{enumerate}
        \begin{jieda}
            \begin{enumerate}
                \item 先从 6 本书中选 2 本给甲，有 ${{C}}_{6}^{2}$ 种方法；再从剩余 4 本书中选 2 本给乙，有 ${{C}}_{4}^{2}$ 种方法； 最后从余下的 2 本书中选 2 本给丙,有 ${{C}}_{2}^{2}$ 种方法，根据分步计数原理,一共有 ${{C}}_{6}^{2}{{C}}_{4}^{2}{{C}}_{2}^{2} = {90}$ 种不同分法.
                
                \item 分给甲、乙、丙三人，每人两本有 ${{C}}_{6}^{2}{{C}}_{4}^{2}{{C}}_{2}^{2}$ 种方法，这个过程可以分两步完成:
                
                第一步分为三份,每份两本,设有 $x$ 种方法;
                
                第二步将这三份分给甲、乙、丙三名同学，有 $P_{3}^{3}$ 种方法.
                
                根据分步计数原理，可得 ${{C}}_{6}^{2}{{C}}_{4}^{2}{{C}}_{2}^{2} = xP_{3}^{3}$ ，所以 $x = \dfrac{{{C}}_{6}^{2}{{C}}_{4}^{2}{{C}}_{2}^{2}}{P_{3}^{3}} = {15}$ ，即共有 15 种不同分法.
                
                \item 先从 6 本书中选 1 本，有 ${{C}}_{6}^{1}$ 种方法；再从剩余 5 本书中选 2 本，有 ${{C}}_{5}^{2}$ 种方法；最后从余下的 3 本书为一份，

                根据分步计数原理，一共有 ${{C}}_{6}^{1}{{C}}_{5}^{2}{{C}}_{3}^{3} = {60}$ 种不同分法.
                
                \item 可以分为三类情况:
                \begin{dingautolist}{172}
                    \item “2、2、2 型”即(1)中的分配情况，有 ${{C}}_{6}^{2}{{C}}_{4}^{2}{{C}}_{2}^{2} = {90}$ 种分法
                    
                    \item “1、2、3 型”即(3)中的分配情况，有 ${{C}}_{6}^{1}{{C}}_{5}^{2}{{C}}_{3}^{3}P_{3}^{3} = {360}$ 种分法
                    
                    \item “1、1、4 型”，有 ${{C}}_{6}^{4}P_{3}^{3} = {90}$ 种分法
                \end{dingautolist}
                所以一共有 ${90} + {360} + {90} = {540}$ 种不同分法.
            \end{enumerate}
        \end{jieda}
        \item 国家教育部为了发展贫困地区教育，在全国重点师范大学免费培养教育专业师范生, 毕业后要分到相应的地区任教.现有 6 个免费培养的教育专业师范毕业生要按照以下要求到 3 所学校去任教, 有多少种不同的分派方法.
        \begin{enumerate}
            \item 6 人分配到三所学校甲学校 1 人、乙学校 2 人、丙学校 3 人
            
            \item 6 人分配到三所学校一校 1 人、一校 1 人、一校 4 人
            
            \item 6 人分配到三所学校每所学校至少一人
        \end{enumerate}
        \begin{jieda}
            \begin{enumerate}
                \item 6 名学生选 1 名到甲学校任教育 ${{C}}_{6}^{1}$ 种方法；从剩余的 5 名学生中选 2 名到乙学校有 ${{C}}_{5}^{2}$ 种方法; 剩余3名学生都分配到丙学校去任教有 ${{C}}_{3}^{3}$ 种方法，则 6 人分配到三所学校甲学校 1 人、乙学校 2 人、丙学校 3 人共有 ${{C}}_{6}^{7}{{C}}_{5}^{2}{{C}}_{3}^{3} = {60}$ 种分配方法

                \item 6 名学生按 1，1，4 分为三个组有 ${{C}}_{6}^{4} = {{C}}_{6}^{2}$ 种方法，则 6 人分配到三所学校一学校 1 人、一学校 1 人、 一学校 4 人共有种 ${{C}}_{6}^{2} \cdot  P_{3}^{3} = {90}$ 分配方法
                
                \item 
                \begin{dingautolist}{172}
                    \item 6 名学生按 1，2，3 分为三个组有 ${{C}}_{6}^{1}{{C}}_{5}^{2}{{C}}_{3}^{3}$ 种方法，则6人分配到三所学校共有 ${{C}}_{6}^{1}{{C}}_{5}^{2}{{C}}_{3}^{3}P_{3}^{3} = {360}$ 种分配方法
                    \item 6名学生按 1，1，4 分为三个组有 ${{C}}_{6}^{4} = {{C}}_{6}^{2}$ 种分法，则 6 人分配到三所学校一学校 1 人、一学校 1 人、一学校 1 人、一学校 4 人共有种 ${{C}}_{6}^{2} \cdot  P_{3}^{3} = {90}$ 分配方法
                    \item 6名学生平均分配到 3 所学校有 ${{C}}_{6}^{2}{{C}}_{4}^{2}{{C}}_{2}^{2} = {90}$ 种方法
                \end{dingautolist}
                则 6 人分配到三所学校每所学校至少一人一共有: ${360} + {90} + {90} = {540}$ 种方法.
            \end{enumerate}
        \end{jieda}
    \end{enumerate}
    
    \item 有关组合数的计算与证明
    \begin{enumerate}

        \item 已知 $m$ 是自然数， $n$ 是正整数，且 $m \leq  n$ . 求证: 
        \begin{enumerate}
            \item ${{C}}_{n}^{m} = {{C}}_{n}^{n - m}$
            \item ${{C}}_{n + 1}^{m} = {{C}}_{n}^{m} + {{C}}_{n}^{m - 1}$
        \end{enumerate}
        \begin{jieda}
            \begin{enumerate}
                \item 根据组合数公式，可以得到 ${{C}}_{n}^{n - m} = \dfrac{n!}{\left( {n - m}\right) !\left\lbrack  {n - \left( {n - m}\right) }\right\rbrack  !} = \dfrac{n!}{m!\left( {n - m}\right) !} = {{C}}_{n}^{m}$ .
                
                \item 根据组合数公式，可以得到 ${{C}}_{n}^{m} + {{C}}_{n}^{m - 1} = \dfrac{n!}{m!\left( {n - m}\right) !} + \dfrac{n!}{\left( {m - 1}\right) !\left( {n - m + 1}\right) !} = \dfrac{n!\left( {n - m + 1}\right) }{m!\left( {n - m + 1}\right) !} + \dfrac{n!m}{m!\left( {n - m + 1}\right) !} = \dfrac{n!\left( {n - m + 1 + m}\right) }{m!\left( {n - m + 1}\right) !} = \dfrac{\left( {n + 1}\right) !}{m!\left( {n - m + 1}\right) !} = {{C}}_{n + 1}^{m}.$
            \end{enumerate}
        \end{jieda}
        \item 式子 ${{C}}_{2n}^{9 - {2n}} + {{C}}_{{10} - n}^{2n}$ 的值为\blkda{27}
        \begin{jieda}
            由题中组合数的形式可知: $\left\{  {\begin{matrix} 0 \leq  9 - {2n} \leq  {2n} \\  {10} - n \geq  {2n} \\  9 - {2n} \in  {N} \\  {2n},{10} - n \in  {{N}}^{ * } \end{matrix} \Longrightarrow  n = 3}\right.$
            
            所以 ${{C}}_{2n}^{9 - {2n}} + {{C}}_{{10} - n}^{2n} = {{C}}_{6}^{3} + {{C}}_{7}^{6} = {27}$ .
        \end{jieda}
        \item $C_6^3+{C}_{7}^{4} + {{C}}_{8}^{5} + {{C}}_{9}^{6} + {{C}}_{10}^{7} =$ \blkda{315}
        \begin{jieda}
            ${{C}}_{6}^{3} + {{C}}_{7}^{4} + {{C}}_{8}^{5} + {{C}}_{9}^{6} + {{C}}_{10}^{7} = {{C}}_{6}^{2} + {{C}}_{6}^{3} + {{C}}_{7}^{4} + {{C}}_{8}^{5} + {{C}}_{9}^{6} + {{C}}_{10}^{7} - {{C}}_{6}^{2} = {{C}}_{11}^{7} - {15} = {330} - {15} = {315}$
        \end{jieda}

        \item 若 ${{C}}_{22}^{m} = {{C}}_{22}^{m + 2}$ ，则 ${{C}}_{3}^{2} + {{C}}_{4}^{2} + \cdots  + {{C}}_{m}^{2}$ 的值为\blkda{164}
        \begin{jieda}
            由 ${{C}}_{22}^{m} = {{C}}_{22}^{m + 2}$ ,故 $m + m + 2 = {22}$ 或 $m = m + 2$ ,故 $m = {10}$
            
            则 ${{C}}_{3}^{2} + {{C}}_{4}^{2} + \cdots  + {{C}}_{m}^{2} = {{C}}_{3}^{3} + {{C}}_{3}^{2} + {{C}}_{4}^{2} + \cdots  + {{C}}_{10}^{2} - 1 = {{C}}_{4}^{3} + {{C}}_{4}^{2} + \cdots  + {{C}}_{10}^{2} - 1 = {{C}}_{5}^{3} + {{C}}_{5}^{2} + \cdots  + {{C}}_{10}^{2} - 1 = \cdots  = {{C}}_{10}^{3} + {{C}}_{10}^{2} - 1 = {{C}}_{11}^{3} - 1 = {164}.$
        \end{jieda}
        
        \item 已知 ${{C}}_{n + 1}^{5} - {{C}}_{n}^{6} = {{C}}_{n}^{5}$ ，则 $n =$ \blkda{10}
        \begin{jieda}
            因为 ${{C}}_{n + 1}^{5} - {{C}}_{n}^{6} = {{C}}_{n}^{5}$ ,所以 ${{C}}_{n + 1}^{5} = {{C}}_{n}^{5} + {{C}}_{n}^{6}$ ,
            
            又 ${{C}}_{n}^{5} + {{C}}_{n}^{6} = {{C}}_{n + 1}^{6}$ ,所以 ${{C}}_{n + 1}^{6} = {{C}}_{n + 1}^{5}$ ,所以 $n + 1 = 5 + 6$ ,解得 $n = {10}$ .
        \end{jieda}
        \item 求值(用数字表示)
        \begin{dingautolist}{172}
            \item $P_{4}^{1} + P_{4}^{2} + P_{4}^{3} + P_{4}^{4}$
            \item ${{C}}_{5}^{3} + {{C}}_{5}^{4}$
            \item ${{C}}_{n}^{5 - n} + P_{n + 1}^{9 - n}$
        \end{dingautolist}
        \begin{jieda}
            \begin{dingautolist}{172}
                \item $P_{4}^{1} + P_{4}^{2} + P_{4}^{3} + P_{4}^{4} = 4 + 4 \times  3 + 4 \times  3 \times  2 + 4 \times  3 \times  2 \times  1 = {64}$
                
                \item ${{C}}_{5}^{3} + {{C}}_{5}^{4} = {{C}}_{5}^{2} + {{C}}_{5}^{1} = \dfrac{5 \times  4}{2 \times  1} + 5 = {15}$
                
                \item 依题意可得 $\left\{  \begin{matrix} 0 \leq  5 - n \leq  n \\  0 < 9 - n \leq  n + 1 \end{matrix}\right.$ ,又 $n \in  {{N}}^{ * }$ ,解得 $n = 4$ 或 $n = 5$
                
                当 $n = 4$ 时， ${{C}}_{n}^{5 - n} + P_{n + 1}^{9 - n} = {{C}}_{4}^{1} + P_{5}^{5} = 4 + 5 \times  4 \times  3 \times  2 \times  1 = {124}$ ；
                
                当 $n = 5$ 时， ${{C}}_{n}^{5 - n} + P_{n + 1}^{9 - n} = {{C}}_{5}^{0} + P_{6}^{4} = 1 + 6 \times  5 \times  4 \times  3 = {361}$ .
            \end{dingautolist}
        \end{jieda}
        
        \item 已知 $m,n,k \in  {{N}}^{ * },m \geq  k \geq  n$ .
        \begin{enumerate}
            \item 证明: ${{C}}_{m}^{k}{{C}}_{m - k}^{n} = {{C}}_{m}^{n}{{C}}_{m - n}^{k}$
            \item 证明: ${{C}}_{m}^{k}{{C}}_{k}^{n} = {{C}}_{m}^{n}{{C}}_{m - n}^{k - n}$
        \end{enumerate}
        \begin{jieda}
            \begin{enumerate}
                \item 因为 ${{C}}_{m}^{k}{{C}}_{m - k}^{n} = \dfrac{m!}{k!\left( {m - k}\right) !} \cdot  \dfrac{\left( {m - k}\right) !}{n!\left( {m - k - n}\right) !} = \dfrac{m!}{k!n!\left( {m - n - k}\right) !}$ ，
                
                ${{C}}_{m}^{n}{{C}}_{m - n}^{k} = \dfrac{m!}{n!\left( {m - n}\right) !} \cdot  \dfrac{\left( {m - n}\right) !}{k!\left( {m - n - k}\right) !} = \dfrac{m!}{k!n!\left( {m - n - k}\right) !}$
                
                所以 ${{C}}_{m}^{k}{{C}}_{m - k}^{n} = {{C}}_{m}^{n}{{C}}_{m - n}^{k}$ ;
                
                \item 因为 ${{C}}_{m}^{k}{{C}}_{k}^{n} = \dfrac{m!}{k!\left( {m - k}\right) !} \cdot  \dfrac{k!}{n!\left( {k - n}\right) !} = \dfrac{m!}{n!\left( {m - k}\right) !\left( {k - n}\right) !}$ ，
                
                ${{C}}_{m}^{n}{{C}}_{m - n}^{k - n} = \dfrac{m!}{n!\left( {m - n}\right) !} \cdot  \dfrac{\left( {m - n}\right) !}{\left( {k - n}\right) !\left( {m - k}\right) !} = \dfrac{m!}{n!\left( {m - k}\right) !\left( {k - n}\right) !}$
                
                所以 ${{C}}_{m}^{k}{{C}}_{k}^{n} = {{C}}_{m}^{n}{{C}}_{m - n}^{k - n}$ .
            \end{enumerate}
        \end{jieda}
    \end{enumerate}
    \item 组合数方程和不等式
    \begin{enumerate}
        \item 关于 $x$ 的方程 ${{C}}_{11}^{2x} = {{C}}_{11}^{{3x} - 4}$ 的解为 \blkda{$3$ 或 $4$}
        \begin{jieda}
            因为 ${{C}}_{11}^{2x} = {{C}}_{11}^{{3x} - 4}$ ,则 ${2x} = {3x} - 4$ 或 ${2x} + {3x} - 4 = {11}$ ,解得 $x = 4$ 或 $x = 3$ , 若 $x = 4$ ，可得 ${{C}}_{11}^{8} = {{C}}_{11}^{8}$ ，符合题意；
            
            若 $x = 3$ ，可得 ${{C}}_{11}^{6} = {{C}}_{11}^{5}$ ，符合题意；
            
            综上所述: $x = 3$ 或 $x = 4$ .
        \end{jieda}
        \item 解关于正整数 $x$ 的方程:
        \begin{dingautolist}{172}
            \item ${{C}}_{13}^{x + 1} = {{C}}_{13}^{{2x} - 3}$
            \item ${{C}}_{16}^{{x}^{2} - x} = {{C}}_{16}^{{5x} - 5}$
            \item ${{C}}_{x + 2}^{x - 2} + {{C}}_{x + 2}^{x - 3} = \dfrac{1}{4}P_{x + 3}^{3}$
        \end{dingautolist}
        \begin{jieda}
            \begin{dingautolist}{172}
                \item 由 ${{C}}_{13}^{x + 1} = {{C}}_{13}^{{2x} - 3}$，即 $x + 1 = {2x} - 3$ 或 $x + 1 + {2x} - 3 = {13}$，解得 $x = 4$ 或 $x = 5$
                \item $x$ 为正整数，由 ${{C}}_{16}^{{x}^{2} - x} = {{C}}_{16}^{{5x} - 5}$ 可得 ${x}^{2} - x = {5x} - 5$ 或 ${x}^{2} - x + {5x} - 5 = {16}$ ,
                
                故 ${x}^{2} - {6x} + 5 = 0$ 或 ${x}^{2} + {4x} - {21} = 0$ ，解得 $x = 1$ 或 $x = 5$ 或 $x = 3$ 或 $x =  - 7$ (舍去)，
                
                又 ${x}^{2} - x,{5x} - 5$ 均为整数，且 $0 \leq  {x}^{2} - x \leq  {16},0 \leq  {5x} - 5 \leq  {16}$ ，
                
                所以 $x = 1$ 或 $x = 3$ 符合要求, $x = 5$ 不符合要求,
                
                故 $x = 1$ 或 $x = 3$
                \item 由组合数的性质可得 ${{C}}_{x + 2}^{x - 2} = {{C}}_{x + 2}^{4},{{C}}_{x + 2}^{x - 3} = {{C}}_{x + 2}^{5},{{C}}_{x + 2}^{4} + {{C}}_{x + 2}^{5} = {{C}}_{x + 3}^{5}$
                
                所以由 ${{C}}_{x + 2}^{x - 2} + {{C}}_{x + 2}^{x - 3} = \dfrac{1}{4}P_{x + 3}^{3}$ 可得 ${{C}}_{x + 3}^{5} = \dfrac{1}{4}P_{x + 3}^{3}$ ,进而可得
                
                \[
                \dfrac{\left( {x + 3}\right) !}{5!\left( {x - 2}\right) !} = \dfrac{1}{4}\dfrac{\left( {x + 3}\right) !}{x!} \Rightarrow  \dfrac{1}{5!} = \dfrac{1}{{4x}\left( {x - 1}\right) } \Rightarrow  x\left( {x - 1}\right)  = {30}\text{ , }
                \]
                
                解得 $x = 6$ 或 $x =  - 5$ (舍去),
                
                由于 $\left\{  \begin{matrix} x + 2 > 0 \\  x + 3 > 0 \\  x - 3 \geq  0 \\  x - 2 \geq  0 \end{matrix}\right.$ ,所以 $x \geq  3$ ,故只取 $x = 6,x =  - 5$ 舍去.
            \end{dingautolist}
        \end{jieda}
        \item 解方程：$\dfrac{1}{{{C}}_{5}^{m}} - \dfrac{1}{{{C}}_{6}^{m}} = \dfrac{7}{4{{C}}_{7}^{m}}\left( {m \geq  0,m \in  {N}}\right)$ .
        \begin{jieda}
            由题意 $0 \leq  m \leq  5,m \in  {N}$ ，$\dfrac{1}{{{C}}_{5}^{m}} - \dfrac{1}{{{C}}_{6}^{m}} = \dfrac{7}{4{{C}}_{7}^{m}}$ ,即 $\dfrac{m!\left( {5 - m}\right) !}{5!} - \dfrac{m!\left( {6 - m}\right) !}{6!} = \dfrac{7}{4} \times  \dfrac{m!\left( {7 - m}\right) !}{7!}$ ,
            
            化简得 ${m}^{2} - {17m} + {42} = 0$ ，解得 $m = {14}$ (舍去)或 $m = 3$ .
            
            故方程的解为 $m = 3$ .
            \end{jieda}
        
        \item 若 ${{C}}_{n}^{4} > {{C}}_{n}^{6}$ ，则 $n$ 的取值集合是\blkda{$\{ 6,7,8,9\}$}
        \begin{jieda}
            因为${{C}}_{n}^{4} > {{C}}_{n}^{6}$，所以$\left\{\begin{matrix} \dfrac{n!}{\left( {4 \times  \left( {n - 4}\right) }\right) !} > \dfrac{n!}{6! \times  \left( {n - 6}\right) !} \\  n \geq  6 \end{matrix}\right.$，即 $\left\{  \begin{matrix} {n}^{2} - {9n} - {10} < 0, \\  n \geq  6, \end{matrix}\right.$， 解得 $6 \leq  n < {10}$，故$n$ 的取值集合为 $\{ 6,7,8,9\}$ .
        \end{jieda}
        \item 使不等式 ${{C}}_{n}^{2} \geq  {{C}}_{n}^{3}$ ( $n$ 为正整数) 成立的 $n$ 的取值不可能是\dotfill{(\kh{D})}
        \xx{3}{4}{5}{6}
        \begin{jieda}
            在 ${{C}}_{n}^{2}$ 中, $n$ 为正整数, $n \geq  2$ ,在 ${{C}}_{n}^{3}$ 中, $n$ 为正整数, $n \geq  3$
            
            因为 ${{C}}_{n}^{2} \geq  {{C}}_{n}^{3}$ ,则有 $\dfrac{n\left( {n - 1}\right) }{2 \times  1} \geq  \dfrac{n\left( {n - 1}\right) \left( {n - 2}\right) }{3 \times  2 \times  1}$ ,即 $n - 2 \leq  3$ ,解得 $n \leq  5$
            
            因此有 $3 \leq  n \leq  5$ , $n$ 为正整数,所以 $n$ 的取值可以是 3 或 4 或 5 .
        \end{jieda}
    \end{enumerate}
    
    \item 元素（位置）有限制的组合+排列问题
    \begin{enumerate}
        \item 有 5 个男生和 3 个女生，从中选出 5 人担任 5 门不同学科的课代表，求分别符合下列条件的选法数.
        \begin{enumerate}
            \item 有女生但人数必须少于男生
            \item 某男生必须包括在内，但不担任数学课代表
            \item 某女生一定要担任语文课代表，某男生必须担任课代表，但不担任数学课代表
        \end{enumerate}
        \begin{jieda}
            \begin{enumerate}
                \item 先选后排，先取可以是 2 女 3 男，也可以是 1 女 4 男，先取有 ${{C}}_{5}^{3}{{C}}_{3}^{2} + {{C}}_{5}^{4}{{C}}_{3}^{1}$ 种，后排有 $P_{5}^{5}$ 种，共 $\left( {{{C}}_{5}^{3}{{C}}_{3}^{2} + {{C}}_{5}^{4}{{C}}_{3}^{1}}\right)  \cdot  P_{5}^{5} = {5400}$ (种) .
                
                \item 先选后排，但先安排该男生，有 ${{C}}_{7}^{4} \cdot  {{C}}_{4}^{1} \cdot  P_{4}^{4} = {3360}$ (种).
                
                \item 先从除去该男生、该女生的 6 人中选 3 人有 ${{C}}_{6}^{3}$ 种，再安排该男生有 ${{C}}_{3}^{1}$ 种，其中 3 人全排有 $P_{3}^{3}$ 种，共 ${{C}}_{6}^{3} \cdot  {{C}}_{3}^{1}$ . $P_{3}^{3} = {360}$ (种) .
            \end{enumerate}
        \end{jieda}
        \item 7 名身高互不相等的学生，分别按下列要求排列，各有多少种不同的排法?
        \begin{enumerate}
            \item 7 人站成一排，要求较高的 3 个学生站在一起
            \item 7 人站成一排，要求最高的站在中间，并向左、右两边看，身高逐个递减
        \end{enumerate}
        \begin{jieda}
            \begin{enumerate}
                \item 将较高的 3 个学生捆绑成一个元素，与另外 4 名学生构成 5 个元素进行自由排列，共有 $P_{5}^{5}$ 种排法;

                捆绑的 3 个学生内部可自由排列，有 $P_{3}^{3}$ 种方法；
                
                共有 $P_{5}^{5}P_{3}^{3} = {120} \times  6 = {720}$ 种

                \item 最高的站在中间，从剩余的 6 名学生中选出 3 名从高到低排列在左边，剩余的 3 名按高低进行排列在右边,共有 ${{C}}_{6}^{3}{{C}}_{3}^{3} = {20}$ 种.
            \end{enumerate}
        \end{jieda}
        \item 某电视节目的主持人邀请年龄互不相同的 5 位嘉宾逐个出场亮相.
        \begin{enumerate}
            \item 其中有 3 位老者要按年龄从大到小的顺序出场，出场顺序有多少种？
            \item 3 位老者与 2 位年轻人都要分别按从小到大的顺序出场，顺序有多少种？
        \end{enumerate}
        \begin{jieda}
            \begin{enumerate}
                \item 5 位嘉宾无约束条件的全排列有 $P_{5}^{5}$ 种，由于 3 位老者的排列顺序已定，因此满足 3 位老者按年龄从大到小的顺序出场,出场顺序有 $\dfrac{P_{5}^{5}}{P_{3}^{3}} = {20}$ (种).
                
                \item 设符合条件的顺序共有 $x$ 种，用(1)的方法可得 $x \cdot  {A}_{3}^{3} \cdot  {A}_{2}^{2} = {A}_{5}^{5}$ ，解得 $x = \dfrac{{A}_{5}^{5}}{{A}_{3}^{3} \cdot  {A}_{2}^{2}} = {10}$ ， 所以符合条件的出场顺序有 10 种.
            \end{enumerate}
        \end{jieda}
    \end{enumerate}
    \item 涂色问题
    \begin{enumerate}
        \item 用五种不同的颜色给下图中的四块区域涂色，要求相邻的区域颜色不同， 则一共有多少种不同的涂色方法?
        
        \begin{center}
        \includegraphics[width=0.2\textwidth]{bodyxb/src/bo_d67fn3v7aajc739ac5e0_16_147_210_321_321_0.jpg}
        \end{center}
        \begin{jieda}
            若选择四种颜色,则有 $P_{5}^{4} = {120}$ 种不同的涂色方法;
            
            若选择三种颜色，则有 ${{C}}_{5}^{3}P_{3}^{3} = {60}$ 种不同的涂色方法,
            
            故一共有 ${120} + {60} = {180}$ 种不同的涂色方法.
        \end{jieda}

        \item 现要用红、橙、黄、绿、青、蓝、紫 7 种颜色对某市的如图的四个区域进行着色, 有公共边的两个区域不涂同一种颜色, 则共有几种不同的涂色方法?
        
        \begin{center}
        \includegraphics[width=0.2\textwidth]{bodyxb/src/bo_d67fn3v7aajc739ac5e0_16_141_950_258_346_0.jpg}
        \end{center}
        \begin{jieda}
            由图形知, I 与 IV 可以同色, 因此涂四个区域可用 3 种颜色, 也可用 4 种颜色,
            
            用 3 种颜色涂色，即 I 与 IV 同色，有$A_7^3$种方法，
            
            用 4 种颜色涂有 $P_{7}^{4}$ 种方法,
            
            所以不同的涂色方法种数是 $P_{7}^{3} + P_{7}^{4} = {210} + {840} = {1050}$ .
        \end{jieda}
        \item 如图，某心形花坛中有 $A$ ， $B$ ， $C$ ， $D$ ， ${E5}$ 个区域，每个区域只种植一种颜色的花.
        
        \begin{center}
        \includegraphics[width=0.2\textwidth]{bodyxb/src/bo_d67fn3v7aajc739ac5e0_16_145_1844_235_236_0.jpg}
        \end{center}
        \begin{enumerate}
            \item 要把 5 种不同颜色的花种植到这 5 个区域中，每种颜色的花都必须种植，共有多少种不同的种植方案？
            \item 要把 4 种不同颜色的花种植到这 5 个区域中，每种颜色的花都必须种植，共有多少种不同的种植方案？
            \item 要把红、黄、蓝、白 4 种不同颜色的花种植到这 5 个区域中，每种颜色的花都必须种植，要求相同颜色的花不能相邻种植，且有两个相邻的区域种植红、黄 2 种不同颜色的花，共有多少种不同的种植方案？
        \end{enumerate}
        \begin{jieda}
            \begin{enumerate}
                \item 由全排列可得，共有 $P_{5}^{5} = {120}$ 种不同的种植方案.
                
                \item 第一步，先将 5 个区域选出 2 个区域种植一种相同颜色的花，共有 ${{C}}_{5}^{2}{{C}}_{4}^{1} = {40}$ 种方案；
                
                第二步，再将剩余的 3 种颜色的花种植到剩下的 3 个区域，共有 $P_{3}^{3} = 6$ 种方案.
                
                所以共有 ${40} \times  6 = {240}$ 种不同的种植方案.
                
                \item 要把 4 种不同颜色的花分别种植到这 5 个区域中，则必然有 2 个区域种植相同颜色的花.
                \begin{dingautolist}{172}
                    \item $E$ 区域种植红色的花, $A,B,C,{D4}$ 个区域中有 2 个区域种植其他相同颜色的花，则相同颜色的花必然种植在 $A,D$ 或 $B,C$ 区域,共有 $1 \times  P_{3}^{1}P_{2}^{1}P_{2}^{2} = {12}$ 种方案.
                    \item $E$ 区域种植黄色的花,同理可得,共有 $1 \times  P_{3}^{1}P_{2}^{1}P_{2}^{2} = {12}$ 种方案.
                    \item $E$ 区域种植蓝色的花,若有 2 个区域种植白色的花，则没有两个相邻的区域种植红、黄 2 种不同颜色的花,所以不可能有 2 个区域种植白色的花，故 2 个区域种植的相同颜色的花是红色或黄色的花,共有 $1 \times  P_{2}^{1}P_{2}^{1}P_{2}^{2} = 8$ 种方案.
                    \item  $E$ 区域种植白色的花，同理可得，共有 $1 \times  P_{2}^{1}P_{2}^{1}P_{2}^{2} = 8$ 种方案.
                \end{dingautolist}
                
                综上，共有 ${12} \times  2 + 8 \times  2 = {40}$ 种不同的种植方案.
            \end{enumerate}
        \end{jieda}
        
        \item 如图，四边形 ${ABCD}$ 的两条对角线 ${AC}$ ， ${BD}$ 相交于 $O$ ，现用五种颜色(其中一种为红色)对图中四个三角形 $\bigtriangleup {ABO}$ ， $\bigtriangleup {BCO}$ ， $\bigtriangleup {CDO}$ ， $\bigtriangleup {ADO}$ 进行染色，且每个三角形用一种颜色染.
        \begin{center}
        \includegraphics[width=0.3\textwidth]{bodyxb/src/bo_d67fn3v7aajc739ac5e0_18_166_222_364_236_0.jpg}
        \end{center}
        \begin{enumerate}
            \item 若必须使用红色，求四个三角形 $\bigtriangleup  {ABO}$ ， $\bigtriangleup  {BCO}$ ， $\bigtriangleup  {CDO}$ ， $\bigtriangleup  {ADO}$ 中有且只有一组相邻三角形同色的染色方法的种数
            \item 若不使用红色，求四个三角形 $\bigtriangleup  {ABO}$ ， $\bigtriangleup  {BCO}$ ， $\bigtriangleup  {CDO}$ ， $\bigtriangleup  {ADO}$ 中所有相邻三角形都不同色的染色方法的种数
        \end{enumerate}
        \begin{jieda}
            \begin{enumerate}
                \item 根据题意，要求四个三角形 $\bigtriangleup {ABO}$ , $\bigtriangleup {BCO}$ , $\bigtriangleup {CDO}$ , $\bigtriangleup {ADO}$ 中有且只有一组相邻三角形同色，而同色的相邻三角形共有 4 种,不妨假设为 $\bigtriangleup {ABO},\bigtriangleup {BCO}$ 同色

                \begin{dingautolist}{172}
                    \item 若 $\bigtriangleup {ABO}$ ， $\bigtriangleup {BCO}$ 同时染红色，则另外两个三角形共有 $P_{4}^{2}$ 种染色方法，因此这种情况共有 $P_{4}^{2} = {12}$ 种染色方法
                    \item 若 $\bigtriangleup  {ABO}$ ， $\bigtriangleup  {BCO}$ 同时染的不是红色，则它们的染色有 4 种，另外两个三角形一个必须染红色，所以这两个三角形共有 $3 \times  2 = 6$ ,因此这种情况共有 $4 \times  6 = {24}$ 种染色方法
                \end{dingautolist}
                
                综上可知有且只有一组相邻三角形同色的染色方法的种数为 $4 \times  \left( {{12} + {24}}\right)  = {144}$ 种;
                
                \item 根据题意，因为不用红色，则只有四种颜色可选，
                
                分 3 种情况讨论:

                \begin{dingautolist}{172}
                    \item 若一共使用了四种颜色，则共有 $P_{4}^{4} = {24}$ 种染色方法
                    \item 若只使用了三种颜色，则必有一种颜色使用了两次，且染在对顶的区域，所以一共有 ${{C}}_{4}^{3} \times  {{C}}_{3}^{1} \times  2 \times  P_{2}^{2} =$ 48 种染色方法
                    \item 若只使用了两种颜色，则两种颜色都使用了两次，且各自染在一组对顶区域，所以共有 ${{C}}_{4}^{2} \times  2 = {12}$ 种染色方法
                \end{dingautolist}
                
                综上可知所有相邻三角形都不同色的染色方法的种数为 ${24} + {48} + {12} = {84}$ 种
            \end{enumerate}
        \end{jieda}
        
        
        \item 如图，用四种不同的颜色给三棱柱 ${ABC} - {A}^{\prime }{B}^{\prime }{C}^{\prime }$ 的六个顶点涂色， 要求每个点涂一种颜色.
        
        \begin{center}
        \includegraphics[width=0.3\textwidth]{bodyxb/src/bo_d67fn3v7aajc739ac5e0_14_138_1900_413_202_0.jpg}
        \end{center}
        \begin{enumerate}
            \item 若每个底面的顶点涂色所使用的颜色不相同，则不同的涂色方法共有多少种？
            \item 若每条棱的两个端点涂不同的颜色，则不同的涂色方法共有多少种？
        \end{enumerate}
        \begin{jieda}
            \begin{enumerate}
                \item 由题得每个底面的顶点涂色所使用的颜色不相同，第一步， $A$ 、 $B$ 、 $C$ 三点所涂颜色各不相同的方法有 $P_{4}^{3} = {24}$ (种)，第二步， ${A}^{\prime }\text{ 、 }{B}^{\prime }\text{ 、 }{C}^{\prime }$ 三点所涂颜色各不相同的方法有 $P_{4}^{3} = {24}$ (种)， 所以由分步计数原理,不同的涂色方法共有 $P_{4}^{3} \cdot  P_{4}^{3} = {24} \times  {24} = {576}$ (种).
                
                \item 考虑平面$BCC'B'$的涂色情况
                \begin{dingautolist}{172}
                    \item $BCC'B'$涂四种不同的颜色，则方案数为$A_4^4 = 24$种. $A$可以从$B',C'$的涂色中选一个，$A'$可以从$B,C$的涂色中选一个，故一共$24\times2\times2 = 96$种.
                    \item $BCC'B'$涂三种不同的颜色，则等价于选择一组相对的顶点并从四个颜色中选择三个排列，共有$C_2^2 A_4^3 = 48$种，不妨设$B,C'$涂同一种颜色，此时若$A$涂第四种颜色，则$A'$仅有唯一选择，若$A$选择和$B'$一样的颜色，则$A'$有两种选择，故一共有$48\times(1+2) = 144$种.
                    \item $BCC'B'$涂两种不同的颜色，则有$A_4^2 = 12$种方案，对$A$和$A'$可以从剩下两种颜色中各选一个，则一共有$12\times2 = 24$种
                \end{dingautolist}
                综上，共有$96+144+24 = 264$种.
            \end{enumerate}
        \end{jieda}
        
    \end{enumerate}
    \item 排列、组合综合
    \begin{enumerate}
        \item 现有 8 名师生站成一排照相，其中老师 2 人，男学生 4 人，女学生 2 人，在下列情况下，各有多少种不同的站法?
        \begin{enumerate}
            \item 老师站在最中间，2名女学生分别在老师的两边且相邻，4 名男学生两边各 2 人；
            \item 4 名男学生互不相邻，男学生甲不能在两端；
            \item 2 名老师之间必要有男女学生各 1 人.
        \end{enumerate}
        \begin{jieda}
            \begin{enumerate}
                \item 由题意可得共 $P_{2}^{2}P_{2}^{2}P_{4}^{4} = 2 \times  2 \times  {24} = {96}$ 种不同的站法.
                
                \item 先排老师和女学生共有 $P_{4}^{4}$ 种站法，再排男学生甲有 ${{C}}_{3}^{1}$ 种站法，最后排剩余的 3 名男学生有 $P_{4}^{3}$ 种站法，所以共有 $P_{4}^{4}{{C}}_{3}^{1}P_{4}^{3} = {24} \times  3 \times  {24} = {1728}$ 种不同的站法.
                
                \item 先任选一男学生一女学生站两位老师中间，有 ${{C}}_{2}^{1}{{C}}_{4}^{1}P_{2}^{2}$ 种站法，两老师的站法有 $P_{2}^{2}$ 种，再将一男学生一女学生两位老师进行捆绑与剩余的 4 个人进行全排列有 $P_{5}^{5}$ 种，所以共有 ${{C}}_{2}^{1}{{C}}_{4}^{1}P_{2}^{2}P_{2}^{2}P_{5}^{5} = 2 \times  4 \times  2 \times  2 \times  {120} = {3840}$ 种不同的站法.
            \end{enumerate}
        \end{jieda}
        \item 某医疗小组有 4 名男性，2 名女性共 6 名医护人员，医护人员甲是其中一名.
        \begin{enumerate}
            \item 若从中任选 2 人参加 $A$ ， $B$ 两项救护活动，每人只能参加其中一项活动，每项活动都要有人参加，求医护人员甲不参加 $A$ 项救护活动的选法种数
            \item 这 6 名医护人员将去 3 个不同的地方参与医疗支援，每人只能去一地，每地有 2 人前往，若 2 名女性不能去往同一个地方，求不同的分配方案种数
        \end{enumerate}
        \begin{jieda}
            \begin{enumerate}
                \item 分两类: \ding{172} 甲参加 B 项救护活动，再从其余 5 人中选一人参加 A，选法数为 ${{C}}_{5}^{1} = 5$ ，  \ding{173} 甲不参加救护活动，则从其余 5 人中任选两人参加救护活动，选法数为 $P_{5}^{2} = {20}$ ，所以共有选法种数为 ${20} + 5 = {25}$ ;
                
                \item 分三步:
                
                第一步先安排两名女性医护人员有:$A_3^2$，
                
                第二步:安排两名女医护人员同去的男医护人员有:$A_4^2$，
                
                第三步:剩余两名男性医护人员去另外一地有: $C_{2}^{2}$ ，
                
                所以共有不同的分配方案数为: $P_{3}^{2}P_{4}^{2}{{C}}_{2}^{2} = {72}$ .
            \end{enumerate}
        \end{jieda}
        
        \item 从 1，2，3，4，5，6 中任取 5 个数字，随机填入如图所示的 5 个空格中.
        
        \begin{center}
        \begin{tabular}{|c|c|c|c|c|}
        \hline
        $A$ & $B$ & $C$ & $D$ & $E$ \\
        \cline{1-5}
        \hline
        \end{tabular}
        \end{center}
        \begin{enumerate}
            \item 若填入的 5 个数字中有 1 和 2 ，且 1 和 2 不能相邻，试问不同的填法有多少种？
            
            \item 若填入的 5 个数字中有 1 和 3，且区域 $A$ ， $B$ ， $C$ 中有奇数，试问不同的填法有多少种？
        \end{enumerate}
        \begin{jieda}
            \begin{enumerate}
                \item 首先从其它 4 个数中任选 3 个并作全排有 $P_{4}^{3} = {24}$ 种，3 个数中共有 4 个空,将 1 和 2 插入其中两个空有 $P_{4}^{2} = {12}$ 种，所以共有 ${24} \times  {12} = {288}$ 种填法.
                \item 若区域 $A,B,C$ 中无奇数，则其它三个数只能为 2、4、6 且在区域 $A,B,C$ 上，

                所以，共有 $P_{3}^{3}P_{2}^{2} = {12}$ 种，
                
                从 $2\text{ 、 }4\text{ 、 }5\text{ 、 }6$ 任选 3 个数有 ${{C}}_{4}^{3}$ 种,再把 5 个数全排有 $P_{5}^{5}$ 种,共有 ${{C}}_{4}^{3}P_{5}^{5} = {480}$ 种,
                
                综上,填入的 5 个数字中有 1 和 3 且区域 $A,B,C$ 中有奇数,共有 ${480} - {12} = {468}$ 种.
            \end{enumerate}
        \end{jieda}
        \item 某学校举行男子乒乓球团体赛，决赛比赛规则采用积分制， 两支决赛的队伍依次进行三场比赛，其中前两场为男子单打比赛，第三场为男子双打的比赛，每位出场队员在决赛中只能参加一场比赛. 某进入决赛的球队共有五名队员，现在需要提交该球队决赛的出场阵容，即三场比赛的出场的队员名单.
        \begin{enumerate}
            \item 一共有多少种不同的出场阵容？
            \item 若队员 $A$ 因为技术原因不能参加男子双打比赛，则一共有多少种不同的出场阵容？
        \end{enumerate}
        \begin{jieda}
            \begin{enumerate}
                \item 出场阵容可以分两步确定:
                
                第 1 步，从 5 名运动员中选择 2 人，分别参加前两场男单比赛，共有 $P_{5}^{2}$ 种；
                
                第 2 步，从剩下的 3 名运动员中选出两人参加男双比赛，共有 ${{C}}_{3}^{2}$ 种，
                
                根据分步乘法计数原理,不同的出场阵容种数为 $N = P_{5}^{2} \times  {{C}}_{3}^{2} = {60}$ .

                \item 队员 $A$ 不能参加男子双打比赛，有两类方案:
                
                第 1 类方案是队员 $A$ 不参加任务比赛，即除了队员 $A$ 之外的 4 人参加本次比赛，只需从 4 人中选出两人，
                
                分别取参加前两场单打比赛，共有 $P_{4}^{2}$ 种，剩余人员参加双打比赛；
                
                第 2 类方案是队员 $A$ 参加单打比赛，可以分 3 个步骤完成:
                
                第 1 步, 确定队员 $A$ 参加的是哪一场单打比赛, 共 2 种;
                
                第 2 步，从剩下 4 名队员中选择一名参加另一场单打比赛，共 4 种；
                
                第 3 步,从剩下的 3 名队员中,选出两人参加男双比赛,共有 ${{C}}_{3}^{2}$ 种,
                
                根据分步乘法计数原理,队员 $A$ 参加单打比赛的不同的出场阵容有 $2 \times  4 \times  {{C}}_{3}^{2}$ 种;
                
                根据分类加法计数原理,队员 $A$ 不参加男子双打比赛的不同的出场阵容种数为 $N = P_{4}^{2} + 2 \times  4 \times  {{C}}_{3}^{2} = {36}$ .
            \end{enumerate}
        \end{jieda}
        \item 从 $A,B,C$ 等 8 人中选出 5 人排成一排.
        \begin{enumerate}
            \item A必须在内，有多少种排法？
            \item $A,B,C$ 三人不全在内，有多少种排法？
            \item $A,B,C$ 都在内，且 $A,B$ 必须相邻， $C$ 与 $A,B$ 都不相邻，都多少种排法？
            \item A 不允许站排头和排尾，B 不允许站在中间(第三位)，有多少种排法？
        \end{enumerate}
        \begin{jieda}
            \begin{enumerate}
                \item 由题意，先从余下的 7 人中选 4 人共有 ${C}_{7}^{4}$ 种不同结果，再将这 4 人与 $A$ 进行全排列有 ${A}_{5}^{5}$ 种不同的排法,故由乘法原理可知共有 ${C}_{7}^{4}{A}_{5}^{5} = {4200}$ 种不同排法

                \item 从8人中任选5人排列共有 ${A}_{8}^{5}$ 种不同排法， $A,B,C$ 三人全在内有 ${C}_{5}^{2}{A}_{5}^{5}$ 种不同排法,由间接法可得 $A,B,C$ 三人不全在内共有 ${A}_{8}^{5} - {C}_{5}^{2}{A}_{5}^{5} = {5520}$ 种不同排法;
                
                \item 因 $A,B,C$ 都在内，所以只需从余下 5 人中选 2 人有 ${C}_{5}^{2}$ 种不同结果， $A,B$ 必须相邻，有 ${A}_{2}^{2}$ 种不同排法，由于 $C$ 与 $A$ ， $B$ 都不相邻，先将选出的 2 人进行全排列共有 ${A}_{2}^{2}$种不同排法,再将 $A\text{ 、 }B$ 这个整体与 $C$ 插入到选出的 2 人所产生的 3 各空位中有 ${A}_{3}^{2}$ 种不同排法,由乘法原理可得共有 ${C}_{5}^{2}{A}_{2}^{2}{A}_{2}^{2}{A}_{3}^{2} = {240}$ 种不同排法

                \item 分四类:
                
                第一类: 所选的 5 人无 $A$ 、 $B$ ，共有 ${A}_{6}^{5} = {720}$ 种排法;
                
                第二类: 所选的 5 人有 $A$ 、无 $B$ ，共有 ${C}_{6}^{4}{C}_{3}^{1}{A}_{4}^{4} = {1080}$ 种排法;
                
                第三类:所选的 5 人无 $A$ 、有 $B$ ，共有 ${C}_{6}^{4}{C}_{4}^{1}{A}_{4}^{4} = {1440}$ 种排法；
                
                第四类: 所选的 5 人有 $A\text{ 、 }B$ ,若 $A$ 排中间时,有 ${C}_{6}^{3}{A}_{4}^{4}$ 种排法,
                
                若 $A$ 不排中间时，有 ${C}_{6}^{3}{C}_{2}^{1}{C}_{3}^{1}{A}_{3}^{3}$ 种排法，共有 ${C}_{6}^{3}\left( {{A}_{4}^{4} + {C}_{2}^{1}{C}_{3}^{1}{A}_{3}^{3}}\right)  = {1200}$ 种排法；
                
                综上, 共有 4440 种不同排法.
            \end{enumerate}
        \end{jieda}
    \end{enumerate}

    \item 求二项展开式
    \begin{enumerate}
        \item 二项式 ${\left( x + 2\right) }^{3}$ 的展开式为\blkda[4]{${x}^{3} + 6{x}^{2} + {12x} + 8$}
        \begin{jieda}
            二项式 ${\left( x + 2\right) }^{3} = {{C}}_{3}^{0}{x}^{3} + {{C}}_{3}^{1}{x}^{2} \times  2 + {{C}}_{3}^{2}x \times  {2}^{2} + {{C}}_{3}^{3} \times  {2}^{3} = {x}^{3} + 6{x}^{2} + {12x} + 8$ .
        \end{jieda}
        
        \item 下列不属于${(x - 2)}^{3}$ 的展开式的项的是\dotfill{(\kh{B})}
        \xx{${x}^{3}$}{$6{x}^{2}$}{${12x}$}{$-8$}
        \begin{jieda}
            由二项式定理可知, ${\left( x - 2\right) }^{3} = {x}^{3} - 6{x}^{2} + {12x} - 8$ ,故 $6{x}^{2}$ 不是展开式的项.
        \end{jieda}
        \item 化简 ${\left( x + 1\right) }^{4} - 4{\left( x + 1\right) }^{3} + 6{\left( x + 1\right) }^{2} - 4\left( {x + 1}\right)  + 1$ 的结果为\blkda{${x}^{4}$}
        \begin{jieda}
            ${\left( x + 1\right) }^{4} - 4{\left( x + 1\right) }^{3} + 6{\left( x + 1\right) }^{2} - 4\left( {x + 1}\right)  + 1 = {\left( x + 1\right) }^{4} + {{C}}_{4}^{1}{\left( x + 1\right) }^{3} \times  \left( {-1}\right)  + {{C}}_{4}^{2}{\left( x + 1\right) }^{2} \times  {\left( -1\right) }^{2} + {{C}}_{4}^{3}\left( {x + 1}\right)  \times  {\left( -1\right) }^{3} + {\left( -1\right) }^{4} = {\left\lbrack  \left( x + 1\right)  - 1\right\rbrack  }^{4} = {x}^{4}$
        \end{jieda}
        \item $C_n^0+C_n^1+\cdots+C_n^n = $\blkda{$2^n$}
        \item ${{C}}_{2022}^{1} + 2{{C}}_{2022}^{2} + 3{{C}}_{2022}^{3} + 4{{C}}_{2022}^{4} + \cdots  + {2022}{{C}}_{2022}^{2022} =$ \blkda{$1011 \times  {2}^{2022}$}
        \begin{jieda}
            \begin{enumerate}
                \item 方法1 

                ${{C}}_{2022}^{1} + 2{{C}}_{2022}^{2} + 3{{C}}_{2022}^{3} + 4{{C}}_{2022}^{4} + \cdots  + {2022}{{C}}_{2022}^{2022} = 2022C_{2021}^{0} + 2022C_{2021}^{1} + \cdots + 2022C_{2021}^{2021} = 2022 \times (C_{2021}^{0} + C_{2021}^{1} + \cdots + C_{2021}^{2021}) = 2022 \times 2^{2021}$

                \item 方法2

                设 ${\left( 1 + x\right) }^{n} = {{C}}_{n}^{0} + {{C}}_{n}^{1}x + {{C}}_{n}^{2}{x}^{2} + {{C}}_{n}^{3}{x}^{3} + {{C}}_{n}^{4}{x}^{4} + \cdots  + {{C}}_{n}^{n}{x}^{n}$ ,
            
                两边求导数, $n{\left( 1 + x\right) }^{n - 1} = {{C}}_{n}^{1} + 2{{C}}_{n}^{2}x + 3{{C}}_{n}^{3}{x}^{2} + 4{{C}}_{n}^{4}{x}^{3} + \cdots  + n{{C}}_{n}^{n}{x}^{n - 1}$ ,
                
                令 $x = 1$ ,得 $n \cdot  {2}^{n - 1} = {{C}}_{n}^{1} + 2{{C}}_{n}^{2} + 3{C}_n^3 + 4{{C}}_{n}^{4} + \cdots  + n{{C}}_{n}^{n}$ ,
                
                取 $n = {2022}$ ，得 ${{C}}_{2022}^{1} + 2{{C}}_{2022}^{2} + 3{{C}}_{2022}^{3} + 4{{C}}_{2022}^{4} + \cdots  + {2022}{{C}}_{2022}^{2022} = {2022} \cdot  {2}^{2021} = {1011} \times  {2}^{2022}$.
            \end{enumerate}
        \end{jieda}

    \end{enumerate}
    
    \item 求展开式的特定项或特定项的系数
    \begin{enumerate}
        \item $(3x - 2)^4$ 展开式中的第 3 项为\blkda{${216}{x}^{2}$}
        \begin{jieda}
            由题意可知: ${\left( 3x - 2\right) }^{4}$ 展开式中的第 3 项为 ${{C}}_{4}^{2}{\left( 3x\right) }^{2} \cdot  {\left( -2\right) }^{2} = {216}{x}^{2}$ .
        \end{jieda}
        
        \item 二项式 ${\left( \dfrac{1}{x} - \sqrt{x}\right) }^{6}$ 的展开式中常数项为第\blkda{5}项.
        \begin{jieda}
            二项式 ${\left( \dfrac{1}{x} - \sqrt{x}\right) }^{6}$ 展开式的通项为:
            
            ${T}_{k + 1} = {{C}}_{6}^{k}{\left( \dfrac{1}{x}\right) }^{6 - k}{\left( -\sqrt{x}\right) }^{k} = {\left( -1\right) }^{k}{{C}}_{6}^{k}{x}^{k - 6}{x}^{\frac{k}{2}} = {\left( -1\right) }^{k}{{C}}_{6}^{k}{x}^{\frac{3k}{2} - 6}$
            
            令 $\dfrac{3k}{2} - 6 = 0$ ，解得， $k = 4$ ，所以二项式 ${\left( \dfrac{1}{x} - \sqrt{x}\right) }^{6}$ 的展开式中常数项为第 5 项 ${T}_{5} = {\left( -1\right) }^{4}{{C}}_{6}^{4} = {15}$ .
        \end{jieda}
        \item 二项式${\left( 2{x}^{3} - \dfrac{1}{\sqrt{x}}\right) }^{7}$的常数项为\blkda{$14$}
        \begin{jieda}
            二项式 ${\left( 2{x}^{3} - \dfrac{1}{\sqrt{x}}\right) }^{7}$ 的展开式的通项为 ${{C}}_{7}^{r}{\left( 2{x}^{3}\right) }^{7 - r}{\left( -\dfrac{1}{\sqrt{x}}\right) }^{r} = {{C}}_{7}^{r}{\left( -1\right) }^{r}{2}^{7 - r}{x}^{{21} - \frac{7}{2}r}$ 
            
            令 ${21} - \dfrac{7}{2}r = 0$ ，解得 $r = 6,{{C}}_{7}^{6} \times  {\left( -1\right) }^{6} \times  {2}^{1} = {14}$
        \end{jieda}
        \item ${\left( \dfrac{x}{2} - \sqrt{x}\right) }^{6}$ 的展开式中， ${x}^{4}$ 的系数为\blkda{$\dfrac{15}{4}$}
        \begin{jieda}
            由 ${T}_{k + 1} = {{C}}_{6}^{k}{\left( \dfrac{x}{2}\right) }^{6 - k}{\left( -\sqrt{x}\right) }^{k} = {\left( -1\right) }^{k}{\left( \dfrac{1}{2}\right) }^{6 - k}{{C}}_{6}^{k}{x}^{6 - \frac{k}{2}}$ ,
            
            当 $6 - \dfrac{k}{2} = 4$ ,解得 $k = 4$ ,
            
            所以 ${x}^{4}$ 的系数为 ${\left( -7\right) }^{4}{\left( \dfrac{1}{2}\right) }^{2}{{C}}_{6}^{4} = \dfrac{1}{4} \times  {15} = \dfrac{15}{4}$
        \end{jieda}
        \item 若 ${\left( 2x + \dfrac{1}{\sqrt[3]{x}}\right) }^{n}$ 的展开式中所有项系数和为 81,则该展开式的常数项为\blkda{8}
        \begin{jieda}
            在 ${\left( 2x + \dfrac{1}{\sqrt[3]{x}}\right) }^{n}$ 的二项展开式中,令 $x = 1$ 得所有项的系数和为 ${3}^{n} = {81}$ ,解得 $n = 4$，于是得 ${\left( 2x + \dfrac{1}{\sqrt[3]{x}}\right) }^{4}$ 展开式的通项为 ${T}_{r + 1} = {C}_{4}^{r}{\left( 2x\right) }^{4 - r} \cdot  {\left( \dfrac{1}{\sqrt[3]{x}}\right) }^{r} = {2}^{4 - r} \cdot  {C}_{4}^{r}{x}^{4 - \frac{4}{3}},r \in  N,r \leq  4$，令 $4 - \dfrac{4}{3}r = 0$ ,得 $r = 3$ ,常数项为 $2{C}_{4}^{3} = 8$ .
        \end{jieda}
    \end{enumerate}

    
    
    \item 根据二项式的特定项求值
    \begin{enumerate}
        \item 若${\left( \sqrt[5]{x} - \dfrac{1}{x}\right) }^{n}$展开式的第四项是常数，则 $n$ 的值为\blkda{$18$}
        \begin{jieda}
            ${\left( \sqrt[5]{x} - \dfrac{1}{x}\right) }^{n}$ 展开式的通项为 ${T}_{k + 1} = {C}_{n}^{k}{\left( {x}^{\frac{1}{5}}\right) }^{n - k}{\left( -1\right) }^{k}{x}^{-k} = {C}_{n}^{k}{\left( -1\right) }^{k}{x}^{\frac{n}{5} - \frac{6k}{5}}$ ,
            
            令 $k = 3$ 可得 ${T}_{4} = {C}_{n}^{3}{\left( -1\right) }^{3}{x}^{\frac{n}{5} - \frac{18}{5}}$ 为常数项,可得 $\dfrac{n}{5} - \dfrac{18}{5} = 0$ ,可得 $n = {18}$ 
        \end{jieda}

        \item 若 ${\left( x + a\right) }^{5}$ 的展开式中 ${x}^{2}$ 的系数是 80，则实数 $a$ 的值是\blkda{2}
        \begin{jieda}
            ${\left( x + a\right) }^{5}$ 展开式的通项为 ${{C}}_{5}^{r}{x}^{5 - r}{a}^{r}$

            令 $5 - r = 2 \Longrightarrow  r = 3$，可得 ${x}^{2}$ 系数为 ${{C}}_{5}^{3}{a}^{3} = {10}{a}^{3} = {80}$，可得 $a = 2$ .
        \end{jieda}
        
        \item 已知 $a > 0$ ，若 ${\left( {x}^{2} + \dfrac{a}{x}\right) }^{6}$ 的展开式中，常数项等于 240, 则 $a =$\blkda{$2$}
        \begin{jieda}
            由二项展开式的通项公式可得 ${T}_{r + 1} = {{C}}_{6}^{r}{\left( {x}^{2}\right) }^{6 - r}{\left( \dfrac{a}{x}\right) }^{r} = {a}^{r}{{C}}_{6}^{r}{x}^{{12} - {3r}}$
            
            令 ${12} - {3r} = 0$ ,解得 $r = 4$，即常数项为 ${T}_{5} = {a}^{4}{{C}}_{6}^{4} = {240}$ ，解得 $a = 2$ .
        \end{jieda}
        \item 已知 ${\left( 2{x}^{3} + \dfrac{1}{x}\right) }^{n}$ 的展开式的常数项是第七项,则正整数 $n$ 的值为\blkda{8}
        \begin{jieda}
            由二项式展开式的通项公式可知展开式的第七项为 ${T}_{6 + 1} = {{C}}_{n}^{6} \cdot  {\left( 2{x}^{3}\right) }^{n - 6} \cdot  {x}^{-6} = {2}^{n - 6} \cdot  {{C}}_{n}^{6} \cdot  {x}^{{3n} - {24}}$，又因为第七项为常数,所以 ${3n} - {24} = 0,n = 8$
        \end{jieda}
        \item 已知在 ${\left( \sqrt[3]{x} - \dfrac{3}{\sqrt[3]{x}}\right) }^{n}$ 的展开式中，第 6 项为常数项.
        \begin{enumerate}
            \item 求$n$
            \item 求含 $x^{2}$ 项的系数
            \item 求展开式中所有的有理项.
        \end{enumerate}
        \begin{jieda}
            \begin{enumerate}
                \item 通项公式为 ${T}_{r + 1} = {C}_{n}^{r} \cdot  {x}^{\frac{n - r}{3}} \cdot  {\left( -3\right) }^{r} \cdot  {x}^{\frac{r}{-3}} = {\left( -3\right) }^{r} \cdot  {C}_{n}^{r} \cdot  {x}^{\frac{n - {2r}}{3}}$ .
                
                因为第 6 项为常数项,所以 $r = 5$ 时,有 $\dfrac{n - {2r}}{3} = 0$ ,解得 $n = {10}$ .
                
                \item 由(1)可知 $n = {10}$ ，令 $\dfrac{n - {2r}}{3} = 2$ ，解得 $r = 2$ . 所以含 ${x}^{2}$ 项的系数为 ${\left( -3\right) }^{2} \cdot  {C}_{10}^{2} = {405}$ .
                
                \item 由题意可知, $\left\{  \begin{matrix} \dfrac{{10} - {2r}}{3} \in  Z \\  0 \leq  r \leq  {10}, \\  r \in  N \end{matrix}\right.$ 则 $r$ 可能的取值为2,5,8.
                
                所以第 3 项，第 6 项，第 9 项为有理项，分别为 ${C}_{10}^{2} \cdot  {\left( -3\right) }^{2} \cdot  {x}^{2},{C}_{10}^{5} \cdot  {\left( -3\right) }^{5},{C}_{10}^{8} \cdot  {\left( -3\right) }^{8} \cdot  {x}^{-2}$.
            \end{enumerate}
        \end{jieda}
    \end{enumerate}

    \item 求展开式中系数最大(小)的项
    \begin{enumerate}
        \item 已知“$(1+3x)^n$”的展开式中前三项的二项式系数和为 79，则展开式中系数最大的项为第\blkda{10}项.
        \begin{jieda}
            ${\left( 1 + 3x\right) }^{n}$ 的展开式中前三项的二项式系数和为 ${{C}}_{n}^{0} + {{C}}_{n}^{1} + {{C}}_{n}^{2} = 1 + n + \dfrac{n\left( {n - 1}\right) }{2} = {79}$
            
            整理可得 ${n}^{2} + n - {156} = 0,\because n \geq  2$ 且 $n \in  {N}^{ * }$ ,解得 $n = {12}$
            
            ${\left( 1 + 3x\right) }^{12}$ 的展开式通项为 ${T}_{k + 1} = {{C}}_{12}^{k} \cdot  {\left( 3x\right) }^{k} = {{C}}_{12}^{k} \cdot  {3}^{k}{x}^{k}\left( {k = 0,1,2,\cdots ,{12}}\right)$
            
            设展开式中第 $r + 1$ 项的系数最大,则 $\left\{  \begin{array}{l} {{C}}_{12}^{r} \cdot  {3}^{r} \geq  {{C}}_{12}^{r + 1} \cdot  {3}^{r + 1} \\  {{C}}_{12}^{r} \cdot  {3}^{r} \geq  {{C}}_{12}^{r - 1} \cdot  {3}^{r - 1} \end{array}\right.$ ,
            
            即 $\left\{  {\begin{array}{l} \dfrac{{12}!}{r!\left( {{12} - r}\right) !} \cdot  {3}^{r} \geq  \dfrac{{12}!}{\left( {r + 1}\right) ! \cdot  \left( {{11} - r}\right) !} \cdot  {3}^{r + 1} \\  \dfrac{{12}!}{r!\left( {{12} - r}\right) !} \cdot  {3}^{r} \geq  \dfrac{{12}!}{\left( {r - 1}\right) !\left( {{13} - r}\right) !} \cdot  {3}^{r - 1} \end{array}\text{ ,解得 }\dfrac{35}{4} \leq  r \leq  \dfrac{39}{4}}\right.$ ,
            
            因为 $r \in  N$ ,故 $r = 9$ ,因此,展开式中系数最大的项为第 10 项.
        \end{jieda}
        
        \item 已知 ${\left( x + \dfrac{1}{2\sqrt{x}}\right) }^{n}$ 的展开式中仅第 4 项的二项式系数最大，则展开式中系数最大的项是第\blkda{3}项
        \begin{jieda}
            由题意二项式系数仅 ${{C}}_{n}^{3}$ 最大,故 $n = 6$ ,
            
            所以二项式为 ${\left( x + \dfrac{1}{2\sqrt{x}}\right) }^{6}$ ,其通项公式为 ${{T}}_{r + 1} = {{C}}_{6}^{r}{x}^{6 - r}{\left( \dfrac{1}{2\sqrt{x}}\right) }^{r} = {2}^{-r}{{C}}_{6}^{r}{x}^{6 - \frac{3}{2}r},r = 0,1,2,3,4,5,6$ ,
            
            设二项式展开式中第 $r + 1$ 项的系数最大,则有 $\left\{  \begin{array}{l} {2}^{-r}{{C}}_{6}^{r} \geq  {2}^{-\left( {r - 1}\right) }{{C}}_{6}^{r - 1} \\  {2}^{-r}{{C}}_{6}^{r} \geq  {2}^{-\left( {r + 1}\right) }{{C}}_{6}^{r + 1} \end{array}\right. \Longrightarrow  \left\{  \begin{array}{l} r \leq  \dfrac{7}{3} \\  r \geq  \dfrac{4}{3} \end{array}\right.$ ,即 $\dfrac{4}{3} \leq  r \leq  \dfrac{7}{3}$ . 故 $r = 2$ ,经经验符合题意,
            
            所以展开式中系数最大的项是第 3 项.
        \end{jieda}
        
        \item 已知 ${\left( x + \dfrac{1}{2{x}^{2}}\right) }^{n},\left( {n \geq  4,n \in  {{N}}^{ * }}\right)$ 的展开式中，第 5 项与第 3 项的二项式系数之比为 15:2.
        \begin{enumerate}
            \item 求 $n$ 的值及展开式中的常数项
            \item 求展开式中系数最大项
        \end{enumerate}
        \begin{jieda}
            \begin{enumerate}
                \item 由题意，可得二项式 ${\left( x + \dfrac{1}{2{x}^{2}}\right) }^{n}$ 展开式的通项为${T}_{r + 1} = {{C}}_{n}^{r}{x}^{n - r} \cdot  {\left( \dfrac{1}{2{x}^{2}}\right) }^{r} = {2}^{-r} \cdot  {{C}}_{n}^{r}{x}^{n - {3r}}$
                
                因为第 5 项与第 3 项的二项式系数之比为 15:2 ,可得 $\dfrac{{{C}}_{n}^{4}}{{{C}}_{n}^{2}} = \dfrac{15}{2}$ ,即 $\dfrac{\left( {n - 2}\right) \left( {n - 3}\right) }{12} = \dfrac{15}{2}$ ,解得 $n = {12}$ (负值舍)
                
                所以 ${T}_{r + 1} = {2}^{-r} \cdot  {{C}}_{12}^{r}{x}^{{12} - {3r}}$ ,令 ${12} - {3r} = 0$ ,得 $r = 4$ ,所以展开式的常数项为 ${T}_{5} = {2}^{-4} \cdot  {{C}}_{12}^{4} = \dfrac{495}{16}$ .
                
                \item 设展开式中第 $r + 1$ 项的系数最大，
                
                则 $\left\{  \begin{array}{l} {{C}}_{12}^{r} \cdot  {2}^{-r} \geq  {{C}}_{12}^{r + 1} \cdot  {2}^{-r - 1} \\  {{C}}_{12}^{r} \cdot  {2}^{-r} \geq  {{C}}_{12}^{r - 1} \cdot  {2}^{-r + 1} \end{array}\right.$ ,可得 $\left\{  \begin{array}{l} 2\left( {r + 1}\right)  \geq  {12} - r \\  {13} - r \geq  {2r} \end{array}\right.$ ,解得 $\dfrac{10}{3} \leq  r \leq  \dfrac{13}{3}$ ,
                
                因为 $r \in  {{N}}^{ * }$ ,所以 $r = 4$ ,所以系数最大的项为 ${T}_{5} = {2}^{-4} \cdot  {{C}}_{12}^{4}{x}^{0} = \dfrac{495}{16}$ .
            \end{enumerate}
        \end{jieda}
        
        \item 在二项式 ${\left( \sqrt[3]{x} + \dfrac{2}{x}\right) }^{n}$ 的展开式中
        \begin{enumerate}
            \item 若第 4 项的系数与第 6 项的系数比为 5 :6，求展开式中的有理项
            \item 若展开式中只有第 5 项的二项式系数最大，求展开式中系数最大的项.
        \end{enumerate}
        \begin{jieda}
            \begin{enumerate}
                \item 由题意得 ${{C}}_{n}^{3} \cdot  {2}^{3} : {{C}}_{n}^{5} \cdot  {2}^{5} = 5 : 6$ 
                
                $\therefore 6{{C}}_{n}^{3} = {20}{{C}}_{n}^{5}$ ,即 ${n}^{2} - {7n} + 6 = 0$ ,解得 $n = 6$ 或 $n = 1$ (舍).
                
                ${T}_{k + 1} = {{C}}_{n}^{k} \cdot  {\left( \sqrt[3]{x}\right) }^{n - k} \cdot  {\left( \frac{2}{x}\right) }^{k} = {{C}}_{6}^{k} \cdot  {2}^{k} \cdot  {x}^{\frac{6 - {4k}}{3}},k = 0,1,2,\ldots 6$，所以 $k = 0,3,6$ 时为有理项，即展开式中的有理项为: ${T}_{1} = {x}^{2},{T}_{4} = \frac{160}{{x}^{2}},{T}_{7} = \frac{64}{{x}^{6}}$ .
                \item 因为展开式中只有第 5 项的二项式系数最大，所以 $n = 8$ ，设第 $r + 1$ 项的展开式系数最大，则$\begin{array}{l} \left\{  \begin{array}{l} {{C}}_{8}^{r} \cdot  {2}^{r} \geq  {{C}}_{8}^{r - 1} \cdot  {2}^{r - 1}, \\  {{C}}_{8}^{r} \cdot  {2}^{r} \geq  {{C}}_{8}^{r + 1} \cdot  {2}^{r + 1} \end{array}\right. \text{ ,解得 }5 \leq  r \leq  6. \end{array}$

                所以展开式中系数最大项为: ${T}_{6} = {{C}}_{8}^{5}{\left( \sqrt[3]{x}\right) }^{3}{\left( \dfrac{2}{x}\right) }^{5} = {1792}{x}^{-4},{T}_{7} = {{C}}_{8}^{6}{\left( \sqrt[3]{x}\right) }^{2}{\left( \dfrac{2}{x}\right) }^{6} = {1792}{x}^{-\frac{16}{3}}$ .
            \end{enumerate}
        \end{jieda}

        \item $({2x} - y)^{6}$ 的展开式中系数最小的项和二项式系数最大的项分别为 \dotfill{(\kh{B})}
        \xx{第 1 项和第 3 项}{第 2 项和第 4 项}{第 3 项和第 1 项}{第 4 项和第 2 项}
        \begin{jieda}
            ${\left( 2x - y\right) }^{6}$ 的展开式的通项为 ${T}_{k + 1} = {\left( -1\right) }^{k}{2}^{6 - k}{{C}}_{6}^{k}{x}^{6 - k}{y}^{k}$
            
            当 $k$ 取奇数时,系数为负值：当 $k = 1$ 时， ${T}_{2} =  - {{192}{x}^{5}}y$ ，当 $k = 3$ 时， ${T}_{4} =  - {{160}{x}^{3}}{y}^{3}$ ，当 $k = 5$ 时， ${T}_{6} =  - {12x}{y}^{5}$ ，所以第 2 项的系数最小;
            
            因为 ${\left( 2x - y\right) }^{6}$ 的展开式有 7 项,所以中间一项的二项式系数最大,即第 $\dfrac{6}{2} + 1 = 4$ 项的二项式系数最大.
        \end{jieda}
        \item ${\left( {x}^{2} + \dfrac{1}{2x}\right) }^{6}$ 的二项展开式中系数最大的项为第\blkda{3}项
        \begin{jieda}
            ${\left( {x}^{2} + \dfrac{1}{2x}\right) }^{6}$ 的展开式通项公式为 ${T}_{r + 1} = {{C}}_{6}^{r}{\left( {x}^{2}\right) }^{6 - r}{\left( \dfrac{1}{2x}\right) }^{r} = \dfrac{{{C}}_{6}^{r}}{{2}^{r}}{x}^{{12} - {3r}}$ ,

            设第 $r + 1$ 项为系数最大的项,则有 $\left\{  \begin{array}{l} \dfrac{{{C}}_{6}^{r}}{{2}^{r}} \geq  \dfrac{{{C}}_{6}^{r + 1}}{{2}^{r + 1}} \\  \dfrac{{{C}}_{6}^{r}}{{2}^{r}} \geq  \dfrac{{{C}}_{6}^{r - 1}}{{2}^{r - 1}} \end{array}\right.$ ,解得 $\dfrac{4}{3} \leq  r \leq  \dfrac{7}{3}$ ,即 $r = 2$ .
        \end{jieda}
        \item 已知$(1-x)^{n}$ 的展开式中，仅有第 5 项的二项式系数最大，则展开式中系数的最小值为\blkda{$-56$}
        \begin{jieda}
            因为 ${\left( 1 - x\right) }^{n}$ 的展开式中,仅有第 5 项的二项式系数最大,所以 $n = 8$ , 所以展开式的通项公式为 ${T}_{r + 1} = {{C}}_{8}^{r}{l}^{8 - r}{\left( -x\right) }^{r} = {{C}}_{8}^{r}{\left( -l\right) }^{r}{x}^{r}$ ,要使展开式中系数的最小值,则 $r$ 为奇数, 取值为1,3,5,7，所以当 $r = 3$ 或 5 时，系数 ${{C}}_{8}^{r}{\left( -1\right) }^{r}$ 最小，则展开式中系数的最小值为 ${{C}}_{8}^{3}{\left( -1\right) }^{3} = {{C}}_{8}^{5}( -$ 1) ${}^{5} =  - {56}$
        \end{jieda}
        
        \item 已知 ${\left( x + \dfrac{1}{2\sqrt{x}}\right) }^{n}$ 的展开式中仅第 4 项的二项式系数最大,则展开式中系数最大的项是第 \blkda{3} 项
        \begin{jieda}
            由题意二项式系数仅 ${{C}}_{n}^{3}$ 最大,故 $n = 6$，所以二项式为 ${\left( x + \dfrac{1}{2\sqrt{x}}\right) }^{6}$ ,其通项公式为 ${{T}}_{r + 1} = {{C}}_{6}^{r}{x}^{6 - r}{\left( \dfrac{1}{2\sqrt{x}}\right) }^{r} = {2}^{-r}{{C}}_{6}^{r}{x}^{6 - \frac{3}{2}r},r = 0,1,2,3,4,5,6$ ,
            
            设二项式展开式中第 $r + 1$ 项的系数最大,则有 $\left\{  \begin{array}{l} {2}^{-r}{{C}}_{6}^{r} \geq  {2}^{-\left( {r - 1}\right) }{{C}}_{6}^{r - 1} \\  {2}^{-r}{{C}}_{6}^{r} \geq  {2}^{-\left( {r + 1}\right) }{{C}}_{6}^{r + 1} \end{array}\right.$ ,
            
            $\Longrightarrow  \left\{  \begin{array}{l} r \leq  \dfrac{7}{3} \\  r \geq  \dfrac{4}{3} \end{array}\right.$ ,即 $\dfrac{4}{3} \leq  r \leq  \dfrac{7}{3}$ ,故 $r = 2$ ,经检验符合题意，所以展开式中系数最大的项是第 3 项.
        \end{jieda}
    \end{enumerate}
    
    \item 多项式积的展开式中的特定项问题
    \begin{enumerate}
        \item 在 $\left( {x - y}\right) {\left( x + y\right) }^{5}$ 的展开式中，含有 ${x}^{2}{y}^{4}$ 项的系数为\blkda{$-5$}
        \begin{jieda}
            由题意,在 $\left( {x - y}\right) {\left( x + y\right) }^{5}$ 的展开式中,
            
            其中 ${x}^{2}{y}^{4}$ 项为 $x \cdot  {{C}}_{5}^{4}x \cdot  {y}^{4} + \left( {-y}\right)  \cdot  {{C}}_{5}^{3}{x}^{2} \cdot  {y}^{3} =  - 5{x}^{2}{y}^{4}$ ,
            
            所以 ${x}^{2}{y}^{4}$ 项的系数为 $-5$
        \end{jieda}
        
        \item 在 ${\left( 1 + x\right) }^{6}{\left( 1 + \dfrac{1}{y}\right) }^{4}$ 的展开式中， $\dfrac{{x}^{3}}{{y}^{2}}$ 的系数为\blkda{$120$}
        \begin{jieda}
            ${\left( 1 + x\right) }^{6}$ 的展开式的二项式通项为 ${T}_{r + 1} = {{C}}_{6}^{r}{x}^{r}$ ,令 $r = 3$ ,则 ${T}_{4} = {{C}}_{6}^{3}{x}^{3} = {20}{x}^{3}$ .
            
            ${\left( 1 + \dfrac{1}{y}\right) }^{4}$ 的展开式的二项式通项为 ${T}_{k + 1} = {{C}}_{4}^{k}{\left( \dfrac{1}{y}\right) }^{k}$，令 $k = 2$ ,可得 ${T}_{3} = {{C}}_{4}^{2}{\left( \dfrac{1}{y}\right) }^{2} = \dfrac{6}{{y}^{2}}$ .
            
            故 $\dfrac{{x}^{3}}{{y}^{2}}$ 项的系数为 ${20} \times  6 = {120}$ .
        \end{jieda}
        
        \item 已知 $(2-x)(1+ax)^{3}$ 的展开式中各项系数之和为 27，则展开式中 $x^{2}$ 项的系数为\blkda{$18$}
        \begin{jieda}
            由题意 $\left( {2 - 1}\right)  \times  {\left( 1 + a\right) }^{3} = {27}$ ,解得 $a = 2$ ,
            
            所以 $\left( {2 - x}\right) {\left( 1 + 2x\right) }^{3}$ 展开式中 ${x}^{2}$ 项的系数为 $2{{C}}_{3}^{2}{2}^{2} - {{C}}_{3}^{1}2 = {18}$ .
        \end{jieda}
        
        \item 若 $\left( {x + \dfrac{m}{x}}\right) {\left( x - \dfrac{1}{x}\right) }^{5}$ 的展开式中常数项是 20,则 $m =$\blkda{3}
        \begin{jieda}
            $\left( {x + \dfrac{m}{x}}\right) {\left( x - \dfrac{1}{x}\right) }^{5} = x{\left( x - \dfrac{1}{x}\right) }^{5} + \dfrac{m}{x}{\left( x - \dfrac{1}{x}\right) }^{5}$ ,
            
            ${\left( x - \dfrac{1}{x}\right) }^{5}$ 的展开式的通项公式为 ${T}_{r + 1} = {{C}}_{5}^{r}{x}^{5 - r}{\left( -\dfrac{1}{x}\right) }^{r} = {{C}}_{5}^{r}{\left( -1\right) }^{r}{x}^{5 - {2r}}$ ,
            
            令 $5 - {2r} =  - 1$ ，解得 $r = 3$ ，则 $x{\left( x - \dfrac{1}{x}\right) }^{5}$ 的展开式的常数项为 $- {{C}}_{5}^{3} =  - {10}$ ；
            
            令 $5 - {2r} = 1$ ，解得 $r = 2$ ，则 $\dfrac{m}{x}{\left( x - \dfrac{1}{x}\right) }^{5}$ 的展开式的常数项为 $m{{C}}_{5}^{2} = {10m}$ ，
            
            因为 $\left( {x + \dfrac{m}{x}}\right) {\left( x - \dfrac{1}{x}\right) }^{5}$ 的展开式中常数项是 20,所以 ${10m} - {10} = {20}$ ,解得 $m = 3$ .
        \end{jieda}

        
        \item 若 $\left( {x + 2}\right) {\left( \dfrac{1}{x} - ax\right) }^{7}$ 展开式的常数项等于 $-280$,则 $a =$ \blkda{2}
        \begin{jieda}
            ${\left( \dfrac{1}{x} - ax\right) }^{7}$ 展开式的通项公式为: ${T}_{k + 1} = {C}_{7}^{k}{\left( \dfrac{1}{x}\right) }^{7 - k}{\left( -ax\right) }^{k} = {C}_{7}^{k}{\left( -a\right) }^{k}{x}^{{2k} - 7}$，${\left( \dfrac{1}{x} - ax\right) }^{7}$ 的展开式无常数项
            
            当 $k = 3$ 时， $\dfrac{1}{x}$ 项的系数为: ${C}_{7}^{3}{\left( -a\right) }^{3}$，所以 $\left( {x + 2}\right) {\left( \dfrac{1}{x} - ax\right) }^{7}$ 展开式的常数项为: ${C}_{7}^{3}{\left( -a\right) }^{3} =  - {280}$ ，解得: $a = 2$
        \end{jieda}
        
    \end{enumerate}
    
    \item 用赋值法求系数和问题
    \begin{enumerate}
        \item 已知 ${\left( x + 2\right) }^{11} = {a}_{0} + {a}_{1}x + \cdots  + {a}_{11}{x}^{11}$ ,则 ${a}_{1} + {a}_{3} + \cdots  + {a}_{11} =$\blkda{$\dfrac{{3}^{11} - 1}{2}$}
        \begin{jieda}
            令 $x = 1$ ,得 ${a}_{0} + {a}_{1} + \cdots  + {a}_{11} = {3}^{11}$
            
            令 $x =  - 1$ ,得 ${a}_{0} - {a}_{1} + \cdots  - {a}_{11} = 1$

            两式做差，得 $2\left( {{a}_{1} + {a}_{3} + \cdots  + {a}_{11}}\right)  = {3}^{11} - 1$ ，即 ${a}_{1} + {a}_{3} + \cdots  + {a}_{11} = \dfrac{{3}^{11} - 1}{2}$ .
        \end{jieda}
        
        \item 若 ${\left( 1 - 2x\right) }^{2024} = {a}_{0} + {a}_{1}x + {a}_{2}{x}^{2} + \cdots  + {a}_{2024}{x}^{2024}$ ,则 $\left| {a}_{0}\right|  +  \left| {a}_{1}\right|  + \left| {a}_{2}\right|  + \cdots \left| {a}_{2024}\right|  =$\blkda{${3}^{2024}$}
        \begin{jieda}
            ${\left( 1 - 2x\right) }^{2024}$ 的展开式的通项公式为 ${T}_{r + 1} = {{C}}_{2020}^{r} \cdot  {\left( -2x\right) }^{r}\left( {r = 0,1,2,\cdots ,{2024}}\right)$ ,
            
            结合 ${\left( 1 - 2x\right) }^{2024} = {a}_{0} + {a}_{1}x + {a}_{2}{x}^{2} + \cdots  + {a}_{2024}{x}^{2024}$ ,知 ${a}_{1},{a}_{3},{a}_{5},\cdots ,{a}_{2023}$ 均为负值,
            
            $\therefore \left| {a}_{0}\right|  + \left| {a}_{1}\right|  + \left| {a}_{2}\right|  + \cdots \left| {a}_{2024}\right|  = {a}_{0} - {a}_{1} + {a}_{2} - {a}_{3} + \cdots  + {a}_{2022} - {a}_{2023} + {a}_{2024}$ ,
            
            令 $x =  - 1,$ 得 ${3}^{2024} = {a}_{0} - {a}_{1} + {a}_{2} - {a}_{3} + \cdots  + {a}_{2022} - {a}_{2023} + {a}_{2024}$ ,
            
            故 $\left| {a}_{0}\right|  + \left| {a}_{1}\right|  + \left| {a}_{2}\right|  + \cdots \left| {a}_{2024}\right|  = {3}^{2024}$ 
        \end{jieda}
        
        \item 已知 $\left( {1 + {2x}}\right) {\left( 2 - x\right) }^{6} = {a}_{0} + {a}_{1}x + {a}_{2}{x}^{2} + \cdots  + {a}_{7}{x}^{7}$ ，那么 $\dfrac{{a}_{0} + {a}_{2} + {a}_{4}}{{a}_{1} + {a}_{3} + {a}_{5} + {a}_{7}}$ 的值为\blkda{$-\dfrac{170}{183}$}
        \begin{jieda}
            因为 $\left( {1 + {2x}}\right) {\left( 2 - x\right) }^{6} = {a}_{0} + {a}_{1}x + {a}_{2}{x}^{2} + \cdots  + {a}_{7}{x}^{7}$ ,
            
            令 $x = 1$ 可得 ${a}_{0} + {a}_{1} + {a}_{2} + {a}_{3} + {a}_{4} + {a}_{5} + {a}_{6} + {a}_{7} = 3$ ,
            
            令 $x =  - 1$ 可得 ${a}_{0} - {a}_{1} + {a}_{2} - {a}_{3} + {a}_{4} - {a}_{5} + {a}_{6} - {a}_{7} =  - {3}^{6} =  - {729}$ ,
            
            所以 ${a}_{0} + {a}_{2} + {a}_{4} + {a}_{6} = \dfrac{3 - {729}}{2} =  - {363},{a}_{1} + {a}_{3} + {a}_{5} + {a}_{7} = \dfrac{3 - \left( {-{729}}\right) }{2} = {366}$ ,
            
            又 $\left( {1 + {2x}}\right) {\left( 2 - x\right) }^{6} = {\left( 2 - x\right) }^{6} + {2x}{\left( 2 - x\right) }^{6}$ ,
            
            其中 ${\left( 2 - x\right) }^{6}$ 展开式的通项为 ${T}_{r + 1} = {2}^{6 - r}{{C}}_{6}^{r}{\left( -x\right) }^{r} = {2}^{6 - r}{\left( -1\right) }^{r}{{C}}_{6}^{r}{x}^{r}\left( {0 \leq  r \leq  6\text{ 且 }r \in  {N}}\right)$ ,
            
            所以 ${a}_{6} = {2}^{0}{{C}}_{6}^{6} + 2 \times  {2}^{1}{{C}}_{6}^{5} \times  {\left( -1\right) }^{5} =  - {23}$ ,
            
            所以 ${a}_{0} + {a}_{2} + {a}_{4} =  - {363} - {a}_{6} =  - {363} + {23} =  - {340}$ ,
            
            所以 $\dfrac{{a}_{0} + {a}_{2} + {a}_{4}}{{a}_{1} + {a}_{3} + {a}_{5} + {a}_{7}} = \dfrac{-{340}}{366} =  - \dfrac{170}{183}$ .
        \end{jieda}
        \item 已知 ${(2x - 1)}^{2025} = {a}_{0} + {a}_{1}x + {a}_{2}{x}^{2} + \cdots  + {a}_{2025}{x}^{2025}$ ，则\dotfill{(\kh{C})}
        \xx{${a}_{0} = 1$}{${a}_{1} + {a}_{2} + \cdots  + {a}_{2025} = 1$}{${a}_{1} + {a}_{3} + {a}_{5} + \cdots  + {a}_{2025} = \dfrac{1 + {3}^{2025}}{2}$}{${a}_{0} + \dfrac{{a}_{1}}{2} + \dfrac{{a}_{2}}{{2}^{2}} + \dfrac{{a}_{3}}{{2}^{3}} + \cdots  + \dfrac{{a}_{2025}}{{2}^{2025}} = 1$}
        \begin{jieda}
            令 $x = 0$ ,得 ${a}_{0} = {\left( -1\right) }^{2025} =  - 1$ ,故 A 不正确;
            
            令 $x = 1$ ,得 ${a}_{0} + {a}_{1} + {a}_{2} + \cdots  + {a}_{2025} = 1$ ,所以 ${a}_{1} + {a}_{2} + \cdots  + {a}_{2025} = 1 + 1 = 2$ ,故 B 不正确;
            
            令 $x =  - 1,$ 得 ${a}_{0} - {a}_{1} + {a}_{2} - \cdots  - {a}_{2025} = {\left( -3\right) }^{2025} =  - {3}^{2025}$ ,
            
            所以 ${a}_{1} + {a}_{3} + {a}_{5} + \cdots  + {a}_{2025} = \dfrac{1 + {3}^{2025}}{2}$ ,故 ${C}$ 正确;
            
            令 $x = \dfrac{1}{2}$ ，得 ${a}_{0} + \dfrac{{a}_{1}}{2} + \dfrac{{a}_{2}}{{2}^{2}} + \dfrac{{a}_{3}}{{2}^{3}} + \cdots  + \dfrac{{a}_{2025}}{{2}^{2025}} = {\left( 2 \times  \dfrac{1}{2} - 1\right) }^{2025} = 0$ ，所以 D 不正确.
        \end{jieda}
    \end{enumerate}
    
    \item 三项展开式的系数问题
    \begin{enumerate}
        \item  ${\left( {x}^{2} - x + 2\right) }^{6}$ 的展开式中， ${x}^{5}$ 的系数为\blkda{$- {732}$}
        \begin{jieda}
            ${\left( {x}^{2} - x + 2\right) }^{6}$ 展开式中的项,可看做有 6 个盒子,每个盒子中 ${x}^{2}, - x,2$ 三个元素,从每个盒子中取出一个元素,再将所得的元素相乘;要得到 ${x}^{5}$ :
            \begin{dingautolist}{172}
                \item 可以取 1 个 ${x}^{2}$ ，3 个 $- x$ ，2 个 2，则为 ${{C}}_{6}^{1}{x}^{2} \cdot  {{C}}_{5}^{3}{\left( -x\right) }^{3} \cdot  {{C}}_{2}^{2} \times  {2}^{2} =  - {240}{x}^{5}$
                \item 可以取2个 ${x}^{2}$ ，1个 $- x$ ，3个2，则为 ${{C}}_{6}^{2}{\left( {x}^{2}\right) }^{2} \cdot  {{C}}_{4}^{1}\left( {-x}\right)  \cdot  {{C}}_{3}^{3} \times  {2}^{3} =  - {480}{x}^{5}$
                \item 可以取0个 ${x}^{2}$ ，5个 $- x$ ，1个2，则为 ${{C}}_{6}^{0}{\left( {x}^{2}\right) }^{0} \cdot  {{C}}_{6}^{5}{\left( -x\right) }^{5} \cdot  {{C}}_{1}^{1} \times  {2}^{1} =  - {12}{x}^{5}$
            \end{dingautolist}
            综上可得 ${x}^{5}$ 的系数为 $- {240} + \left( {-{480}}\right)  + \left( {-{12}}\right)  =  - {732}$ .
        \end{jieda}
        \item ${\left( 1 - 2x + 3{x}^{2}\right) }^{8}$展开式中$x^2$的系数为\blkda{$136$}
        \begin{jieda}
            ${\left( 1 - 2x + 3{x}^{2}\right) }^{8}$ 的展开式的通项为 ${T}_{r + 1} = {{C}}_{8}^{r}{\left( 1 - 2x\right) }^{8 - r}{3}^{r}{x}^{2r}$ ,
            
            要得到 ${x}^{2}$ 项，必有 ${2r} \leq  2$ ，所以 $r \leq  1$ ，所以 $r = 0$ ，或 $r = 1$ .
            
            当 $r = 0$ 时， ${T}_{1} = {{C}}_{8}^{0}{\left( 1 - 2x\right) }^{8}$ ，而 ${\left( 1 - 2x\right) }^{8}$ 展开式中的 ${x}^{2}$ 项为 ${{C}}_{8}^{2}{\left( -2x\right) }^{2} = {112}{x}^{2}$ ，
            
            故 ${T}_{1}$ 中 ${x}^{2}$ 项的系数为 ${112}{{C}}_{8}^{0} = {112}$ ;
            
            当 $r = 1$ 时， ${T}_{2} = {{C}}_{8}^{1}{\left( 1 - 2x\right) }^{7}{3}^{1}{x}^{2}$ ，而 ${\left( 1 - 2x\right) }^{7}$ 中的常数项为 1，
            
            故 ${T}_{2}$ 中 ${x}^{2}$ 项的系数为 ${{C}}_{8}^{1} \times  1 \times  3 = {24}$ ，所以所求 ${x}^{2}$ 项的系数为 ${112} + {24} = {136}$ .
        \end{jieda}
        \item 多项式 ${\left( {x}^{2} + x + y\right) }^{5}$的展开式中所有项的系数之和为\blkda{$243$}
        \begin{jieda}
            由多项式 ${\left( {x}^{2} + x + y\right) }^{5}$ ,令 $x = y = 1$ ,可得所有项的系数之和为 ${3}^{5} = {243}$ . 
        \end{jieda}
        \item ${\left( x - 2\sqrt{x} + 7\right) }^{3}$ 的展开式中， ${x}^{2}$ 的系数为\blkda{$15$}
        \begin{jieda}
            因为 ${\left( x - 2\sqrt{x} + 1\right) }^{3} = {\left\lbrack  {\left( \sqrt{x}\right) }^{2} - 2\sqrt{x} + 1\right\rbrack  }^{3} = {\left( \sqrt{x} - 1\right) }^{6}$ ,
            
            其中 ${\left( \sqrt{x} - 1\right) }^{6}$ 展开式的通项为 ${T}_{r + 1} = {\left( -1\right) }^{r} \cdot  {{C}}_{6}^{r}{\left( \sqrt{x}\right) }^{6 - r} = {\left( -1\right) }^{r} \cdot  {{C}}_{6}^{r}{x}^{\dfrac{6 - r}{2}}\left( {0 \leq  r \leq  6\text{ 且 }r \in  {N}}\right)$ ,
            
            令 $\dfrac{6 - r}{2} = 2$ ,解得 $r = 2$ ,
            
            所以 ${T}_{3} = {\left( -1\right) }^{2} \cdot  {{C}}_{6}^{2}{x}^{2} = {15}{x}^{2}$ ,
            
            即 ${\left( x - 2\sqrt{x} + 1\right) }^{3}$ 的展开式中 ${x}^{2}$ 的系数为 15.
        \end{jieda}
    \end{enumerate}
    
    \item 证明组合恒等式
    \begin{enumerate}
        \item 证明: ${C}_{n}^{0} + {C}_{n}^{2} + {C}_{n}^{4} + \ldots  + {C}_{n}^{n} = {2}^{n - 1}$ ($n$ 是偶数).
        \begin{jieda}
            ${\left( a + b\right) }^{n} = {C}_{n}^{0}{a}^{n} + {C}_{n}^{1}{a}^{n - 1}b + {C}_{n}^{2}{a}^{n - 2}{b}^{2} + \cdots  + {C}_{n}^{n}{b}^{n}$ ,
            
            令 $a = 1,b = 1$ ,得 ${C}_{n}^{0} + {C}_{n}^{1} + {C}_{n}^{2} + \cdots  + {C}_{n}^{n} = {2}^{n}$ ,
            
            令 $a = 1,b =  - 1$ ,得 ${C}_{n}^{0} - {C}_{n}^{1} + {C}_{n}^{2} - \cdots  + {C}_{n}^{n} = 0$ ,
            
            两式相加得 $2\left( {{C}_{n}^{0} + {C}_{n}^{2} + {C}_{n}^{4} + \cdots  + {C}_{n}^{n}}\right)  = {2}^{n}$ ,
            
            $\therefore {C}_{n}^{0} + {C}_{n}^{2} + {C}_{n}^{4} + \cdots  + {C}_{n}^{n} = {2}^{n - 1}$ .
        \end{jieda}
        \item 求证: ${2}^{n} - {C}_{n}^{1} \times  {2}^{n - 1} + {C}_{n}^{2} \times  {2}^{n - 2} + \ldots  + {\left( -1\right) }^{n - 1}{C}_{n}^{n - 1} \times  2 + (-1)^{n} = 1$ .
        \begin{jieda}
            左边 $= {2}^{n} - {C}_{n}^{1} \times  {2}^{n - 1} + {C}_{n}^{2} \times  {2}^{n - 2} + \ldots  + {\left( -1\right) }^{n - 1}{C}_{n}^{n - 1} \times  2 + {\left( -1\right) }^{n} = {C}_{n}^{0} \times  {2}^{n} \times  {\left( -1\right) }^{0} + {C}_{n}^{1} \times  {2}^{n - 1} \times  {\left( -1\right) }^{1} + {C}_{n}^{2} \times  {2}^{n - 2}{\left( -1\right) }^{2} + \ldots  + {C}_{n}^{n - 1} \times  {2}^{1} \times  {\left( -1\right) }^{n - 1} + {C}_{n}^{n} \times  {2}^{0} \times  (-{1})^{n} = {\left( 2 - 1\right) }^{n} = 1 =$ 右边.
        \end{jieda}
        
        \item 求证: $1 + 4{{C}}_{n}^{1} + 7{{C}}_{n}^{2} + \cdots  + \left( {{3n} + 1}\right) {{C}}_{n}^{n} = \left( {{3n} + 2}\right)  \cdot  {2}^{n - 1}$
        \begin{jieda}
            令 $S = 1 + 4{{C}}_{n}^{1} + 7{{C}}_{n}^{2} + \cdots  + \left( {{3n} + 1}\right) {{C}}_{n}^{n}$ ，则 $S = \left( {{3n} + 1}\right) {{C}}_{n}^{n} + \cdots  + 7{{C}}_{n}^{2} + 4{{C}}_{n}^{1} + 1$ ；
            
            两式相加可得 ${2S} = \left( {{3n} + 2}\right) \left( {1 + {{C}}_{n}^{1} + {{C}}_{n}^{2} + \cdots  + {{C}}_{n}^{n}}\right)  = \left( {{3n} + 2}\right)  \cdot  {2}^{n}$ ,
            
            所以 $S = \dfrac{\left( {{3n} + 2}\right) \left( {1 + {{C}}_{n}^{1} + {{C}}_{n}^{2} + \cdots  + {{C}}_{n}^{n}}\right) }{2} = \dfrac{\left( {{3n} + 2}\right)  \cdot  {2}^{n}}{2} = \left( {{3n} + 2}\right)  \cdot  {2}^{n - 1}$.
        \end{jieda}
    \end{enumerate}
    \item 利用二项式定理证明整除问题或求余数
    \begin{enumerate}
        \item $3^{86}$ 被 8 除所得的余数为\blkda{1}
        \begin{jieda}
            ${3}^{86} = {9}^{43} = {\left( 8 + 1\right) }^{43} = {{C}}_{43}^{0} \times  {8}^{43} \times  {1}^{0} + {{C}}_{43}^{1} \times  {8}^{42} \times  {1}^{1} + \cdots  + {{C}}_{43}^{42} \times  {8}^{1} \times  {1}^{42} + {{C}}_{43}^{43} \times  {8}^{0} \times  {1}^{43}  = 8 \times  \left( {{{C}}_{43}^{0} \times  {8}^{42} + {{C}}_{43}^{1} \times  {8}^{41} + \cdots  + {{C}}_{43}^{42}}\right)  + 1,$
            
            因为 $8 \times  \left( {{{C}}_{43}^{0} \times  {8}^{42} + {{C}}_{43}^{1} \times  {8}^{41} + \cdots  + {{C}}_{43}^{42}}\right)$ 能被 8 整除,
            
            所以 ${3}^{86}$ 被 8 除所得的余数为 1 .
        \end{jieda}
        
        \item ${{C}}_{2024}^{1} + {{C}}_{2024}^{3} + {{C}}_{2024}^{5} + \cdots  + {{C}}_{2024}^{2023}$ 被 3 除的余数为\blkda{2}
        \begin{jieda}
            由二项式定理得 ${\left( 1 + x\right) }^{2024} = {{C}}_{2024}^{0} + {{C}}_{2024}^{1}x + {{C}}_{2024}^{2}{x}^{2} + \cdots  + {{C}}_{2024}^{2023}{x}^{2023} + {{C}}_{2024}^{2024}{x}^{2024}$
            
            令 $x = 1$ 得, ${2}^{2024} = {{C}}_{2024}^{0} + {{C}}_{2024}^{1} + {{C}}_{2024}^{2} + {{C}}_{2024}^{3}\cdots  + {{C}}_{2024}^{2023} + {{C}}_{2024}^{2024}$
            
            令 $x =  - 1$ 得, $0 = {{C}}_{2024}^{0} - {{C}}_{2024}^{1} + {{C}}_{2024}^{2} - {{C}}_{2024}^{3}\cdots  - {{C}}_{2024}^{2023} + {{C}}_{2024}^{2024}$
            
            两式作差可得${{C}}_{2024}^{1} + {{C}}_{2024}^{3} + {{C}}_{2024}^{5} + \cdots  + {{C}}_{2024}^{2023} = {2}^{2023}$
            
            由 ${2}^{2023} = {\left( 3 - 1\right) }^{2023} = {{C}}_{2023}^{0} \times  {3}^{2023} \cdot  {\left( -1\right) }^{0} + {{C}}_{2023}^{1} \times  {3}^{2022} \cdot  {\left( -1\right) }^{1} + {{C}}_{2023}^{2} \times  {3}^{2021} \cdot  {\left( -1\right) }^{2} + \cdots + {{C}}_{2023}^{2022} \times  3 \cdot  {\left( -1\right) }^{2022} + {{C}}_{2023}^{2023}{\left( -1\right) }^{2023} = 3\left\lbrack  {{{C}}_{2023}^{0} \times  {3}^{2022} \cdot  {\left( -1\right) }^{0} + {{C}}_{2023}^{1} \times  {3}^{2021} \cdot  {\left( -1\right) }^{1} + {{C}}_{2023}^{2} \times  {3}^{2020} \cdot  {\left( -1\right) }^{2} + \cdots  + {{C}}_{2023}^{2022} \cdot  {\left( -1\right) }^{2022}}\right\rbrack   - 1$
            
            故 ${{C}}_{2024}^{1} + {{C}}_{2024}^{3} + {{C}}_{2024}^{5} + \cdots  + {{C}}_{2024}^{2023}$ 被 3 除的余数为 2 .
        \end{jieda}
        \item 若 ${2}^{27} + a$ 既能被 9 整除又能被 7 整除，则正整数 $a$ 的最小值为\dotfill{(\kh{C})}
        \xx{6}{10}{55}{63}
        \begin{jieda}
            因为 ${2}^{27} = {\left( {2}^{3}\right) }^{9} = {8}^{9} = {\left( 1 + 7\right) }^{9}$ ,
            
            所以 ${2}^{27} + a = {{C}}_{9}^{0} + {{C}}_{9}^{1}{7}^{1} + {{C}}_{9}^{2}{7}^{2} + {{C}}_{9}^{3}{7}^{3} + {{C}}_{9}^{4}{7}^{4} + {{C}}_{9}^{5}{7}^{5} + {{C}}_{9}^{6}{7}^{6} + {{C}}_{9}^{7}{7}^{7} + {{C}}_{9}^{8}{7}^{8} + {{C}}_{9}^{9}{7}^{9} + a$
            
            $= {{C}}_{9}^{1}{7}^{1} + {{C}}_{9}^{2}{7}^{2} + {{C}}_{9}^{3}{7}^{3} + {{C}}_{9}^{4}{7}^{4} + {{C}}_{9}^{5}{7}^{5} + {{C}}_{9}^{6}{7}^{6} + {{C}}_{9}^{7}{7}^{7} + {{C}}_{9}^{8}{7}^{8} + {{C}}_{9}^{9}{7}^{9} + 1 + a$ ,
            
            所以若 ${2}^{27} + a$ 既能被 7 整除,则 $1 + a = {7m}\left( {m \in  \mathbf{N}}\right)$ ,故 $a = {7m} - 1\left( {m \in  \mathbf{N}}\right)$
            
            又 ${2}^{27} = {\left( {2}^{3}\right) }^{9} = {8}^{9} = {\left( 9 - 1\right) }^{9}$ ,
            
            所以 ${2}^{27} + a = {{C}}_{9}^{0}{9}^{9} - {{C}}_{9}^{1}{9}^{8} + {{C}}_{9}^{2}{9}^{7} - {{C}}_{9}^{3}{9}^{6} + {{C}}_{9}^{4}{9}^{5} - {{C}}_{9}^{5}{9}^{4} + {{C}}_{9}^{6}{9}^{3} - {{C}}_{9}^{7}{9}^{2} + {{C}}_{9}^{8}{9}^{1} - {{C}}_{9}^{9} + a$
            
            $= {{C}}_{9}^{0}{9}^{9} - {{C}}_{9}^{1}{9}^{8} + {{C}}_{9}^{2}{9}^{7} - {{C}}_{9}^{3}{9}^{6} + {{C}}_{9}^{4}{9}^{5} - {{C}}_{9}^{5}{9}^{4} + {{C}}_{9}^{6}{9}^{3} - {{C}}_{9}^{7}{9}^{2} + {{C}}_{9}^{8}{9}^{1} - 1 + a$ ,
            
            所以若 ${2}^{27} + a$ 既能被 9 整除,则 $- 1 + a = {9n}\left( {n \in  {N}}\right)$ ,故 $a = {9n} + 1\left( {n \in  \mathbf{N}}\right)$
        \end{jieda}
        \item 中国南北朝时期的著作《孙子算经》中，对同余除法有较深的研究,设 $a,b,m\left( {m > 0}\right)$ 均为整数,若 $a$ 和 $b$ 被 $m$ 除得的余数相同,则称 $a$ 和 $b$ 对模 $m$ 同余,记为 $a \equiv  b\left( {\operatorname{mod}m}\right)$ ,如 9 和 21 被 6 除得的余数都是 3,则记 $g \equiv  {21}\left( {\operatorname{mod}6}\right)$ . 若 $a \equiv  b\left( {\operatorname{mod}{10}}\right)$ ,且 $a = {{C}}_{20}^{0} +  {{C}}_{20}^{1} \cdot  2 + {{C}}_{20}^{2} \cdot  {2}^{2} + \cdots  + {{C}}_{20}^{20} \cdot  {2}^{20}$ ,则 $b$ 的值可以是\dotfill{(\kh{B})}
        \xx{2010}{2021}{2019}{1997}
        \begin{jieda}
            因为 $a = {{C}}_{20}^{0} + {{C}}_{20}^{1} \cdot  2 + {{C}}_{20}^{2} \cdot  {2}^{2} + \cdots  + {{C}}_{20}^{20} \cdot  {2}^{20} = {\left( 1 + 2\right) }^{20} = {3}^{20} = {9}^{10} = {\left( {10} - 1\right) }^{10}$
            
            又 ${\left( {10} - 1\right) }^{10} = {{C}}_{10}^{0} \cdot  {10}^{10} - {{C}}_{10}^{1} \cdot  {10}^{9} + \cdots  - {{C}}_{10}^{9} \cdot  {10} + {{C}}_{10}^{10}$ ,故 $a \equiv  1\left( {\operatorname{mod}{10}}\right)$ ,
            
            又 ${2010} \div  {10} = {201},{2021} \div  {10} = {202}\cdots 1,{2019} \div  {10} = {201}\cdots 9$ ,
            
            ${1997} \div  {10} = {199}\cdots 7$ ,结合选项可知只有 ${B}$ 符合题意.
        \end{jieda}
    \end{enumerate}
    
    \item 杨辉三角问题
    \begin{enumerate}
        \item 杨辉三角(如下图所示)是数学史上的一个伟大成就，杨辉三角中从第 2 行到第 2024 行，每行的第 3 个数字之和为\blkda{${{C}}_{2025}^{3}$}
        
        \begin{center}
        \includegraphics[width=0.3\textwidth]{bodyxb/src/bo_d67fn677aajc739ac5f0_14_502_986_363_279_0.jpg}
        \end{center}
        \begin{jieda}
            由题意可知,从第 2 行开始,第 $n$ 行的第 3 个数字为 ${{C}}_{n}^{2}$ ,
            
            故从第 2 行到第 2024 行,每行的第 3 个数字之和为
            
            ${{C}}_{2}^{2} + {{C}}_{3}^{2} + {{C}}_{4}^{2} + \cdots  + {{C}}_{2024}^{2} = {{C}}_{3}^{3} + {{C}}_{3}^{2} + {{C}}_{4}^{2} + \cdots  + {{C}}_{2024}^{2} = {{C}}_{4}^{3} + {{C}}_{4}^{2} + \cdots  + {{C}}_{2024}^{2} = \cdots  = {{C}}_{2024}^{3} + {{C}}_{2024}^{2} = {{C}}_{2025}^{3}.$
        \end{jieda}
        
        \item “杨辉三角”是中国古代数学文化的现宝之一，它揭示了二项式展开式中的组合数在三角形数表中的一种几何排列规律，如图所示，则下列关于“杨辉三角”的结论正确的是\dotfill{(\kh{D})}
        \begin{center}
        \includegraphics[width=0.6\textwidth]{bodyxb/src/bo_d67fn677aajc739ac5f0_15_137_233_710_407_0.jpg}
        \end{center}
        \xx{在第 10 行中第 5 个数最大}{第 2023 行中第 1011 个数和第 1012 个数相等}{${{C}}_{3}^{2} + {{C}}_{4}^{2} + {{C}}_{5}^{2} + \cdots  + {{C}}_{20}^{2} = {120}$}{第 6 行的第 7 个数、第 7 行的第 7 个数及第 8 行的第 7 个数之和等于 9 行的第 8 个数}
        \begin{jieda}
            对于 ${A}$ ,因“杨辉三角”的第 10 行中第 5 个数是 ${{C}}_{10}^{4}$ ,又 ${{C}}_{10}^{4} < {{C}}_{10}^{5}$ ,故 ${A}$ 错误;
            
            对于 ${B}$ ,因“杨辉三角”的第 2023 行中第 1011 个数和第 1012 个数分别为 ${{C}}_{2023}^{1010}$ 和 ${{C}}_{2023}^{1011}$ ,
            
            因 1010+1011 $\neq$ 2023,故 ${{C}}_{2023}^{1010} \neq  {{C}}_{2023}^{1011}$ ,故 ${B}$ 错误;
            
            对于 ${C},$ 因 ${{C}}_{3}^{3} + {{C}}_{3}^{2} + {{C}}_{4}^{2} + {{C}}_{5}^{2} + \cdots  + {{C}}_{20}^{2} = {{C}}_{4}^{3} + {{C}}_{4}^{2} + {{C}}_{5}^{2} + \cdots  + {{C}}_{20}^{2}$
            
            $= {{C}}_{5}^{3} + {{C}}_{5}^{2} + \cdots  + {{C}}_{20}^{2} = {{C}}_{20}^{3} + {{C}}_{20}^{2} = {{C}}_{21}^{3} = \dfrac{{21} \times  {20} \times  {19}}{6} = {1330}$ ,
            
            则 ${{C}}_{3}^{2} + {{C}}_{4}^{2} + {{C}}_{5}^{2} + \cdots  + {{C}}_{20}^{2} = {1330} - 1 = {1329}$ ,故 ${C}$ 错误;
            
            对于 ${D}$ ,因 ${{C}}_{6}^{6} + {{C}}_{7}^{6} + {{C}}_{8}^{6} = 1 + 7 + {28} = {36}$ ,而 ${{C}}_{9}^{7} = {{C}}_{9}^{2} = {36}$ ,故 ${D}$ 正确.
        \end{jieda}
        
        \item 我国古代数学史上一位著述丰富的数学家，他提出的杨辉三角是我国古代数学重大成就之一. 图为杨辉三角的部分内容. 设杨辉三角中第 $n$ 行的第 $r$ 个数为 ${{C}}_{n}^{r - 1}$ ,观察题图可知, 相邻两行中三角形的两个腰都是由数字 1 组成的, 其余的数都等于它肩上的两个数相加.
        \begin{center}
        \includegraphics[width=0.3\textwidth]{bodyxb/src/bo_d67fn677aajc739ac5f0_15_238_1750_405_250_0.jpg}
        \end{center}
        \begin{enumerate}
            \item 用公式表示出题目中叙述的规律，并加以证明.
            \item 在杨辉三角中是否存在某一行，使该行中三个相邻的数之比为 $3 : 8 : {14}$ ? 若存在，试求出这三个数；若不存在, 请说明理由.
        \end{enumerate}
        \begin{jieda}
            \begin{enumerate}
                \item 观察得到 ${{C}}_{n}^{r} = {{C}}_{n - 1}^{r - 1} + {{C}}_{n - 1}^{r}$ .
                
                利用组合相关公式证明如下: ${{C}}_{n - 1}^{r - 1} + {{C}}_{n - 1}^{r} = \dfrac{\left( {n - 1}\right) !}{\left( {r - 1}\right) !\left( {n - r}\right) !} + \dfrac{\left( {n - 1}\right) !}{r!\left( {n - 1 - r}\right) !} = \dfrac{\left( {n - 1}\right) !}{r!\left( {n - r}\right) !}\left\lbrack  {r + \left( {n - r}\right) }\right\rbrack   = \dfrac{n!}{r!\left( {n - r}\right) !} = {{C}}_{n}^{r}$ , 故原式得证.

                \item 存在,理由如下:
                
                设在第 $n$ 行存在连续三项 ${{C}}_{n}^{k - 1}$ ， ${{C}}_{n}^{k}$ ， ${{C}}_{n}^{k + 1}$ ，其中 $n \in  {\mathbf{N}}^{ * }$ 且 $n \geq  2$ ， $k \in  {\mathbf{N}}^{ * }$ 且 $k \geq  2$
                
                有 $\dfrac{{{C}}_{n}^{k - 1}}{{{C}}_{n}^{k}} = \dfrac{3}{8}$ 且 $\dfrac{{{C}}_{n}^{k}}{{{C}}_{n}^{k + 1}} = \dfrac{8}{14}$ ,化简得 $\dfrac{k}{n - k + 1} = \dfrac{3}{8}$ 且 $\dfrac{k + 1}{n - k} = \dfrac{8}{14}$ ,
                
                即 $\left\{  \begin{matrix} {3n} + 3 = {11k} \\  {22k} - {8n} + {14} = 0 \end{matrix}\right.$ ，解得 $k = 3$ ， $n = {10}$ ，
                
                故三个数依次是 45, 120, 210.
            \end{enumerate}
        \end{jieda}
        \item 杨辉是中国南宋末年的一位杰出的数学家、教育家，杨辉三角是杨辉的一项重要研究成果.杨辉三角中蕴藏了许多优美的规律，它的许多性质与组合数的性质有关, 图 1 为杨辉三角的部分内容, 图 2 为杨辉三角的改写形式
        
        \begin{center}
        \includegraphics[width=1.0\textwidth]{bodyxb/src/bo_d67fn677aajc739ac5f0_16_140_1324_1193_503_0.jpg}
        \end{center}
        \begin{enumerate}
            \item 求图 2 中第 11 行的各数之和；
            
            \item 从图 2 第 2 行开始，取每一行的第 3 个数一直取到第 100 行的第 3 个数，求取出的所有数之和；
            
            \item 在杨辉三角中是否存在某一行，使该行中三个相邻的数之比为 3:8:14？若存在，试求出这三个数；若不存在，请说明理由.
        \end{enumerate}
        \begin{jieda}
            \begin{enumerate}
                \item 第 11 行的各数之和为 ${{C}}_{11}^{0} + {{C}}_{11}^{1} + {{C}}_{11}^{2} + \cdots  + {{C}}_{11}^{11} = {2}^{11} = {2048}$ 
                
                \item 杨辉三角中第 2 行到第 100 行，各行第 3 个数之和为
                
                ${{C}}_{2}^{2} + {{C}}_{3}^{2} + {{C}}_{4}^{2} + \cdots  + {{C}}_{100}^{2} = {{C}}_{3}^{3} + {{C}}_{3}^{2} + {{C}}_{4}^{2} + \cdots  + {{C}}_{100}^{2} = {{C}}_{4}^{3} + {{C}}_{4}^{2} + \cdots  + {{C}}_{100}^{2} = {{C}}_{101}^{3}  = {{C}}_{101}^{3} = \dfrac{{101} \times  {100} \times  {99}}{3 \times  2 \times  1} = {166650}$ ;
                
                \item 存在,理由如下:
                
                设在第 $n$ 行存在三个相邻的数 ${{C}}_{n}^{k - 1},{{C}}_{n}^{k},{{C}}_{n}^{k + 1}$ ，其中 $k,n \in  {{N}}^{ * }$ ，且 $k \neq  1 \leq  n,n \geq  2$ ， ${{C}}_{n}^{k - 1},{{C}}_{n}^{k},{{C}}_{n}^{k + 1}$ 之比为 $3 : 8 : {14}$ ,
                
                故 $\dfrac{{{C}}_{n}^{k - 1}}{{{C}}_{n}^{k}} = \dfrac{3}{8},\dfrac{{{C}}_{n}^{k}}{{{C}}_{n}^{k + 1}} = \dfrac{8}{14}$ ,化简得 $\dfrac{k}{n - k + 1} = \dfrac{3}{8},\dfrac{k + 1}{n - k} = \dfrac{8}{14}$ ,
                
                即 $\left\{  \begin{matrix} {8k} = {3n} - {3k} + 3 \\  {14k} + {14} = {8n} - {8k} \end{matrix}\right.$ ,解得 $\left\{  \begin{array}{l} k = 3 \\  n = {10} \end{array}\right.$ ,
                
                所以这三个数为 ${{C}}_{10}^{2} = {45},{{C}}_{10}^{3} = {120},{{C}}_{10}^{4} = {210}$ .
            \end{enumerate}
        \end{jieda}
    \end{enumerate}
\end{enumerate}